\documentclass[12pt]{report}
\usepackage[utf8x]{inputenc}

%Warning: tipa declares many non-standard macros used by utf8x to
%interpret utf8 characters but extra packages might have to be added
%such as "textgreek" for Greek letters not already in tipa
%or "stmaryrd" for mathematical symbols.
%Utf8 codes missing a LaTeX interpretation can be defined by using
%\DeclareUnicodeCharacter{code}{interpretation}.
%Use coqdoc's option -p to add new packages or declarations.
\usepackage{tipa}

\usepackage[T1]{fontenc}
\usepackage{fullpage}
\usepackage{coqdoc}
\usepackage{amsmath,amssymb}
\usepackage{url}
\begin{document}
%%%%%%%%%%%%%%%%%%%%%%%%%%%%%%%%%%%%%%%%%%%%%%%%%%%%%%%%%%%%%%%%%
%% This file has been automatically generated with the command
%% /nix/store/5nhkcrnn3hzlx1rsgy6zxlg2lrla4f5y-coq-with-packages-9.1.1/bin-orig/coqdoc -latex -toc -interpolate -utf8 -Q . SECF Maps.v Imp.v Equiv.v Hoare.v Hoare2.v Preface.v Noninterference.v StaticIFC.v SpecCT.v Postscript.v Bib.v MapsTest.v ImpTest.v EquivTest.v HoareTest.v Hoare2Test.v PrefaceTest.v NoninterferenceTest.v StaticIFCTest.v SpecCTTest.v PostscriptTest.v BibTest.v 
%%%%%%%%%%%%%%%%%%%%%%%%%%%%%%%%%%%%%%%%%%%%%%%%%%%%%%%%%%%%%%%%%
\coqlibrary{SECF.Equiv}{Library }{SECF.Equiv}

\begin{coqdoccode}
\end{coqdoccode}
\section{Equiv: Program Equivalence}


\begin{coqdoccode}
\coqdocemptyline
\coqdocnoindent
\coqdockw{Set} \coqdocvar{Warnings} "-notation-overridden".\coqdoceol
\coqdocnoindent
\coqdockw{From} \coqdocvar{SECF} \coqdockw{Require} \coqdockw{Import} \coqref{SECF.Maps}{\coqdoclibrary{Maps}}.\coqdoceol
\coqdocnoindent
\coqdockw{From} \coqdocvar{Stdlib} \coqdockw{Require} \coqdockw{Import} \coqdoclibrary{Bool}.\coqdoceol
\coqdocnoindent
\coqdockw{From} \coqdocvar{Stdlib} \coqdockw{Require} \coqdockw{Import} \coqdoclibrary{Arith}.\coqdoceol
\coqdocnoindent
\coqdockw{From} \coqdocvar{Stdlib} \coqdockw{Require} \coqdockw{Import} \coqdoclibrary{Init.Nat}.\coqdoceol
\coqdocnoindent
\coqdockw{From} \coqdocvar{Stdlib} \coqdockw{Require} \coqdockw{Import} \coqdoclibrary{PeanoNat}. \coqdockw{Import} \coqdocmodule{Nat}.\coqdoceol
\coqdocnoindent
\coqdockw{From} \coqdocvar{Stdlib} \coqdockw{Require} \coqdockw{Import} \coqdoclibrary{EqNat}.\coqdoceol
\coqdocnoindent
\coqdockw{From} \coqdocvar{Stdlib} \coqdockw{Require} \coqdockw{Import} \coqdoclibrary{Lia}.\coqdoceol
\coqdocnoindent
\coqdockw{From} \coqdocvar{Stdlib} \coqdockw{Require} \coqdockw{Import} \coqdoclibrary{List}. \coqdockw{Import} \coqdocmodule{ListNotations}.\coqdoceol
\coqdocnoindent
\coqdockw{From} \coqdocvar{Stdlib} \coqdockw{Require} \coqdockw{Import} \coqdoclibrary{FunctionalExtensionality}.\coqdoceol
\coqdocnoindent
\coqdockw{From} \coqdocvar{SECF} \coqdockw{Require} \coqdockw{Export} \coqref{SECF.Imp}{\coqdoclibrary{Imp}}.\coqdoceol
\coqdocemptyline
\end{coqdoccode}
\subsubsection{Before You Get Started:}





\begin{itemize}
\item  Create a fresh directory for this volume. Do not try to mix the
      files from this volume with files from \textit{Logical Foundations} in
      the same directory: the result will not make you happy.


      You can either start with an empty directory and populate it
      with the files listed below (and others as you go along), or
      else download the whole PLF zip file and unzip it.



\item  The new directory should contain at least the following files:

\begin{itemize}
\item  \coqdocvar{Imp.v} (make sure it is the one from the PLF distribution,
           not the one from LF: they are slightly different);

\item  \coqdocvar{Maps.v} (ditto)

\item  \coqdocvar{Equiv.v} (this file)

\item  \coqdocvar{\_CoqProject}, containing the following line:



\begin{itemize}
\item Q . PLF



\end{itemize}

\end{itemize}

\item  If you see errors like this...


           Compiled library PLF.Maps (in file .../plf/Maps.vo)
           makes inconsistent assumptions over library Rocq.Init.Logic


      ... it may mean something went wrong with the above steps.
      Doing ``\coqdocvar{make} \coqdocvar{clean}'' (or manually removing everything except
      .\coqdocvar{v} and \coqdocvar{\_CoqProject} files) may help.



\item  If you are using VSCode with the VSCoq extension, you'll then
      want to open a new window in VSCode, click \coqdockw{Open} \coqdocvar{Folder} > \coqdocvar{plf},
      and run \coqdocvar{make}.  If you get an error like ``Cannot find a
      physical path...'' error, it may be because you didn't open plf
      directly (you instead opened a folder containing both lf and
      plf, for example). 
\end{itemize}


\subsubsection{Advice for Working on Exercises:}





\begin{itemize}
\item  Most of the Rocq proofs we ask you to do in this chapter are
      similar to proofs that we've provided.  Before starting to work
      on exercises, take the time to work through our proofs (both
      informally and in Rocq) and make sure you understand them in
      detail.  This will save you a lot of effort.



\item  The Rocq proofs we're doing now are sufficiently complicated that
      it is more or less impossible to complete them by random
      experimentation or following your nose.  You need to start with
      an idea about why the property is true and how the proof is
      going to go.  The best way to do this is to write out at least a
      sketch of an informal proof on paper -- one that intuitively
      convinces you of the truth of the theorem -- before starting to
      work on the formal one.  Alternately, grab a friend and try to
      convince them that the theorem is true; then try to formalize
      your explanation.



\item  Use automation to save work!  The proofs in this chapter can get
      pretty long if you try to write out all the cases explicitly. 

\end{itemize}
\begin{coqdoccode}
\coqdocemptyline
\end{coqdoccode}
\section{Behavioral Equivalence}



 In an earlier chapter, we investigated the correctness of a very
    simple program transformation: the \coqref{SECF.Imp.AExp.optimize 0plus}{\coqdocdefinition{optimize\_0plus}} function.  The
    programming language we were considering was the first version of
    the language of arithmetic expressions -- with no variables -- so
    in that setting it was very easy to define what it means for a
    program transformation to be correct: it should always yield a
    program that evaluates to the same number as the original.


    To talk about the correctness of program transformations for the
    full Imp language -- in particular, assignment -- we need to
    consider the role of mutable state and develop a more
    sophisticated notion of correctness, which we'll call \textit{behavioral
    equivalence}. 

 For example:

\begin{itemize}
\item  \coqref{SECF.SpecCT.X}{\coqdocdefinition{X}} + 2 is behaviorally equivalent to 1 + \coqref{SECF.SpecCT.X}{\coqdocdefinition{X}} + 1

\item  \coqref{SECF.SpecCT.X}{\coqdocdefinition{X}} - \coqref{SECF.SpecCT.X}{\coqdocdefinition{X}} is behaviorally equivalent to 0

\item  (\coqref{SECF.SpecCT.X}{\coqdocdefinition{X}} - 1) + 1 is \textit{not} behaviorally equivalent to \coqref{SECF.SpecCT.X}{\coqdocdefinition{X}}   

\end{itemize}
\begin{coqdoccode}
\coqdocemptyline
\end{coqdoccode}
\subsection{Definitions}



 For \coqref{SECF.Imp.aexp}{\coqdocinductive{aexp}}s and \coqref{SECF.Imp.bexp}{\coqdocinductive{bexp}}s with variables, the definition we want is
    clear: Two \coqref{SECF.Imp.aexp}{\coqdocinductive{aexp}}s or \coqref{SECF.Imp.bexp}{\coqdocinductive{bexp}}s are ``behaviorally equivalent'' if
    they evaluate to the same result in every state. 
\begin{coqdoccode}
\coqdocemptyline
\coqdocnoindent
\coqdockw{Definition} \coqdef{SECF.Equiv.aequiv}{aequiv}{\coqdocdefinition{aequiv}} (\coqdef{SECF.Equiv.a1:1}{a1}{\coqdocbinder{a1}} \coqdef{SECF.Equiv.a2:2}{a2}{\coqdocbinder{a2}} : \coqref{SECF.Imp.aexp}{\coqdocinductive{aexp}}) : \coqdockw{Prop} :=\coqdoceol
\coqdocindent{1.00em}
\coqdockw{\ensuremath{\forall}} (\coqdef{SECF.Equiv.st:3}{st}{\coqdocbinder{st}} : \coqref{SECF.Imp.state}{\coqdocdefinition{state}}),\coqdoceol
\coqdocindent{2.00em}
\coqref{SECF.Imp.aeval}{\coqdocdefinition{aeval}} \coqref{SECF.Equiv.st:3}{\coqdocvariable{st}} \coqref{SECF.Equiv.a1:1}{\coqdocvariable{a1}} \coqexternalref{::type scope:x '=' x}{http://rocq-prover.org/doc/V9.1.1/stdlib//Corelib.Init.Logic}{\coqdocnotation{=}} \coqref{SECF.Imp.aeval}{\coqdocdefinition{aeval}} \coqref{SECF.Equiv.st:3}{\coqdocvariable{st}} \coqref{SECF.Equiv.a2:2}{\coqdocvariable{a2}}.\coqdoceol
\coqdocemptyline
\coqdocnoindent
\coqdockw{Definition} \coqdef{SECF.Equiv.bequiv}{bequiv}{\coqdocdefinition{bequiv}} (\coqdef{SECF.Equiv.b1:4}{b1}{\coqdocbinder{b1}} \coqdef{SECF.Equiv.b2:5}{b2}{\coqdocbinder{b2}} : \coqref{SECF.Imp.bexp}{\coqdocinductive{bexp}}) : \coqdockw{Prop} :=\coqdoceol
\coqdocindent{1.00em}
\coqdockw{\ensuremath{\forall}} (\coqdef{SECF.Equiv.st:6}{st}{\coqdocbinder{st}} : \coqref{SECF.Imp.state}{\coqdocdefinition{state}}),\coqdoceol
\coqdocindent{2.00em}
\coqref{SECF.Imp.beval}{\coqdocdefinition{beval}} \coqref{SECF.Equiv.st:6}{\coqdocvariable{st}} \coqref{SECF.Equiv.b1:4}{\coqdocvariable{b1}} \coqexternalref{::type scope:x '=' x}{http://rocq-prover.org/doc/V9.1.1/stdlib//Corelib.Init.Logic}{\coqdocnotation{=}} \coqref{SECF.Imp.beval}{\coqdocdefinition{beval}} \coqref{SECF.Equiv.st:6}{\coqdocvariable{st}} \coqref{SECF.Equiv.b2:5}{\coqdocvariable{b2}}.\coqdoceol
\coqdocemptyline
\end{coqdoccode}
Here are some simple examples of equivalences of arithmetic
    and boolean expressions. 
\begin{coqdoccode}
\coqdocemptyline
\coqdocnoindent
\coqdockw{Theorem} \coqdef{SECF.Equiv.aequiv example}{aequiv\_example}{\coqdoclemma{aequiv\_example}}:\coqdoceol
\coqdocindent{1.00em}
\coqref{SECF.Equiv.aequiv}{\coqdocdefinition{aequiv}}\coqdoceol
\coqdocindent{2.00em}
\coqref{SECF.Imp.::com scope:'<x7B' x 'x7D>'}{\coqdocnotation{<\{}} \coqref{SECF.Imp.X}{\coqdocdefinition{X}} \coqref{SECF.Imp.:com::x '-' x}{\coqdocnotation{-}} \coqref{SECF.Imp.X}{\coqdocdefinition{X}} \coqref{SECF.Imp.::com scope:'<x7B' x 'x7D>'}{\coqdocnotation{\}>}}\coqdoceol
\coqdocindent{2.00em}
\coqref{SECF.Imp.::com scope:'<x7B' x 'x7D>'}{\coqdocnotation{<\{}} 0 \coqref{SECF.Imp.::com scope:'<x7B' x 'x7D>'}{\coqdocnotation{\}>}}.\coqdoceol
\coqdocnoindent
\coqdockw{Proof}.\coqdoceol
\coqdocindent{1.00em}
\coqdoctac{intros} \coqdocvar{st}. \coqdoctac{simpl}. \coqdocvar{lia}.\coqdoceol
\coqdocnoindent
\coqdockw{Qed}.\coqdoceol
\coqdocemptyline
\coqdocnoindent
\coqdockw{Theorem} \coqdef{SECF.Equiv.bequiv example}{bequiv\_example}{\coqdoclemma{bequiv\_example}}:\coqdoceol
\coqdocindent{1.00em}
\coqref{SECF.Equiv.bequiv}{\coqdocdefinition{bequiv}}\coqdoceol
\coqdocindent{2.00em}
\coqref{SECF.Imp.::com scope:'<x7B' x 'x7D>'}{\coqdocnotation{<\{}} \coqref{SECF.Imp.X}{\coqdocdefinition{X}} \coqref{SECF.Imp.:com::x '-' x}{\coqdocnotation{-}} \coqref{SECF.Imp.X}{\coqdocdefinition{X}} \coqref{SECF.Imp.:com::x '=' x}{\coqdocnotation{=}} 0 \coqref{SECF.Imp.::com scope:'<x7B' x 'x7D>'}{\coqdocnotation{\}>}}\coqdoceol
\coqdocindent{2.00em}
\coqref{SECF.Imp.::com scope:'<x7B' x 'x7D>'}{\coqdocnotation{<\{}} \coqref{SECF.Imp.:com::'true'}{\coqdocnotation{true}} \coqref{SECF.Imp.::com scope:'<x7B' x 'x7D>'}{\coqdocnotation{\}>}}.\coqdoceol
\coqdocnoindent
\coqdockw{Proof}.\coqdoceol
\coqdocindent{1.00em}
\coqdoctac{intros} \coqdocvar{st}. \coqdoctac{unfold} \coqref{SECF.Imp.beval}{\coqdocdefinition{beval}}.\coqdoceol
\coqdocindent{1.00em}
\coqdoctac{rewrite} \coqref{SECF.Equiv.aequiv example}{\coqdoclemma{aequiv\_example}}. \coqdoctac{reflexivity}.\coqdoceol
\coqdocnoindent
\coqdockw{Qed}.\coqdoceol
\coqdocemptyline
\end{coqdoccode}
For commands, the situation is a little more subtle.  We
    can't simply say ``two commands are behaviorally equivalent if they
    evaluate to the same ending state whenever they are started in the
    same initial state,'' because some commands, when run in some
    starting states, don't terminate in any final state at all!


    What we need instead is this: two commands are behaviorally
    equivalent if, for any given starting state, they either (1) both
    diverge or else (2) both terminate in the same final state.  A
    compact way to express this is ``if the first one terminates in a
    particular state then so does the second, and vice versa.'' 
\begin{coqdoccode}
\coqdocemptyline
\coqdocnoindent
\coqdockw{Definition} \coqdef{SECF.Equiv.cequiv}{cequiv}{\coqdocdefinition{cequiv}} (\coqdef{SECF.Equiv.c1:7}{c1}{\coqdocbinder{c1}} \coqdef{SECF.Equiv.c2:8}{c2}{\coqdocbinder{c2}} : \coqref{SECF.Imp.com}{\coqdocinductive{com}}) : \coqdockw{Prop} :=\coqdoceol
\coqdocindent{1.00em}
\coqdockw{\ensuremath{\forall}} (\coqdef{SECF.Equiv.st:9}{st}{\coqdocbinder{st}} \coqdef{SECF.Equiv.st':10}{st'}{\coqdocbinder{st'}} : \coqref{SECF.Imp.state}{\coqdocdefinition{state}}),\coqdoceol
\coqdocindent{2.00em}
\coqexternalref{::type scope:x '<->' x}{http://rocq-prover.org/doc/V9.1.1/stdlib//Corelib.Init.Logic}{\coqdocnotation{(}}\coqref{SECF.Equiv.st:9}{\coqdocvariable{st}} \coqref{SECF.Imp.:::x '=[' x ']=>' x}{\coqdocnotation{=[}} \coqref{SECF.Equiv.c1:7}{\coqdocvariable{c1}} \coqref{SECF.Imp.:::x '=[' x ']=>' x}{\coqdocnotation{]=>}} \coqref{SECF.Equiv.st':10}{\coqdocvariable{st'}}\coqexternalref{::type scope:x '<->' x}{http://rocq-prover.org/doc/V9.1.1/stdlib//Corelib.Init.Logic}{\coqdocnotation{)}} \coqexternalref{::type scope:x '<->' x}{http://rocq-prover.org/doc/V9.1.1/stdlib//Corelib.Init.Logic}{\coqdocnotation{\ensuremath{\leftrightarrow}}} \coqexternalref{::type scope:x '<->' x}{http://rocq-prover.org/doc/V9.1.1/stdlib//Corelib.Init.Logic}{\coqdocnotation{(}}\coqref{SECF.Equiv.st:9}{\coqdocvariable{st}} \coqref{SECF.Imp.:::x '=[' x ']=>' x}{\coqdocnotation{=[}} \coqref{SECF.Equiv.c2:8}{\coqdocvariable{c2}} \coqref{SECF.Imp.:::x '=[' x ']=>' x}{\coqdocnotation{]=>}} \coqref{SECF.Equiv.st':10}{\coqdocvariable{st'}}\coqexternalref{::type scope:x '<->' x}{http://rocq-prover.org/doc/V9.1.1/stdlib//Corelib.Init.Logic}{\coqdocnotation{)}}.\coqdoceol
\coqdocemptyline
\end{coqdoccode}
\subsection{Simple Examples}



 For examples of command equivalence, let's start by looking at
    a trivial equivalence involving \coqdocvar{skip}: 
\begin{coqdoccode}
\coqdocemptyline
\coqdocnoindent
\coqdockw{Theorem} \coqdef{SECF.Equiv.skip left}{skip\_left}{\coqdoclemma{skip\_left}} : \coqdockw{\ensuremath{\forall}} \coqdef{SECF.Equiv.c:11}{c}{\coqdocbinder{c}},\coqdoceol
\coqdocindent{1.00em}
\coqref{SECF.Equiv.cequiv}{\coqdocdefinition{cequiv}}\coqdoceol
\coqdocindent{2.00em}
\coqref{SECF.Imp.::com scope:'<x7B' x 'x7D>'}{\coqdocnotation{<\{}} \coqref{SECF.Imp.:com:com scope:'skip'}{\coqdocnotation{skip}}\coqref{SECF.Imp.:com:com scope:x ';' x}{\coqdocnotation{;}} \coqref{SECF.Equiv.c:11}{\coqdocvariable{c}} \coqref{SECF.Imp.::com scope:'<x7B' x 'x7D>'}{\coqdocnotation{\}>}}\coqdoceol
\coqdocindent{2.00em}
\coqref{SECF.Equiv.c:11}{\coqdocvariable{c}}.\coqdoceol
\coqdocnoindent
\coqdockw{Proof}.\coqdoceol
\coqdocindent{1.00em}
\coqdoctac{intros} \coqdocvar{c} \coqdocvar{st} \coqdocvar{st'}.\coqdoceol
\coqdocindent{1.00em}
\coqdoctac{split}; \coqdoctac{intros} \coqdocvar{H}.\coqdoceol
\coqdocindent{1.00em}
- \coqdoceol
\coqdocindent{2.00em}
\coqdoctac{inversion} \coqdocvar{H}. \coqdoctac{subst}.\coqdoceol
\coqdocindent{2.00em}
\coqdoctac{inversion} \coqdocvar{H2}. \coqdoctac{subst}.\coqdoceol
\coqdocindent{2.00em}
\coqdoctac{assumption}.\coqdoceol
\coqdocindent{1.00em}
- \coqdoceol
\coqdocindent{2.00em}
\coqdoctac{apply} \coqref{SECF.Imp.E Seq}{\coqdocconstructor{E\_Seq}} \coqdockw{with} \coqdocvar{st}.\coqdoceol
\coqdocindent{2.00em}
+ \coqdoctac{apply} \coqref{SECF.Imp.E Skip}{\coqdocconstructor{E\_Skip}}.\coqdoceol
\coqdocindent{2.00em}
+ \coqdoctac{assumption}.\coqdoceol
\coqdocnoindent
\coqdockw{Qed}.\coqdoceol
\coqdocemptyline
\end{coqdoccode}
\paragraph{Exercise: 2 stars, standard (skip\_right)}




    Prove that adding a \coqdocvar{skip} \textit{after} a command also results in an
    equivalent program 
\begin{coqdoccode}
\coqdocemptyline
\coqdocnoindent
\coqdockw{Theorem} \coqdef{SECF.Equiv.skip right}{skip\_right}{\coqdoclemma{skip\_right}} : \coqdockw{\ensuremath{\forall}} \coqdef{SECF.Equiv.c:12}{c}{\coqdocbinder{c}},\coqdoceol
\coqdocindent{1.00em}
\coqref{SECF.Equiv.cequiv}{\coqdocdefinition{cequiv}}\coqdoceol
\coqdocindent{2.00em}
\coqref{SECF.Imp.::com scope:'<x7B' x 'x7D>'}{\coqdocnotation{<\{}} \coqref{SECF.Equiv.c:12}{\coqdocvariable{c}}\coqref{SECF.Imp.:com:com scope:x ';' x}{\coqdocnotation{;}} \coqref{SECF.Imp.:com:com scope:'skip'}{\coqdocnotation{skip}} \coqref{SECF.Imp.::com scope:'<x7B' x 'x7D>'}{\coqdocnotation{\}>}}\coqdoceol
\coqdocindent{2.00em}
\coqref{SECF.Equiv.c:12}{\coqdocvariable{c}}.\coqdoceol
\coqdocnoindent
\coqdockw{Proof}.\coqdoceol
 \coqdocvar{Admitted}.\coqdoceol
\end{coqdoccode}
\ensuremath{\Box} 

 Similarly, here is a simple equivalence that optimizes \coqdockw{if}
    commands: 
\begin{coqdoccode}
\coqdocemptyline
\coqdocnoindent
\coqdockw{Theorem} \coqdef{SECF.Equiv.if true simple}{if\_true\_simple}{\coqdoclemma{if\_true\_simple}} : \coqdockw{\ensuremath{\forall}} \coqdef{SECF.Equiv.c1:13}{c1}{\coqdocbinder{c1}} \coqdef{SECF.Equiv.c2:14}{c2}{\coqdocbinder{c2}},\coqdoceol
\coqdocindent{1.00em}
\coqref{SECF.Equiv.cequiv}{\coqdocdefinition{cequiv}}\coqdoceol
\coqdocindent{2.00em}
\coqref{SECF.Imp.::com scope:'<x7B' x 'x7D>'}{\coqdocnotation{<\{}} \coqref{SECF.Imp.:com:com scope:'if' x 'then' x 'else' x 'end'}{\coqdocnotation{if}} \coqref{SECF.Imp.:com::'true'}{\coqdocnotation{true}} \coqref{SECF.Imp.:com:com scope:'if' x 'then' x 'else' x 'end'}{\coqdocnotation{then}} \coqref{SECF.Equiv.c1:13}{\coqdocvariable{c1}} \coqref{SECF.Imp.:com:com scope:'if' x 'then' x 'else' x 'end'}{\coqdocnotation{else}} \coqref{SECF.Equiv.c2:14}{\coqdocvariable{c2}} \coqref{SECF.Imp.:com:com scope:'if' x 'then' x 'else' x 'end'}{\coqdocnotation{end}} \coqref{SECF.Imp.::com scope:'<x7B' x 'x7D>'}{\coqdocnotation{\}>}}\coqdoceol
\coqdocindent{2.00em}
\coqref{SECF.Equiv.c1:13}{\coqdocvariable{c1}}.\coqdoceol
\coqdocnoindent
\coqdockw{Proof}.\coqdoceol
\coqdocindent{1.00em}
\coqdoctac{intros} \coqdocvar{c1} \coqdocvar{c2}.\coqdoceol
\coqdocindent{1.00em}
\coqdoctac{split}; \coqdoctac{intros} \coqdocvar{H}.\coqdoceol
\coqdocindent{1.00em}
- \coqdoceol
\coqdocindent{2.00em}
\coqdoctac{inversion} \coqdocvar{H}; \coqdoctac{subst}.\coqdoceol
\coqdocindent{2.00em}
+ \coqdoctac{assumption}.\coqdoceol
\coqdocindent{2.00em}
+ \coqdoctac{discriminate}.\coqdoceol
\coqdocindent{1.00em}
- \coqdoceol
\coqdocindent{2.00em}
\coqdoctac{apply} \coqref{SECF.Imp.E IfTrue}{\coqdocconstructor{E\_IfTrue}}.\coqdoceol
\coqdocindent{2.00em}
+ \coqdoctac{reflexivity}.\coqdoceol
\coqdocindent{2.00em}
+ \coqdoctac{assumption}. \coqdockw{Qed}.\coqdoceol
\coqdocemptyline
\end{coqdoccode}
Of course, no programmer would write a conditional whose condition
    is literally \coqexternalref{true}{http://rocq-prover.org/doc/V9.1.1/stdlib//Corelib.Init.Datatypes}{\coqdocconstructor{true}}.  (At least, no human programmer -- compilers
    and macro preprocessors do this sort of thing internally all the
    time!) But they might write one whose condition is \textit{equivalent} to
    true: 

 \textit{Theorem}: If \coqdocvar{b} is equivalent to \coqexternalref{true}{http://rocq-prover.org/doc/V9.1.1/stdlib//Corelib.Init.Datatypes}{\coqdocconstructor{true}}, then \coqdockw{if} \coqdocvar{b} \coqdockw{then} \coqdocvar{c1}
    \coqdockw{else} \coqdocvar{c2} \coqdockw{end} is equivalent to \coqdocvar{c1}.
   \textit{Proof}:

\begin{itemize}
\item  (\ensuremath{\rightarrow}) We must show, for all \coqdocvar{st} and \coqdocvar{st'}, that if \coqdocvar{st} =[
       \coqdockw{if} \coqdocvar{b} \coqdockw{then} \coqdocvar{c1} \coqdockw{else} \coqdocvar{c2} \coqdockw{end} ]=> \coqdocvar{st'} then \coqdocvar{st} =[ \coqdocvar{c1} ]=> \coqdocvar{st'}.


       Proceed by cases on the rules that could possibly have been
       used to show \coqdocvar{st} =[ \coqdockw{if} \coqdocvar{b} \coqdockw{then} \coqdocvar{c1} \coqdockw{else} \coqdocvar{c2} \coqdockw{end} ]=> \coqdocvar{st'}, namely
       \coqref{SECF.Hoare.HoareAssertAssume.E IfTrue}{\coqdocconstructor{E\_IfTrue}} and \coqref{SECF.Hoare.HoareAssertAssume.E IfFalse}{\coqdocconstructor{E\_IfFalse}}.



\begin{itemize}
\item  Suppose the final rule in the derivation of \coqdocvar{st} =[ \coqdockw{if} \coqdocvar{b}
         \coqdockw{then} \coqdocvar{c1} \coqdockw{else} \coqdocvar{c2} \coqdockw{end} ]=> \coqdocvar{st'} was \coqref{SECF.Hoare.HoareAssertAssume.E IfTrue}{\coqdocconstructor{E\_IfTrue}}.  We then have,
         by the premises of \coqref{SECF.Hoare.HoareAssertAssume.E IfTrue}{\coqdocconstructor{E\_IfTrue}}, that \coqdocvar{st} =[ \coqdocvar{c1} ]=> \coqdocvar{st'}.
         This is exactly what we set out to prove.



\item  On the other hand, suppose the final rule in the derivation
         of \coqdocvar{st} =[ \coqdockw{if} \coqdocvar{b} \coqdockw{then} \coqdocvar{c1} \coqdockw{else} \coqdocvar{c2} \coqdockw{end} ]=> \coqdocvar{st'} was \coqref{SECF.Hoare.HoareAssertAssume.E IfFalse}{\coqdocconstructor{E\_IfFalse}}.
         We then know that \coqref{SECF.Imp.beval}{\coqdocdefinition{beval}} \coqdocvar{st} \coqdocvar{b} = \coqexternalref{false}{http://rocq-prover.org/doc/V9.1.1/stdlib//Corelib.Init.Datatypes}{\coqdocconstructor{false}} and \coqdocvar{st} =[ \coqdocvar{c2} ]=> \coqdocvar{st'}.


         Recall that \coqdocvar{b} is equivalent to \coqexternalref{true}{http://rocq-prover.org/doc/V9.1.1/stdlib//Corelib.Init.Datatypes}{\coqdocconstructor{true}}, i.e., forall \coqdocvar{st},
         \coqref{SECF.Imp.beval}{\coqdocdefinition{beval}} \coqdocvar{st} \coqdocvar{b} = \coqref{SECF.Imp.beval}{\coqdocdefinition{beval}} \coqdocvar{st} <\{\coqexternalref{true}{http://rocq-prover.org/doc/V9.1.1/stdlib//Corelib.Init.Datatypes}{\coqdocconstructor{true}}\}> .  In particular, this means
         that \coqref{SECF.Imp.beval}{\coqdocdefinition{beval}} \coqdocvar{st} \coqdocvar{b} = \coqexternalref{true}{http://rocq-prover.org/doc/V9.1.1/stdlib//Corelib.Init.Datatypes}{\coqdocconstructor{true}}, since \coqref{SECF.Imp.beval}{\coqdocdefinition{beval}} \coqdocvar{st} <\{\coqexternalref{true}{http://rocq-prover.org/doc/V9.1.1/stdlib//Corelib.Init.Datatypes}{\coqdocconstructor{true}}\}> = \coqexternalref{true}{http://rocq-prover.org/doc/V9.1.1/stdlib//Corelib.Init.Datatypes}{\coqdocconstructor{true}}.  But
         this is a contradiction, since \coqref{SECF.Hoare.HoareAssertAssume.E IfFalse}{\coqdocconstructor{E\_IfFalse}} requires that
         \coqref{SECF.Imp.beval}{\coqdocdefinition{beval}} \coqdocvar{st} \coqdocvar{b} = \coqexternalref{false}{http://rocq-prover.org/doc/V9.1.1/stdlib//Corelib.Init.Datatypes}{\coqdocconstructor{false}}.  Thus, the final rule could not have
         been \coqref{SECF.Hoare.HoareAssertAssume.E IfFalse}{\coqdocconstructor{E\_IfFalse}}.



\end{itemize}

\item  (\ensuremath{\leftarrow}) We must show, for all \coqdocvar{st} and \coqdocvar{st'}, that if
       \coqdocvar{st} =[ \coqdocvar{c1} ]=> \coqdocvar{st'} then
       \coqdocvar{st} =[ \coqdockw{if} \coqdocvar{b} \coqdockw{then} \coqdocvar{c1} \coqdockw{else} \coqdocvar{c2} \coqdockw{end} ]=> \coqdocvar{st'}.


       Since \coqdocvar{b} is equivalent to \coqexternalref{true}{http://rocq-prover.org/doc/V9.1.1/stdlib//Corelib.Init.Datatypes}{\coqdocconstructor{true}}, we know that \coqref{SECF.Imp.beval}{\coqdocdefinition{beval}} \coqdocvar{st} \coqdocvar{b} =
       \coqref{SECF.Imp.beval}{\coqdocdefinition{beval}} \coqdocvar{st} <\{\coqexternalref{true}{http://rocq-prover.org/doc/V9.1.1/stdlib//Corelib.Init.Datatypes}{\coqdocconstructor{true}}\}>  = \coqexternalref{true}{http://rocq-prover.org/doc/V9.1.1/stdlib//Corelib.Init.Datatypes}{\coqdocconstructor{true}}.  Together with the assumption that
       \coqdocvar{st} =[ \coqdocvar{c1} ]=> \coqdocvar{st'}, we can apply \coqref{SECF.Hoare.HoareAssertAssume.E IfTrue}{\coqdocconstructor{E\_IfTrue}} to derive
       \coqdocvar{st} =[ \coqdockw{if} \coqdocvar{b} \coqdockw{then} \coqdocvar{c1} \coqdockw{else} \coqdocvar{c2} \coqdockw{end} ]=> \coqdocvar{st'}.  \ensuremath{\Box}

\end{itemize}


   Here is the formal version of this proof: 
\begin{coqdoccode}
\coqdocemptyline
\coqdocnoindent
\coqdockw{Theorem} \coqdef{SECF.Equiv.if true}{if\_true}{\coqdoclemma{if\_true}}: \coqdockw{\ensuremath{\forall}} \coqdef{SECF.Equiv.b:15}{b}{\coqdocbinder{b}} \coqdef{SECF.Equiv.c1:16}{c1}{\coqdocbinder{c1}} \coqdef{SECF.Equiv.c2:17}{c2}{\coqdocbinder{c2}},\coqdoceol
\coqdocindent{1.00em}
\coqref{SECF.Equiv.bequiv}{\coqdocdefinition{bequiv}} \coqref{SECF.Equiv.b:15}{\coqdocvariable{b}} \coqref{SECF.Imp.::com scope:'<x7B' x 'x7D>'}{\coqdocnotation{<\{}}\coqref{SECF.Imp.:com::'true'}{\coqdocnotation{true}}\coqref{SECF.Imp.::com scope:'<x7B' x 'x7D>'}{\coqdocnotation{\}>}}  \coqexternalref{::type scope:x '->' x}{http://rocq-prover.org/doc/V9.1.1/stdlib//Corelib.Init.Logic}{\coqdocnotation{\ensuremath{\rightarrow}}}\coqdoceol
\coqdocindent{1.00em}
\coqref{SECF.Equiv.cequiv}{\coqdocdefinition{cequiv}}\coqdoceol
\coqdocindent{2.00em}
\coqref{SECF.Imp.::com scope:'<x7B' x 'x7D>'}{\coqdocnotation{<\{}} \coqref{SECF.Imp.:com:com scope:'if' x 'then' x 'else' x 'end'}{\coqdocnotation{if}} \coqref{SECF.Equiv.b:15}{\coqdocvariable{b}} \coqref{SECF.Imp.:com:com scope:'if' x 'then' x 'else' x 'end'}{\coqdocnotation{then}} \coqref{SECF.Equiv.c1:16}{\coqdocvariable{c1}} \coqref{SECF.Imp.:com:com scope:'if' x 'then' x 'else' x 'end'}{\coqdocnotation{else}} \coqref{SECF.Equiv.c2:17}{\coqdocvariable{c2}} \coqref{SECF.Imp.:com:com scope:'if' x 'then' x 'else' x 'end'}{\coqdocnotation{end}} \coqref{SECF.Imp.::com scope:'<x7B' x 'x7D>'}{\coqdocnotation{\}>}}\coqdoceol
\coqdocindent{2.00em}
\coqref{SECF.Equiv.c1:16}{\coqdocvariable{c1}}.\coqdoceol
\coqdocnoindent
\coqdockw{Proof}.\coqdoceol
\coqdocindent{1.00em}
\coqdoctac{intros} \coqdocvar{b} \coqdocvar{c1} \coqdocvar{c2} \coqdocvar{Hb}.\coqdoceol
\coqdocindent{1.00em}
\coqdoctac{split}; \coqdoctac{intros} \coqdocvar{H}.\coqdoceol
\coqdocindent{1.00em}
- \coqdoceol
\coqdocindent{2.00em}
\coqdoctac{inversion} \coqdocvar{H}; \coqdoctac{subst}.\coqdoceol
\coqdocindent{2.00em}
+ \coqdoceol
\coqdocindent{3.00em}
\coqdoctac{assumption}.\coqdoceol
\coqdocindent{2.00em}
+ \coqdoceol
\coqdocindent{3.00em}
\coqdoctac{unfold} \coqref{SECF.Equiv.bequiv}{\coqdocdefinition{bequiv}} \coqdoctac{in} \coqdocvar{Hb}. \coqdoctac{simpl} \coqdoctac{in} \coqdocvar{Hb}.\coqdoceol
\coqdocindent{3.00em}
\coqdoctac{rewrite} \coqdocvar{Hb} \coqdoctac{in} \coqdocvar{H5}.\coqdoceol
\coqdocindent{3.00em}
\coqdoctac{discriminate}.\coqdoceol
\coqdocindent{1.00em}
- \coqdoceol
\coqdocindent{2.00em}
\coqdoctac{apply} \coqref{SECF.Imp.E IfTrue}{\coqdocconstructor{E\_IfTrue}}; \coqdoctac{try} \coqdoctac{assumption}.\coqdoceol
\coqdocindent{2.00em}
\coqdoctac{unfold} \coqref{SECF.Equiv.bequiv}{\coqdocdefinition{bequiv}} \coqdoctac{in} \coqdocvar{Hb}. \coqdoctac{simpl} \coqdoctac{in} \coqdocvar{Hb}.\coqdoceol
\coqdocindent{2.00em}
\coqdoctac{apply} \coqdocvar{Hb}. \coqdockw{Qed}.\coqdoceol
\coqdocemptyline
\end{coqdoccode}
\paragraph{Exercise: 2 stars, standard, especially useful (if\_false)}


\begin{coqdoccode}
\coqdocnoindent
\coqdockw{Theorem} \coqdef{SECF.Equiv.if false}{if\_false}{\coqdoclemma{if\_false}} : \coqdockw{\ensuremath{\forall}} \coqdef{SECF.Equiv.b:18}{b}{\coqdocbinder{b}} \coqdef{SECF.Equiv.c1:19}{c1}{\coqdocbinder{c1}} \coqdef{SECF.Equiv.c2:20}{c2}{\coqdocbinder{c2}},\coqdoceol
\coqdocindent{1.00em}
\coqref{SECF.Equiv.bequiv}{\coqdocdefinition{bequiv}} \coqref{SECF.Equiv.b:18}{\coqdocvariable{b}} \coqref{SECF.Imp.::com scope:'<x7B' x 'x7D>'}{\coqdocnotation{<\{}}\coqref{SECF.Imp.:com::'false'}{\coqdocnotation{false}}\coqref{SECF.Imp.::com scope:'<x7B' x 'x7D>'}{\coqdocnotation{\}>}} \coqexternalref{::type scope:x '->' x}{http://rocq-prover.org/doc/V9.1.1/stdlib//Corelib.Init.Logic}{\coqdocnotation{\ensuremath{\rightarrow}}}\coqdoceol
\coqdocindent{1.00em}
\coqref{SECF.Equiv.cequiv}{\coqdocdefinition{cequiv}}\coqdoceol
\coqdocindent{2.00em}
\coqref{SECF.Imp.::com scope:'<x7B' x 'x7D>'}{\coqdocnotation{<\{}} \coqref{SECF.Imp.:com:com scope:'if' x 'then' x 'else' x 'end'}{\coqdocnotation{if}} \coqref{SECF.Equiv.b:18}{\coqdocvariable{b}} \coqref{SECF.Imp.:com:com scope:'if' x 'then' x 'else' x 'end'}{\coqdocnotation{then}} \coqref{SECF.Equiv.c1:19}{\coqdocvariable{c1}} \coqref{SECF.Imp.:com:com scope:'if' x 'then' x 'else' x 'end'}{\coqdocnotation{else}} \coqref{SECF.Equiv.c2:20}{\coqdocvariable{c2}} \coqref{SECF.Imp.:com:com scope:'if' x 'then' x 'else' x 'end'}{\coqdocnotation{end}} \coqref{SECF.Imp.::com scope:'<x7B' x 'x7D>'}{\coqdocnotation{\}>}}\coqdoceol
\coqdocindent{2.00em}
\coqref{SECF.Equiv.c2:20}{\coqdocvariable{c2}}.\coqdoceol
\coqdocnoindent
\coqdockw{Proof}.\coqdoceol
 \coqdocvar{Admitted}.\coqdoceol
\end{coqdoccode}
\ensuremath{\Box} 

\paragraph{Exercise: 3 stars, standard (swap\_if\_branches)}




    Show that we can swap the branches of an \coqdockw{if} if we also negate its
    condition. 
\begin{coqdoccode}
\coqdocemptyline
\coqdocnoindent
\coqdockw{Theorem} \coqdef{SECF.Equiv.swap if branches}{swap\_if\_branches}{\coqdoclemma{swap\_if\_branches}} : \coqdockw{\ensuremath{\forall}} \coqdef{SECF.Equiv.b:21}{b}{\coqdocbinder{b}} \coqdef{SECF.Equiv.c1:22}{c1}{\coqdocbinder{c1}} \coqdef{SECF.Equiv.c2:23}{c2}{\coqdocbinder{c2}},\coqdoceol
\coqdocindent{1.00em}
\coqref{SECF.Equiv.cequiv}{\coqdocdefinition{cequiv}}\coqdoceol
\coqdocindent{2.00em}
\coqref{SECF.Imp.::com scope:'<x7B' x 'x7D>'}{\coqdocnotation{<\{}} \coqref{SECF.Imp.:com:com scope:'if' x 'then' x 'else' x 'end'}{\coqdocnotation{if}} \coqref{SECF.Equiv.b:21}{\coqdocvariable{b}}   \coqref{SECF.Imp.:com:com scope:'if' x 'then' x 'else' x 'end'}{\coqdocnotation{then}} \coqref{SECF.Equiv.c1:22}{\coqdocvariable{c1}} \coqref{SECF.Imp.:com:com scope:'if' x 'then' x 'else' x 'end'}{\coqdocnotation{else}} \coqref{SECF.Equiv.c2:23}{\coqdocvariable{c2}} \coqref{SECF.Imp.:com:com scope:'if' x 'then' x 'else' x 'end'}{\coqdocnotation{end}} \coqref{SECF.Imp.::com scope:'<x7B' x 'x7D>'}{\coqdocnotation{\}>}}\coqdoceol
\coqdocindent{2.00em}
\coqref{SECF.Imp.::com scope:'<x7B' x 'x7D>'}{\coqdocnotation{<\{}} \coqref{SECF.Imp.:com:com scope:'if' x 'then' x 'else' x 'end'}{\coqdocnotation{if}} \coqref{SECF.Imp.:com::'x7E' x}{\coqdocnotation{\ensuremath{\lnot}}} \coqref{SECF.Equiv.b:21}{\coqdocvariable{b}} \coqref{SECF.Imp.:com:com scope:'if' x 'then' x 'else' x 'end'}{\coqdocnotation{then}} \coqref{SECF.Equiv.c2:23}{\coqdocvariable{c2}} \coqref{SECF.Imp.:com:com scope:'if' x 'then' x 'else' x 'end'}{\coqdocnotation{else}} \coqref{SECF.Equiv.c1:22}{\coqdocvariable{c1}} \coqref{SECF.Imp.:com:com scope:'if' x 'then' x 'else' x 'end'}{\coqdocnotation{end}} \coqref{SECF.Imp.::com scope:'<x7B' x 'x7D>'}{\coqdocnotation{\}>}}.\coqdoceol
\coqdocnoindent
\coqdockw{Proof}.\coqdoceol
 \coqdocvar{Admitted}.\coqdoceol
\end{coqdoccode}
\ensuremath{\Box} 

 For \coqdocvar{while} loops, we can give a similar pair of theorems.  A loop
    whose guard is equivalent to \coqexternalref{false}{http://rocq-prover.org/doc/V9.1.1/stdlib//Corelib.Init.Datatypes}{\coqdocconstructor{false}} is equivalent to \coqdocvar{skip},
    while a loop whose guard is equivalent to \coqexternalref{true}{http://rocq-prover.org/doc/V9.1.1/stdlib//Corelib.Init.Datatypes}{\coqdocconstructor{true}} is equivalent to
    \coqdocvar{while} \coqexternalref{true}{http://rocq-prover.org/doc/V9.1.1/stdlib//Corelib.Init.Datatypes}{\coqdocconstructor{true}} \coqdoctac{do} \coqdocvar{skip} \coqdockw{end} (or any other non-terminating program). 

 The first of these facts is easy. 
\begin{coqdoccode}
\coqdocemptyline
\coqdocnoindent
\coqdockw{Theorem} \coqdef{SECF.Equiv.while false}{while\_false}{\coqdoclemma{while\_false}} : \coqdockw{\ensuremath{\forall}} \coqdef{SECF.Equiv.b:24}{b}{\coqdocbinder{b}} \coqdef{SECF.Equiv.c:25}{c}{\coqdocbinder{c}},\coqdoceol
\coqdocindent{1.00em}
\coqref{SECF.Equiv.bequiv}{\coqdocdefinition{bequiv}} \coqref{SECF.Equiv.b:24}{\coqdocvariable{b}} \coqref{SECF.Imp.::com scope:'<x7B' x 'x7D>'}{\coqdocnotation{<\{}}\coqref{SECF.Imp.:com::'false'}{\coqdocnotation{false}}\coqref{SECF.Imp.::com scope:'<x7B' x 'x7D>'}{\coqdocnotation{\}>}} \coqexternalref{::type scope:x '->' x}{http://rocq-prover.org/doc/V9.1.1/stdlib//Corelib.Init.Logic}{\coqdocnotation{\ensuremath{\rightarrow}}}\coqdoceol
\coqdocindent{1.00em}
\coqref{SECF.Equiv.cequiv}{\coqdocdefinition{cequiv}}\coqdoceol
\coqdocindent{2.00em}
\coqref{SECF.Imp.::com scope:'<x7B' x 'x7D>'}{\coqdocnotation{<\{}} \coqref{SECF.Imp.:com:com scope:'while' x 'do' x 'end'}{\coqdocnotation{while}} \coqref{SECF.Equiv.b:24}{\coqdocvariable{b}} \coqref{SECF.Imp.:com:com scope:'while' x 'do' x 'end'}{\coqdocnotation{do}} \coqref{SECF.Equiv.c:25}{\coqdocvariable{c}} \coqref{SECF.Imp.:com:com scope:'while' x 'do' x 'end'}{\coqdocnotation{end}} \coqref{SECF.Imp.::com scope:'<x7B' x 'x7D>'}{\coqdocnotation{\}>}}\coqdoceol
\coqdocindent{2.00em}
\coqref{SECF.Imp.::com scope:'<x7B' x 'x7D>'}{\coqdocnotation{<\{}} \coqref{SECF.Imp.:com:com scope:'skip'}{\coqdocnotation{skip}} \coqref{SECF.Imp.::com scope:'<x7B' x 'x7D>'}{\coqdocnotation{\}>}}.\coqdoceol
\coqdocnoindent
\coqdockw{Proof}.\coqdoceol
\coqdocindent{1.00em}
\coqdoctac{intros} \coqdocvar{b} \coqdocvar{c} \coqdocvar{Hb}. \coqdoctac{split}; \coqdoctac{intros} \coqdocvar{H}.\coqdoceol
\coqdocindent{1.00em}
- \coqdoceol
\coqdocindent{2.00em}
\coqdoctac{inversion} \coqdocvar{H}; \coqdoctac{subst}.\coqdoceol
\coqdocindent{2.00em}
+ \coqdoceol
\coqdocindent{3.00em}
\coqdoctac{apply} \coqref{SECF.Imp.E Skip}{\coqdocconstructor{E\_Skip}}.\coqdoceol
\coqdocindent{2.00em}
+ \coqdoceol
\coqdocindent{3.00em}
\coqdoctac{rewrite} \coqdocvar{Hb} \coqdoctac{in} \coqdocvar{H2}. \coqdoctac{discriminate}.\coqdoceol
\coqdocindent{1.00em}
- \coqdoceol
\coqdocindent{2.00em}
\coqdoctac{inversion} \coqdocvar{H}. \coqdoctac{subst}.\coqdoceol
\coqdocindent{2.00em}
\coqdoctac{apply} \coqref{SECF.Imp.E WhileFalse}{\coqdocconstructor{E\_WhileFalse}}.\coqdoceol
\coqdocindent{2.00em}
\coqdoctac{apply} \coqdocvar{Hb}. \coqdockw{Qed}.\coqdoceol
\coqdocemptyline
\end{coqdoccode}
\paragraph{Exercise: 2 stars, advanced, optional (while\_false\_informal)}




    Write an informal proof of \coqref{SECF.Equiv.while false}{\coqdoclemma{while\_false}}.


 \ensuremath{\Box} 

 To prove the second fact, we need an auxiliary lemma stating that
    \coqdocvar{while} loops whose guards are equivalent to \coqexternalref{true}{http://rocq-prover.org/doc/V9.1.1/stdlib//Corelib.Init.Datatypes}{\coqdocconstructor{true}} never
    terminate. 

 \textit{Lemma}: If \coqdocvar{b} is equivalent to \coqexternalref{true}{http://rocq-prover.org/doc/V9.1.1/stdlib//Corelib.Init.Datatypes}{\coqdocconstructor{true}}, then it cannot be
    the case that \coqdocvar{st} =[ \coqdocvar{while} \coqdocvar{b} \coqdoctac{do} \coqdocvar{c} \coqdockw{end} ]=> \coqdocvar{st'}.


    \textit{Proof}: Suppose that \coqdocvar{st} =[ \coqdocvar{while} \coqdocvar{b} \coqdoctac{do} \coqdocvar{c} \coqdockw{end} ]=> \coqdocvar{st'}.  We show,
    by induction on a derivation of \coqdocvar{st} =[ \coqdocvar{while} \coqdocvar{b} \coqdoctac{do} \coqdocvar{c} \coqdockw{end} ]=> \coqdocvar{st'},
    that this assumption leads to a contradiction. The only two cases
    to consider are \coqref{SECF.Hoare.HoareAssertAssume.E WhileFalse}{\coqdocconstructor{E\_WhileFalse}} and \coqref{SECF.Hoare.Himp.E WhileTrue}{\coqdocconstructor{E\_WhileTrue}}; the others
    are contradictory.



\begin{itemize}
\item  Suppose \coqdocvar{st} =[ \coqdocvar{while} \coqdocvar{b} \coqdoctac{do} \coqdocvar{c} \coqdockw{end} ]=> \coqdocvar{st'} is proved using rule
      \coqref{SECF.Hoare.HoareAssertAssume.E WhileFalse}{\coqdocconstructor{E\_WhileFalse}}.  Then by assumption \coqref{SECF.Imp.beval}{\coqdocdefinition{beval}} \coqdocvar{st} \coqdocvar{b} = \coqexternalref{false}{http://rocq-prover.org/doc/V9.1.1/stdlib//Corelib.Init.Datatypes}{\coqdocconstructor{false}}.  But
      this contradicts the assumption that \coqdocvar{b} is equivalent to
      \coqexternalref{true}{http://rocq-prover.org/doc/V9.1.1/stdlib//Corelib.Init.Datatypes}{\coqdocconstructor{true}}.



\item  Suppose \coqdocvar{st} =[ \coqdocvar{while} \coqdocvar{b} \coqdoctac{do} \coqdocvar{c} \coqdockw{end} ]=> \coqdocvar{st'} is proved using rule
      \coqref{SECF.Hoare.Himp.E WhileTrue}{\coqdocconstructor{E\_WhileTrue}}.  We must have:


      1. \coqref{SECF.Imp.beval}{\coqdocdefinition{beval}} \coqdocvar{st} \coqdocvar{b} = \coqexternalref{true}{http://rocq-prover.org/doc/V9.1.1/stdlib//Corelib.Init.Datatypes}{\coqdocconstructor{true}}, and
      2. there is some \coqdocvar{st0} such that \coqdocvar{st} =[ \coqdocvar{c} ]=> \coqdocvar{st0} and
         \coqdocvar{st0} =[ \coqdocvar{while} \coqdocvar{b} \coqdoctac{do} \coqdocvar{c} \coqdockw{end} ]=> \coqdocvar{st'}.
      3. Also, we are given an induction hypothesis saying that
         \coqdocvar{st0} =[ \coqdocvar{while} \coqdocvar{b} \coqdoctac{do} \coqdocvar{c} \coqdockw{end} ]=> \coqdocvar{st'} leads to a contradiction,


      We obtain a contradiction by 2 and 3. \ensuremath{\Box} 

\end{itemize}
\begin{coqdoccode}
\coqdocemptyline
\coqdocnoindent
\coqdockw{Lemma} \coqdef{SECF.Equiv.while true nonterm}{while\_true\_nonterm}{\coqdoclemma{while\_true\_nonterm}} : \coqdockw{\ensuremath{\forall}} \coqdef{SECF.Equiv.b:26}{b}{\coqdocbinder{b}} \coqdef{SECF.Equiv.c:27}{c}{\coqdocbinder{c}} \coqdef{SECF.Equiv.st:28}{st}{\coqdocbinder{st}} \coqdef{SECF.Equiv.st':29}{st'}{\coqdocbinder{st'}},\coqdoceol
\coqdocindent{1.00em}
\coqref{SECF.Equiv.bequiv}{\coqdocdefinition{bequiv}} \coqref{SECF.Equiv.b:26}{\coqdocvariable{b}} \coqref{SECF.Imp.::com scope:'<x7B' x 'x7D>'}{\coqdocnotation{<\{}}\coqref{SECF.Imp.:com::'true'}{\coqdocnotation{true}}\coqref{SECF.Imp.::com scope:'<x7B' x 'x7D>'}{\coqdocnotation{\}>}} \coqexternalref{::type scope:x '->' x}{http://rocq-prover.org/doc/V9.1.1/stdlib//Corelib.Init.Logic}{\coqdocnotation{\ensuremath{\rightarrow}}}\coqdoceol
\coqdocindent{1.00em}
\coqexternalref{::type scope:'x7E' x}{http://rocq-prover.org/doc/V9.1.1/stdlib//Corelib.Init.Logic}{\coqdocnotation{\~{}(}} \coqref{SECF.Equiv.st:28}{\coqdocvariable{st}} \coqref{SECF.Imp.:::x '=[' x ']=>' x}{\coqdocnotation{=[}} \coqref{SECF.Imp.:com:com scope:'while' x 'do' x 'end'}{\coqdocnotation{while}} \coqref{SECF.Equiv.b:26}{\coqdocvariable{b}} \coqref{SECF.Imp.:com:com scope:'while' x 'do' x 'end'}{\coqdocnotation{do}} \coqref{SECF.Equiv.c:27}{\coqdocvariable{c}} \coqref{SECF.Imp.:com:com scope:'while' x 'do' x 'end'}{\coqdocnotation{end}} \coqref{SECF.Imp.:::x '=[' x ']=>' x}{\coqdocnotation{]=>}} \coqref{SECF.Equiv.st':29}{\coqdocvariable{st'}} \coqexternalref{::type scope:'x7E' x}{http://rocq-prover.org/doc/V9.1.1/stdlib//Corelib.Init.Logic}{\coqdocnotation{)}}.\coqdoceol
\coqdocnoindent
\coqdockw{Proof}.\coqdoceol
\coqdocindent{1.00em}
\coqdoctac{intros} \coqdocvar{b} \coqdocvar{c} \coqdocvar{st} \coqdocvar{st'} \coqdocvar{Hb}.\coqdoceol
\coqdocindent{1.00em}
\coqdoctac{intros} \coqdocvar{H}.\coqdoceol
\coqdocindent{1.00em}
\coqdocvar{remember} \coqref{SECF.Imp.::com scope:'<x7B' x 'x7D>'}{\coqdocnotation{<\{}} \coqref{SECF.Imp.:com:com scope:'while' x 'do' x 'end'}{\coqdocnotation{while}} \coqdocvar{b} \coqref{SECF.Imp.:com:com scope:'while' x 'do' x 'end'}{\coqdocnotation{do}} \coqdocvar{c} \coqref{SECF.Imp.:com:com scope:'while' x 'do' x 'end'}{\coqdocnotation{end}} \coqref{SECF.Imp.::com scope:'<x7B' x 'x7D>'}{\coqdocnotation{\}>}} \coqdockw{as} \coqdocvar{cw} \coqdocvar{eqn}:\coqdocvar{Heqcw}.\coqdoceol
\coqdocindent{1.00em}
\coqdoctac{induction} \coqdocvar{H};\coqdoceol
\coqdocindent{1.00em}
\coqdoceol
\coqdocindent{1.00em}
\coqdoctac{inversion} \coqdocvar{Heqcw}; \coqdoctac{subst}; \coqdoctac{clear} \coqdocvar{Heqcw}.\coqdoceol
\coqdocindent{1.00em}
-  \coqdoceol
\coqdocindent{2.00em}
\coqdoctac{unfold} \coqref{SECF.Equiv.bequiv}{\coqdocdefinition{bequiv}} \coqdoctac{in} \coqdocvar{Hb}.\coqdoceol
\coqdocindent{2.00em}
\coqdoctac{rewrite} \coqdocvar{Hb} \coqdoctac{in} \coqdocvar{H}. \coqdoctac{discriminate}.\coqdoceol
\coqdocindent{1.00em}
-  \coqdoceol
\coqdocindent{2.00em}
\coqdoctac{apply} \coqdocvar{IHceval2}. \coqdoctac{reflexivity}. \coqdockw{Qed}.\coqdoceol
\coqdocemptyline
\end{coqdoccode}
\paragraph{Exercise: 2 stars, standard, optional (while\_true\_nonterm\_informal)}




    Explain what the lemma \coqref{SECF.Equiv.while true nonterm}{\coqdoclemma{while\_true\_nonterm}} means in English.


 \ensuremath{\Box} 

\paragraph{Exercise: 2 stars, standard, especially useful (while\_true)}




    Prove the following theorem. \textit{Hint}: You'll want to use
    \coqref{SECF.Equiv.while true nonterm}{\coqdoclemma{while\_true\_nonterm}} here. 
\begin{coqdoccode}
\coqdocemptyline
\coqdocnoindent
\coqdockw{Theorem} \coqdef{SECF.Equiv.while true}{while\_true}{\coqdoclemma{while\_true}} : \coqdockw{\ensuremath{\forall}} \coqdef{SECF.Equiv.b:30}{b}{\coqdocbinder{b}} \coqdef{SECF.Equiv.c:31}{c}{\coqdocbinder{c}},\coqdoceol
\coqdocindent{1.00em}
\coqref{SECF.Equiv.bequiv}{\coqdocdefinition{bequiv}} \coqref{SECF.Equiv.b:30}{\coqdocvariable{b}} \coqref{SECF.Imp.::com scope:'<x7B' x 'x7D>'}{\coqdocnotation{<\{}}\coqref{SECF.Imp.:com::'true'}{\coqdocnotation{true}}\coqref{SECF.Imp.::com scope:'<x7B' x 'x7D>'}{\coqdocnotation{\}>}}  \coqexternalref{::type scope:x '->' x}{http://rocq-prover.org/doc/V9.1.1/stdlib//Corelib.Init.Logic}{\coqdocnotation{\ensuremath{\rightarrow}}}\coqdoceol
\coqdocindent{1.00em}
\coqref{SECF.Equiv.cequiv}{\coqdocdefinition{cequiv}}\coqdoceol
\coqdocindent{2.00em}
\coqref{SECF.Imp.::com scope:'<x7B' x 'x7D>'}{\coqdocnotation{<\{}} \coqref{SECF.Imp.:com:com scope:'while' x 'do' x 'end'}{\coqdocnotation{while}} \coqref{SECF.Equiv.b:30}{\coqdocvariable{b}} \coqref{SECF.Imp.:com:com scope:'while' x 'do' x 'end'}{\coqdocnotation{do}} \coqref{SECF.Equiv.c:31}{\coqdocvariable{c}} \coqref{SECF.Imp.:com:com scope:'while' x 'do' x 'end'}{\coqdocnotation{end}} \coqref{SECF.Imp.::com scope:'<x7B' x 'x7D>'}{\coqdocnotation{\}>}}\coqdoceol
\coqdocindent{2.00em}
\coqref{SECF.Imp.::com scope:'<x7B' x 'x7D>'}{\coqdocnotation{<\{}} \coqref{SECF.Imp.:com:com scope:'while' x 'do' x 'end'}{\coqdocnotation{while}} \coqref{SECF.Imp.:com::'true'}{\coqdocnotation{true}} \coqref{SECF.Imp.:com:com scope:'while' x 'do' x 'end'}{\coqdocnotation{do}} \coqref{SECF.Imp.:com:com scope:'skip'}{\coqdocnotation{skip}} \coqref{SECF.Imp.:com:com scope:'while' x 'do' x 'end'}{\coqdocnotation{end}} \coqref{SECF.Imp.::com scope:'<x7B' x 'x7D>'}{\coqdocnotation{\}>}}.\coqdoceol
\coqdocnoindent
\coqdockw{Proof}.\coqdoceol
 \coqdocvar{Admitted}.\coqdoceol
\end{coqdoccode}
\ensuremath{\Box} 

 A more interesting fact about \coqdocvar{while} commands is that any number
    of copies of the body can be ``unrolled'' without changing meaning.


    Loop unrolling is an important transformation in any real
    compiler, so its correctness is of more than just academic
    interest! 
\begin{coqdoccode}
\coqdocemptyline
\coqdocnoindent
\coqdockw{Theorem} \coqdef{SECF.Equiv.loop unrolling}{loop\_unrolling}{\coqdoclemma{loop\_unrolling}} : \coqdockw{\ensuremath{\forall}} \coqdef{SECF.Equiv.b:32}{b}{\coqdocbinder{b}} \coqdef{SECF.Equiv.c:33}{c}{\coqdocbinder{c}},\coqdoceol
\coqdocindent{1.00em}
\coqref{SECF.Equiv.cequiv}{\coqdocdefinition{cequiv}}\coqdoceol
\coqdocindent{2.00em}
\coqref{SECF.Imp.::com scope:'<x7B' x 'x7D>'}{\coqdocnotation{<\{}} \coqref{SECF.Imp.:com:com scope:'while' x 'do' x 'end'}{\coqdocnotation{while}} \coqref{SECF.Equiv.b:32}{\coqdocvariable{b}} \coqref{SECF.Imp.:com:com scope:'while' x 'do' x 'end'}{\coqdocnotation{do}} \coqref{SECF.Equiv.c:33}{\coqdocvariable{c}} \coqref{SECF.Imp.:com:com scope:'while' x 'do' x 'end'}{\coqdocnotation{end}} \coqref{SECF.Imp.::com scope:'<x7B' x 'x7D>'}{\coqdocnotation{\}>}}\coqdoceol
\coqdocindent{2.00em}
\coqref{SECF.Imp.::com scope:'<x7B' x 'x7D>'}{\coqdocnotation{<\{}} \coqref{SECF.Imp.:com:com scope:'if' x 'then' x 'else' x 'end'}{\coqdocnotation{if}} \coqref{SECF.Equiv.b:32}{\coqdocvariable{b}} \coqref{SECF.Imp.:com:com scope:'if' x 'then' x 'else' x 'end'}{\coqdocnotation{then}} \coqref{SECF.Equiv.c:33}{\coqdocvariable{c}} \coqref{SECF.Imp.:com:com scope:x ';' x}{\coqdocnotation{;}} \coqref{SECF.Imp.:com:com scope:'while' x 'do' x 'end'}{\coqdocnotation{while}} \coqref{SECF.Equiv.b:32}{\coqdocvariable{b}} \coqref{SECF.Imp.:com:com scope:'while' x 'do' x 'end'}{\coqdocnotation{do}} \coqref{SECF.Equiv.c:33}{\coqdocvariable{c}} \coqref{SECF.Imp.:com:com scope:'while' x 'do' x 'end'}{\coqdocnotation{end}} \coqref{SECF.Imp.:com:com scope:'if' x 'then' x 'else' x 'end'}{\coqdocnotation{else}} \coqref{SECF.Imp.:com:com scope:'skip'}{\coqdocnotation{skip}} \coqref{SECF.Imp.:com:com scope:'if' x 'then' x 'else' x 'end'}{\coqdocnotation{end}} \coqref{SECF.Imp.::com scope:'<x7B' x 'x7D>'}{\coqdocnotation{\}>}}.\coqdoceol
\coqdocnoindent
\coqdockw{Proof}.\coqdoceol
\coqdocindent{1.00em}
\coqdoctac{intros} \coqdocvar{b} \coqdocvar{c} \coqdocvar{st} \coqdocvar{st'}.\coqdoceol
\coqdocindent{1.00em}
\coqdoctac{split}; \coqdoctac{intros} \coqdocvar{Hce}.\coqdoceol
\coqdocindent{1.00em}
- \coqdoceol
\coqdocindent{2.00em}
\coqdoctac{inversion} \coqdocvar{Hce}; \coqdoctac{subst}.\coqdoceol
\coqdocindent{2.00em}
+ \coqdoceol
\coqdocindent{3.00em}
\coqdoctac{apply} \coqref{SECF.Imp.E IfFalse}{\coqdocconstructor{E\_IfFalse}}.\coqdoceol
\coqdocindent{3.00em}
\ensuremath{\times} \coqdoctac{assumption}.\coqdoceol
\coqdocindent{3.00em}
\ensuremath{\times} \coqdoctac{apply} \coqref{SECF.Imp.E Skip}{\coqdocconstructor{E\_Skip}}.\coqdoceol
\coqdocindent{2.00em}
+ \coqdoceol
\coqdocindent{3.00em}
\coqdoctac{apply} \coqref{SECF.Imp.E IfTrue}{\coqdocconstructor{E\_IfTrue}}.\coqdoceol
\coqdocindent{3.00em}
\ensuremath{\times} \coqdoctac{assumption}.\coqdoceol
\coqdocindent{3.00em}
\ensuremath{\times} \coqdoctac{apply} \coqref{SECF.Imp.E Seq}{\coqdocconstructor{E\_Seq}} \coqdockw{with} (\coqdocvar{st'} := \coqdocvar{st'0}).\coqdoceol
\coqdocindent{4.00em}
-- \coqdoctac{assumption}.\coqdoceol
\coqdocindent{4.00em}
-- \coqdoctac{assumption}.\coqdoceol
\coqdocindent{1.00em}
- \coqdoceol
\coqdocindent{2.00em}
\coqdoctac{inversion} \coqdocvar{Hce}; \coqdoctac{subst}.\coqdoceol
\coqdocindent{2.00em}
+ \coqdoceol
\coqdocindent{3.00em}
\coqdoctac{inversion} \coqdocvar{H5}; \coqdoctac{subst}.\coqdoceol
\coqdocindent{3.00em}
\coqdoctac{apply} \coqref{SECF.Imp.E WhileTrue}{\coqdocconstructor{E\_WhileTrue}} \coqdockw{with} (\coqdocvar{st'} := \coqdocvar{st'0}).\coqdoceol
\coqdocindent{3.00em}
\ensuremath{\times} \coqdoctac{assumption}.\coqdoceol
\coqdocindent{3.00em}
\ensuremath{\times} \coqdoctac{assumption}.\coqdoceol
\coqdocindent{3.00em}
\ensuremath{\times} \coqdoctac{assumption}.\coqdoceol
\coqdocindent{2.00em}
+ \coqdoceol
\coqdocindent{3.00em}
\coqdoctac{inversion} \coqdocvar{H5}; \coqdoctac{subst}. \coqdoctac{apply} \coqref{SECF.Imp.E WhileFalse}{\coqdocconstructor{E\_WhileFalse}}. \coqdoctac{assumption}. \coqdockw{Qed}.\coqdoceol
\coqdocemptyline
\end{coqdoccode}
\paragraph{Exercise: 2 stars, standard, optional (seq\_assoc)}


\begin{coqdoccode}
\coqdocnoindent
\coqdockw{Theorem} \coqdef{SECF.Equiv.seq assoc}{seq\_assoc}{\coqdoclemma{seq\_assoc}} : \coqdockw{\ensuremath{\forall}} \coqdef{SECF.Equiv.c1:34}{c1}{\coqdocbinder{c1}} \coqdef{SECF.Equiv.c2:35}{c2}{\coqdocbinder{c2}} \coqdef{SECF.Equiv.c3:36}{c3}{\coqdocbinder{c3}},\coqdoceol
\coqdocindent{1.00em}
\coqref{SECF.Equiv.cequiv}{\coqdocdefinition{cequiv}} \coqref{SECF.Imp.::com scope:'<x7B' x 'x7D>'}{\coqdocnotation{<\{}}\coqref{SECF.Imp.:com::'(' x ')'}{\coqdocnotation{(}}\coqref{SECF.Equiv.c1:34}{\coqdocvariable{c1}}\coqref{SECF.Imp.:com:com scope:x ';' x}{\coqdocnotation{;}}\coqref{SECF.Equiv.c2:35}{\coqdocvariable{c2}}\coqref{SECF.Imp.:com::'(' x ')'}{\coqdocnotation{)}}\coqref{SECF.Imp.:com:com scope:x ';' x}{\coqdocnotation{;}}\coqref{SECF.Equiv.c3:36}{\coqdocvariable{c3}}\coqref{SECF.Imp.::com scope:'<x7B' x 'x7D>'}{\coqdocnotation{\}>}} \coqref{SECF.Imp.::com scope:'<x7B' x 'x7D>'}{\coqdocnotation{<\{}}\coqref{SECF.Equiv.c1:34}{\coqdocvariable{c1}}\coqref{SECF.Imp.:com:com scope:x ';' x}{\coqdocnotation{;}}\coqref{SECF.Imp.:com::'(' x ')'}{\coqdocnotation{(}}\coqref{SECF.Equiv.c2:35}{\coqdocvariable{c2}}\coqref{SECF.Imp.:com:com scope:x ';' x}{\coqdocnotation{;}}\coqref{SECF.Equiv.c3:36}{\coqdocvariable{c3}}\coqref{SECF.Imp.:com::'(' x ')'}{\coqdocnotation{)}}\coqref{SECF.Imp.::com scope:'<x7B' x 'x7D>'}{\coqdocnotation{\}>}}.\coqdoceol
\coqdocnoindent
\coqdockw{Proof}.\coqdoceol
 \coqdocvar{Admitted}.\coqdoceol
\end{coqdoccode}
\ensuremath{\Box} 

 Proving program properties involving assignments is one place
    where the fact that we are treating equality on program states
    extensionally (e.g., \coqdocvar{x} !-> \coqdocvar{m} \coqdocvar{x} ; \coqdocvar{m} and \coqdocvar{m} are equal maps) comes
    in handy. 
\begin{coqdoccode}
\coqdocemptyline
\coqdocnoindent
\coqdockw{Theorem} \coqdef{SECF.Equiv.identity assignment}{identity\_assignment}{\coqdoclemma{identity\_assignment}} : \coqdockw{\ensuremath{\forall}} \coqdef{SECF.Equiv.x:37}{x}{\coqdocbinder{x}},\coqdoceol
\coqdocindent{1.00em}
\coqref{SECF.Equiv.cequiv}{\coqdocdefinition{cequiv}}\coqdoceol
\coqdocindent{2.00em}
\coqref{SECF.Imp.::com scope:'<x7B' x 'x7D>'}{\coqdocnotation{<\{}} \coqref{SECF.Equiv.x:37}{\coqdocvariable{x}} \coqref{SECF.Imp.:com:com scope:x ':=' x}{\coqdocnotation{:=}} \coqref{SECF.Equiv.x:37}{\coqdocvariable{x}} \coqref{SECF.Imp.::com scope:'<x7B' x 'x7D>'}{\coqdocnotation{\}>}}\coqdoceol
\coqdocindent{2.00em}
\coqref{SECF.Imp.::com scope:'<x7B' x 'x7D>'}{\coqdocnotation{<\{}} \coqref{SECF.Imp.:com:com scope:'skip'}{\coqdocnotation{skip}} \coqref{SECF.Imp.::com scope:'<x7B' x 'x7D>'}{\coqdocnotation{\}>}}.\coqdoceol
\coqdocnoindent
\coqdockw{Proof}.\coqdoceol
\coqdocindent{1.00em}
\coqdoctac{intros}.\coqdoceol
\coqdocindent{1.00em}
\coqdoctac{split}; \coqdoctac{intro} \coqdocvar{H}; \coqdoctac{inversion} \coqdocvar{H}; \coqdoctac{subst}; \coqdoctac{clear} \coqdocvar{H}.\coqdoceol
\coqdocindent{1.00em}
- \coqdoceol
\coqdocindent{2.00em}
\coqdoctac{rewrite} \coqref{SECF.Maps.t update same}{\coqdocaxiom{t\_update\_same}}.\coqdoceol
\coqdocindent{2.00em}
\coqdoctac{apply} \coqref{SECF.Imp.E Skip}{\coqdocconstructor{E\_Skip}}.\coqdoceol
\coqdocindent{1.00em}
- \coqdoceol
\coqdocindent{2.00em}
\coqdoctac{assert} (\coqdocvar{Hx} : \coqdocvar{st'} \coqref{SECF.Imp.:::x '=[' x ']=>' x}{\coqdocnotation{=[}} \coqdocvar{x} \coqref{SECF.Imp.:com:com scope:x ':=' x}{\coqdocnotation{:=}} \coqdocvar{x} \coqref{SECF.Imp.:::x '=[' x ']=>' x}{\coqdocnotation{]=>}} \coqref{SECF.Imp.:::x '=[' x ']=>' x}{\coqdocnotation{(}}\coqdocvar{x} \coqref{SECF.Maps.:::x '!->' x ';' x}{\coqdocnotation{!->}} \coqdocvar{st'} \coqdocvar{x} \coqref{SECF.Maps.:::x '!->' x ';' x}{\coqdocnotation{;}} \coqdocvar{st'}\coqref{SECF.Imp.:::x '=[' x ']=>' x}{\coqdocnotation{)}}).\coqdoceol
\coqdocindent{2.00em}
\{ \coqdoctac{apply} \coqref{SECF.Imp.E Asgn}{\coqdocconstructor{E\_Asgn}}. \coqdoctac{reflexivity}. \}\coqdoceol
\coqdocindent{2.00em}
\coqdoctac{rewrite} \coqref{SECF.Maps.t update same}{\coqdocaxiom{t\_update\_same}} \coqdoctac{in} \coqdocvar{Hx}.\coqdoceol
\coqdocindent{2.00em}
\coqdoctac{apply} \coqdocvar{Hx}.\coqdoceol
\coqdocnoindent
\coqdockw{Qed}.\coqdoceol
\coqdocemptyline
\end{coqdoccode}
\paragraph{Exercise: 2 stars, standard, especially useful (assign\_aequiv)}


\begin{coqdoccode}
\coqdocnoindent
\coqdockw{Theorem} \coqdef{SECF.Equiv.assign aequiv}{assign\_aequiv}{\coqdoclemma{assign\_aequiv}} : \coqdockw{\ensuremath{\forall}} (\coqdef{SECF.Equiv.X:38}{X}{\coqdocbinder{X}} : \coqdocinductive{string}) (\coqdef{SECF.Equiv.a:39}{a}{\coqdocbinder{a}} : \coqref{SECF.Imp.aexp}{\coqdocinductive{aexp}}),\coqdoceol
\coqdocindent{1.00em}
\coqref{SECF.Equiv.aequiv}{\coqdocdefinition{aequiv}} \coqref{SECF.Imp.::com scope:'<x7B' x 'x7D>'}{\coqdocnotation{<\{}} \coqref{SECF.Equiv.X:38}{\coqdocvariable{X}} \coqref{SECF.Imp.::com scope:'<x7B' x 'x7D>'}{\coqdocnotation{\}>}} \coqref{SECF.Equiv.a:39}{\coqdocvariable{a}} \coqexternalref{::type scope:x '->' x}{http://rocq-prover.org/doc/V9.1.1/stdlib//Corelib.Init.Logic}{\coqdocnotation{\ensuremath{\rightarrow}}}\coqdoceol
\coqdocindent{1.00em}
\coqref{SECF.Equiv.cequiv}{\coqdocdefinition{cequiv}} \coqref{SECF.Imp.::com scope:'<x7B' x 'x7D>'}{\coqdocnotation{<\{}} \coqref{SECF.Imp.:com:com scope:'skip'}{\coqdocnotation{skip}} \coqref{SECF.Imp.::com scope:'<x7B' x 'x7D>'}{\coqdocnotation{\}>}} \coqref{SECF.Imp.::com scope:'<x7B' x 'x7D>'}{\coqdocnotation{<\{}} \coqref{SECF.Equiv.X:38}{\coqdocvariable{X}} \coqref{SECF.Imp.:com:com scope:x ':=' x}{\coqdocnotation{:=}} \coqref{SECF.Equiv.a:39}{\coqdocvariable{a}} \coqref{SECF.Imp.::com scope:'<x7B' x 'x7D>'}{\coqdocnotation{\}>}}.\coqdoceol
\coqdocnoindent
\coqdockw{Proof}.\coqdoceol
 \coqdocvar{Admitted}.\coqdoceol
\end{coqdoccode}
\ensuremath{\Box} 

\paragraph{Exercise: 2 stars, standard, optional (equiv\_classes)}



 Given the following programs, group together those that are
    equivalent in Imp. Your answer should be given as a list of lists,
    where each sub-list represents a group of equivalent programs. For
    example, if you think programs (a) through (h) are all equivalent
    to each other, but not to (i), your answer should look like this:


        [\coqref{SECF.Equiv.prog a}{\coqdocdefinition{prog\_a}};\coqref{SECF.Equiv.prog b}{\coqdocdefinition{prog\_b}};\coqref{SECF.Equiv.prog c}{\coqdocdefinition{prog\_c}};\coqref{SECF.Equiv.prog d}{\coqdocdefinition{prog\_d}};\coqref{SECF.Equiv.prog e}{\coqdocdefinition{prog\_e}};\coqref{SECF.Equiv.prog f}{\coqdocdefinition{prog\_f}};\coqref{SECF.Equiv.prog g}{\coqdocdefinition{prog\_g}};\coqref{SECF.Equiv.prog h}{\coqdocdefinition{prog\_h}}] ;
         [\coqref{SECF.Equiv.prog i}{\coqdocdefinition{prog\_i}}] 


    Write down your answer below in the definition of
    \coqref{SECF.Equiv.equiv classes}{\coqdocdefinition{equiv\_classes}}. 
\begin{coqdoccode}
\coqdocemptyline
\coqdocnoindent
\coqdockw{Definition} \coqdef{SECF.Equiv.prog a}{prog\_a}{\coqdocdefinition{prog\_a}} : \coqref{SECF.Imp.com}{\coqdocinductive{com}} :=\coqdoceol
\coqdocindent{1.00em}
\coqref{SECF.Imp.::com scope:'<x7B' x 'x7D>'}{\coqdocnotation{<\{}} \coqref{SECF.Imp.:com:com scope:'while' x 'do' x 'end'}{\coqdocnotation{while}} \coqref{SECF.Imp.X}{\coqdocdefinition{X}} \coqref{SECF.Imp.:com::x '>' x}{\coqdocnotation{>}} 0 \coqref{SECF.Imp.:com:com scope:'while' x 'do' x 'end'}{\coqdocnotation{do}}\coqdoceol
\coqdocindent{3.50em}
\coqref{SECF.Imp.X}{\coqdocdefinition{X}} \coqref{SECF.Imp.:com:com scope:x ':=' x}{\coqdocnotation{:=}} \coqref{SECF.Imp.X}{\coqdocdefinition{X}} \coqref{SECF.Imp.:com::x '+' x}{\coqdocnotation{+}} 1\coqdoceol
\coqdocindent{2.50em}
\coqref{SECF.Imp.:com:com scope:'while' x 'do' x 'end'}{\coqdocnotation{end}} \coqref{SECF.Imp.::com scope:'<x7B' x 'x7D>'}{\coqdocnotation{\}>}}.\coqdoceol
\coqdocemptyline
\coqdocnoindent
\coqdockw{Definition} \coqdef{SECF.Equiv.prog b}{prog\_b}{\coqdocdefinition{prog\_b}} : \coqref{SECF.Imp.com}{\coqdocinductive{com}} :=\coqdoceol
\coqdocindent{1.00em}
\coqref{SECF.Imp.::com scope:'<x7B' x 'x7D>'}{\coqdocnotation{<\{}} \coqref{SECF.Imp.:com:com scope:'if' x 'then' x 'else' x 'end'}{\coqdocnotation{if}} \coqref{SECF.Imp.:com::'(' x ')'}{\coqdocnotation{(}}\coqref{SECF.Imp.X}{\coqdocdefinition{X}} \coqref{SECF.Imp.:com::x '=' x}{\coqdocnotation{=}} 0\coqref{SECF.Imp.:com::'(' x ')'}{\coqdocnotation{)}} \coqref{SECF.Imp.:com:com scope:'if' x 'then' x 'else' x 'end'}{\coqdocnotation{then}}\coqdoceol
\coqdocindent{3.50em}
\coqref{SECF.Imp.X}{\coqdocdefinition{X}} \coqref{SECF.Imp.:com:com scope:x ':=' x}{\coqdocnotation{:=}} \coqref{SECF.Imp.X}{\coqdocdefinition{X}} \coqref{SECF.Imp.:com::x '+' x}{\coqdocnotation{+}} 1\coqref{SECF.Imp.:com:com scope:x ';' x}{\coqdocnotation{;}}\coqdoceol
\coqdocindent{3.50em}
\coqref{SECF.Imp.Y}{\coqdocdefinition{Y}} \coqref{SECF.Imp.:com:com scope:x ':=' x}{\coqdocnotation{:=}} 1\coqdoceol
\coqdocindent{2.50em}
\coqref{SECF.Imp.:com:com scope:'if' x 'then' x 'else' x 'end'}{\coqdocnotation{else}}\coqdoceol
\coqdocindent{3.50em}
\coqref{SECF.Imp.Y}{\coqdocdefinition{Y}} \coqref{SECF.Imp.:com:com scope:x ':=' x}{\coqdocnotation{:=}} 0\coqdoceol
\coqdocindent{2.50em}
\coqref{SECF.Imp.:com:com scope:'if' x 'then' x 'else' x 'end'}{\coqdocnotation{end}}\coqref{SECF.Imp.:com:com scope:x ';' x}{\coqdocnotation{;}}\coqdoceol
\coqdocindent{2.50em}
\coqref{SECF.Imp.X}{\coqdocdefinition{X}} \coqref{SECF.Imp.:com:com scope:x ':=' x}{\coqdocnotation{:=}} \coqref{SECF.Imp.X}{\coqdocdefinition{X}} \coqref{SECF.Imp.:com::x '-' x}{\coqdocnotation{-}} \coqref{SECF.Imp.Y}{\coqdocdefinition{Y}}\coqref{SECF.Imp.:com:com scope:x ';' x}{\coqdocnotation{;}}\coqdoceol
\coqdocindent{2.50em}
\coqref{SECF.Imp.Y}{\coqdocdefinition{Y}} \coqref{SECF.Imp.:com:com scope:x ':=' x}{\coqdocnotation{:=}} 0 \coqref{SECF.Imp.::com scope:'<x7B' x 'x7D>'}{\coqdocnotation{\}>}}.\coqdoceol
\coqdocemptyline
\coqdocnoindent
\coqdockw{Definition} \coqdef{SECF.Equiv.prog c}{prog\_c}{\coqdocdefinition{prog\_c}} : \coqref{SECF.Imp.com}{\coqdocinductive{com}} :=\coqdoceol
\coqdocindent{1.00em}
\coqref{SECF.Imp.::com scope:'<x7B' x 'x7D>'}{\coqdocnotation{<\{}} \coqref{SECF.Imp.:com:com scope:'skip'}{\coqdocnotation{skip}} \coqref{SECF.Imp.::com scope:'<x7B' x 'x7D>'}{\coqdocnotation{\}>}} .\coqdoceol
\coqdocemptyline
\coqdocnoindent
\coqdockw{Definition} \coqdef{SECF.Equiv.prog d}{prog\_d}{\coqdocdefinition{prog\_d}} : \coqref{SECF.Imp.com}{\coqdocinductive{com}} :=\coqdoceol
\coqdocindent{1.00em}
\coqref{SECF.Imp.::com scope:'<x7B' x 'x7D>'}{\coqdocnotation{<\{}} \coqref{SECF.Imp.:com:com scope:'while' x 'do' x 'end'}{\coqdocnotation{while}} \coqref{SECF.Imp.X}{\coqdocdefinition{X}} \coqref{SECF.Imp.:com::x '<>' x}{\coqdocnotation{\ensuremath{\not=}}} 0 \coqref{SECF.Imp.:com:com scope:'while' x 'do' x 'end'}{\coqdocnotation{do}}\coqdoceol
\coqdocindent{3.50em}
\coqref{SECF.Imp.X}{\coqdocdefinition{X}} \coqref{SECF.Imp.:com:com scope:x ':=' x}{\coqdocnotation{:=}} \coqref{SECF.Imp.:com::'(' x ')'}{\coqdocnotation{(}}\coqref{SECF.Imp.X}{\coqdocdefinition{X}} \coqref{SECF.Imp.:com::x '*' x}{\coqdocnotation{\ensuremath{\times}}} \coqref{SECF.Imp.Y}{\coqdocdefinition{Y}}\coqref{SECF.Imp.:com::'(' x ')'}{\coqdocnotation{)}} \coqref{SECF.Imp.:com::x '+' x}{\coqdocnotation{+}} 1\coqdoceol
\coqdocindent{2.50em}
\coqref{SECF.Imp.:com:com scope:'while' x 'do' x 'end'}{\coqdocnotation{end}} \coqref{SECF.Imp.::com scope:'<x7B' x 'x7D>'}{\coqdocnotation{\}>}}.\coqdoceol
\coqdocemptyline
\coqdocnoindent
\coqdockw{Definition} \coqdef{SECF.Equiv.prog e}{prog\_e}{\coqdocdefinition{prog\_e}} : \coqref{SECF.Imp.com}{\coqdocinductive{com}} :=\coqdoceol
\coqdocindent{1.00em}
\coqref{SECF.Imp.::com scope:'<x7B' x 'x7D>'}{\coqdocnotation{<\{}} \coqref{SECF.Imp.Y}{\coqdocdefinition{Y}} \coqref{SECF.Imp.:com:com scope:x ':=' x}{\coqdocnotation{:=}} 0 \coqref{SECF.Imp.::com scope:'<x7B' x 'x7D>'}{\coqdocnotation{\}>}}.\coqdoceol
\coqdocemptyline
\coqdocnoindent
\coqdockw{Definition} \coqdef{SECF.Equiv.prog f}{prog\_f}{\coqdocdefinition{prog\_f}} : \coqref{SECF.Imp.com}{\coqdocinductive{com}} :=\coqdoceol
\coqdocindent{1.00em}
\coqref{SECF.Imp.::com scope:'<x7B' x 'x7D>'}{\coqdocnotation{<\{}} \coqref{SECF.Imp.Y}{\coqdocdefinition{Y}} \coqref{SECF.Imp.:com:com scope:x ':=' x}{\coqdocnotation{:=}} \coqref{SECF.Imp.X}{\coqdocdefinition{X}} \coqref{SECF.Imp.:com::x '+' x}{\coqdocnotation{+}} 1\coqref{SECF.Imp.:com:com scope:x ';' x}{\coqdocnotation{;}}\coqdoceol
\coqdocindent{2.50em}
\coqref{SECF.Imp.:com:com scope:'while' x 'do' x 'end'}{\coqdocnotation{while}} \coqref{SECF.Imp.X}{\coqdocdefinition{X}} \coqref{SECF.Imp.:com::x '<>' x}{\coqdocnotation{\ensuremath{\not=}}} \coqref{SECF.Imp.Y}{\coqdocdefinition{Y}} \coqref{SECF.Imp.:com:com scope:'while' x 'do' x 'end'}{\coqdocnotation{do}}\coqdoceol
\coqdocindent{3.50em}
\coqref{SECF.Imp.Y}{\coqdocdefinition{Y}} \coqref{SECF.Imp.:com:com scope:x ':=' x}{\coqdocnotation{:=}} \coqref{SECF.Imp.X}{\coqdocdefinition{X}} \coqref{SECF.Imp.:com::x '+' x}{\coqdocnotation{+}} 1\coqdoceol
\coqdocindent{2.50em}
\coqref{SECF.Imp.:com:com scope:'while' x 'do' x 'end'}{\coqdocnotation{end}} \coqref{SECF.Imp.::com scope:'<x7B' x 'x7D>'}{\coqdocnotation{\}>}}.\coqdoceol
\coqdocemptyline
\coqdocnoindent
\coqdockw{Definition} \coqdef{SECF.Equiv.prog g}{prog\_g}{\coqdocdefinition{prog\_g}} : \coqref{SECF.Imp.com}{\coqdocinductive{com}} :=\coqdoceol
\coqdocindent{1.00em}
\coqref{SECF.Imp.::com scope:'<x7B' x 'x7D>'}{\coqdocnotation{<\{}} \coqref{SECF.Imp.:com:com scope:'while' x 'do' x 'end'}{\coqdocnotation{while}} \coqref{SECF.Imp.:com::'true'}{\coqdocnotation{true}} \coqref{SECF.Imp.:com:com scope:'while' x 'do' x 'end'}{\coqdocnotation{do}}\coqdoceol
\coqdocindent{3.50em}
\coqref{SECF.Imp.:com:com scope:'skip'}{\coqdocnotation{skip}}\coqdoceol
\coqdocindent{2.50em}
\coqref{SECF.Imp.:com:com scope:'while' x 'do' x 'end'}{\coqdocnotation{end}} \coqref{SECF.Imp.::com scope:'<x7B' x 'x7D>'}{\coqdocnotation{\}>}}.\coqdoceol
\coqdocemptyline
\coqdocnoindent
\coqdockw{Definition} \coqdef{SECF.Equiv.prog h}{prog\_h}{\coqdocdefinition{prog\_h}} : \coqref{SECF.Imp.com}{\coqdocinductive{com}} :=\coqdoceol
\coqdocindent{1.00em}
\coqref{SECF.Imp.::com scope:'<x7B' x 'x7D>'}{\coqdocnotation{<\{}} \coqref{SECF.Imp.:com:com scope:'while' x 'do' x 'end'}{\coqdocnotation{while}} \coqref{SECF.Imp.X}{\coqdocdefinition{X}} \coqref{SECF.Imp.:com::x '<>' x}{\coqdocnotation{\ensuremath{\not=}}} \coqref{SECF.Imp.X}{\coqdocdefinition{X}} \coqref{SECF.Imp.:com:com scope:'while' x 'do' x 'end'}{\coqdocnotation{do}}\coqdoceol
\coqdocindent{3.50em}
\coqref{SECF.Imp.X}{\coqdocdefinition{X}} \coqref{SECF.Imp.:com:com scope:x ':=' x}{\coqdocnotation{:=}} \coqref{SECF.Imp.X}{\coqdocdefinition{X}} \coqref{SECF.Imp.:com::x '+' x}{\coqdocnotation{+}} 1\coqdoceol
\coqdocindent{2.50em}
\coqref{SECF.Imp.:com:com scope:'while' x 'do' x 'end'}{\coqdocnotation{end}} \coqref{SECF.Imp.::com scope:'<x7B' x 'x7D>'}{\coqdocnotation{\}>}}.\coqdoceol
\coqdocemptyline
\coqdocnoindent
\coqdockw{Definition} \coqdef{SECF.Equiv.prog i}{prog\_i}{\coqdocdefinition{prog\_i}} : \coqref{SECF.Imp.com}{\coqdocinductive{com}} :=\coqdoceol
\coqdocindent{1.00em}
\coqref{SECF.Imp.::com scope:'<x7B' x 'x7D>'}{\coqdocnotation{<\{}} \coqref{SECF.Imp.:com:com scope:'while' x 'do' x 'end'}{\coqdocnotation{while}} \coqref{SECF.Imp.X}{\coqdocdefinition{X}} \coqref{SECF.Imp.:com::x '<>' x}{\coqdocnotation{\ensuremath{\not=}}} \coqref{SECF.Imp.Y}{\coqdocdefinition{Y}} \coqref{SECF.Imp.:com:com scope:'while' x 'do' x 'end'}{\coqdocnotation{do}}\coqdoceol
\coqdocindent{3.50em}
\coqref{SECF.Imp.X}{\coqdocdefinition{X}} \coqref{SECF.Imp.:com:com scope:x ':=' x}{\coqdocnotation{:=}} \coqref{SECF.Imp.Y}{\coqdocdefinition{Y}} \coqref{SECF.Imp.:com::x '+' x}{\coqdocnotation{+}} 1\coqdoceol
\coqdocindent{2.50em}
\coqref{SECF.Imp.:com:com scope:'while' x 'do' x 'end'}{\coqdocnotation{end}} \coqref{SECF.Imp.::com scope:'<x7B' x 'x7D>'}{\coqdocnotation{\}>}}.\coqdoceol
\coqdocemptyline
\coqdocnoindent
\coqdockw{Definition} \coqdef{SECF.Equiv.equiv classes}{equiv\_classes}{\coqdocdefinition{equiv\_classes}} : \coqexternalref{list}{http://rocq-prover.org/doc/V9.1.1/stdlib//Corelib.Init.Datatypes}{\coqdocinductive{list}} (\coqexternalref{list}{http://rocq-prover.org/doc/V9.1.1/stdlib//Corelib.Init.Datatypes}{\coqdocinductive{list}} \coqref{SECF.Imp.com}{\coqdocinductive{com}})\coqdoceol
\coqdocindent{1.00em}
. \coqdocvar{Admitted}.\coqdoceol
\coqdocemptyline
\coqdocnoindent
\coqdockw{Definition} \coqdef{SECF.Equiv.manual grade for equiv classes}{manual\_grade\_for\_equiv\_classes}{\coqdocdefinition{manual\_grade\_for\_equiv\_classes}} : \coqexternalref{option}{http://rocq-prover.org/doc/V9.1.1/stdlib//Corelib.Init.Datatypes}{\coqdocinductive{option}} (\coqexternalref{nat}{http://rocq-prover.org/doc/V9.1.1/stdlib//Corelib.Init.Datatypes}{\coqdocinductive{nat}}\coqexternalref{::type scope:x '*' x}{http://rocq-prover.org/doc/V9.1.1/stdlib//Corelib.Init.Datatypes}{\coqdocnotation{\ensuremath{\times}}}\coqdocinductive{string}) := \coqexternalref{None}{http://rocq-prover.org/doc/V9.1.1/stdlib//Corelib.Init.Datatypes}{\coqdocconstructor{None}}.\coqdoceol
\end{coqdoccode}
\ensuremath{\Box} 
\begin{coqdoccode}
\coqdocemptyline
\end{coqdoccode}
\section{Properties of Behavioral Equivalence}



 We next consider some fundamental properties of program
    equivalence. 
\begin{coqdoccode}
\coqdocemptyline
\end{coqdoccode}
\subsection{Behavioral Equivalence is an Equivalence}



 First, let's verify that the equivalences on \coqdocvar{aexps}, \coqdocvar{bexps}, and
    \coqref{SECF.SpecCT.Div.com}{\coqdocinductive{com}}s really are \textit{equivalences} -- i.e., that they are reflexive,
    symmetric, and transitive.  The proofs are all easy. 
\begin{coqdoccode}
\coqdocemptyline
\coqdocnoindent
\coqdockw{Lemma} \coqdef{SECF.Equiv.refl aequiv}{refl\_aequiv}{\coqdoclemma{refl\_aequiv}} : \coqdockw{\ensuremath{\forall}} (\coqdef{SECF.Equiv.a:40}{a}{\coqdocbinder{a}} : \coqref{SECF.Imp.aexp}{\coqdocinductive{aexp}}),\coqdoceol
\coqdocindent{1.00em}
\coqref{SECF.Equiv.aequiv}{\coqdocdefinition{aequiv}} \coqref{SECF.Equiv.a:40}{\coqdocvariable{a}} \coqref{SECF.Equiv.a:40}{\coqdocvariable{a}}.\coqdoceol
\coqdocnoindent
\coqdockw{Proof}.\coqdoceol
\coqdocindent{1.00em}
\coqdoctac{intros} \coqdocvar{a} \coqdocvar{st}. \coqdoctac{reflexivity}. \coqdockw{Qed}.\coqdoceol
\coqdocemptyline
\coqdocnoindent
\coqdockw{Lemma} \coqdef{SECF.Equiv.sym aequiv}{sym\_aequiv}{\coqdoclemma{sym\_aequiv}} : \coqdockw{\ensuremath{\forall}} (\coqdef{SECF.Equiv.a1:41}{a1}{\coqdocbinder{a1}} \coqdef{SECF.Equiv.a2:42}{a2}{\coqdocbinder{a2}} : \coqref{SECF.Imp.aexp}{\coqdocinductive{aexp}}),\coqdoceol
\coqdocindent{1.00em}
\coqref{SECF.Equiv.aequiv}{\coqdocdefinition{aequiv}} \coqref{SECF.Equiv.a1:41}{\coqdocvariable{a1}} \coqref{SECF.Equiv.a2:42}{\coqdocvariable{a2}} \coqexternalref{::type scope:x '->' x}{http://rocq-prover.org/doc/V9.1.1/stdlib//Corelib.Init.Logic}{\coqdocnotation{\ensuremath{\rightarrow}}} \coqref{SECF.Equiv.aequiv}{\coqdocdefinition{aequiv}} \coqref{SECF.Equiv.a2:42}{\coqdocvariable{a2}} \coqref{SECF.Equiv.a1:41}{\coqdocvariable{a1}}.\coqdoceol
\coqdocnoindent
\coqdockw{Proof}.\coqdoceol
\coqdocindent{1.00em}
\coqdoctac{intros} \coqdocvar{a1} \coqdocvar{a2} \coqdocvar{H}. \coqdoctac{intros} \coqdocvar{st}. \coqdoctac{symmetry}. \coqdoctac{apply} \coqdocvar{H}. \coqdockw{Qed}.\coqdoceol
\coqdocemptyline
\coqdocnoindent
\coqdockw{Lemma} \coqdef{SECF.Equiv.trans aequiv}{trans\_aequiv}{\coqdoclemma{trans\_aequiv}} : \coqdockw{\ensuremath{\forall}} (\coqdef{SECF.Equiv.a1:43}{a1}{\coqdocbinder{a1}} \coqdef{SECF.Equiv.a2:44}{a2}{\coqdocbinder{a2}} \coqdef{SECF.Equiv.a3:45}{a3}{\coqdocbinder{a3}} : \coqref{SECF.Imp.aexp}{\coqdocinductive{aexp}}),\coqdoceol
\coqdocindent{1.00em}
\coqref{SECF.Equiv.aequiv}{\coqdocdefinition{aequiv}} \coqref{SECF.Equiv.a1:43}{\coqdocvariable{a1}} \coqref{SECF.Equiv.a2:44}{\coqdocvariable{a2}} \coqexternalref{::type scope:x '->' x}{http://rocq-prover.org/doc/V9.1.1/stdlib//Corelib.Init.Logic}{\coqdocnotation{\ensuremath{\rightarrow}}} \coqref{SECF.Equiv.aequiv}{\coqdocdefinition{aequiv}} \coqref{SECF.Equiv.a2:44}{\coqdocvariable{a2}} \coqref{SECF.Equiv.a3:45}{\coqdocvariable{a3}} \coqexternalref{::type scope:x '->' x}{http://rocq-prover.org/doc/V9.1.1/stdlib//Corelib.Init.Logic}{\coqdocnotation{\ensuremath{\rightarrow}}} \coqref{SECF.Equiv.aequiv}{\coqdocdefinition{aequiv}} \coqref{SECF.Equiv.a1:43}{\coqdocvariable{a1}} \coqref{SECF.Equiv.a3:45}{\coqdocvariable{a3}}.\coqdoceol
\coqdocnoindent
\coqdockw{Proof}.\coqdoceol
\coqdocindent{1.00em}
\coqdoctac{unfold} \coqref{SECF.Equiv.aequiv}{\coqdocdefinition{aequiv}}. \coqdoctac{intros} \coqdocvar{a1} \coqdocvar{a2} \coqdocvar{a3} \coqdocvar{H12} \coqdocvar{H23} \coqdocvar{st}.\coqdoceol
\coqdocindent{1.00em}
\coqdoctac{rewrite} (\coqdocvar{H12} \coqdocvar{st}). \coqdoctac{rewrite} (\coqdocvar{H23} \coqdocvar{st}). \coqdoctac{reflexivity}. \coqdockw{Qed}.\coqdoceol
\coqdocemptyline
\coqdocnoindent
\coqdockw{Lemma} \coqdef{SECF.Equiv.refl bequiv}{refl\_bequiv}{\coqdoclemma{refl\_bequiv}} : \coqdockw{\ensuremath{\forall}} (\coqdef{SECF.Equiv.b:46}{b}{\coqdocbinder{b}} : \coqref{SECF.Imp.bexp}{\coqdocinductive{bexp}}),\coqdoceol
\coqdocindent{1.00em}
\coqref{SECF.Equiv.bequiv}{\coqdocdefinition{bequiv}} \coqref{SECF.Equiv.b:46}{\coqdocvariable{b}} \coqref{SECF.Equiv.b:46}{\coqdocvariable{b}}.\coqdoceol
\coqdocnoindent
\coqdockw{Proof}.\coqdoceol
\coqdocindent{1.00em}
\coqdoctac{unfold} \coqref{SECF.Equiv.bequiv}{\coqdocdefinition{bequiv}}. \coqdoctac{intros} \coqdocvar{b} \coqdocvar{st}. \coqdoctac{reflexivity}. \coqdockw{Qed}.\coqdoceol
\coqdocemptyline
\coqdocnoindent
\coqdockw{Lemma} \coqdef{SECF.Equiv.sym bequiv}{sym\_bequiv}{\coqdoclemma{sym\_bequiv}} : \coqdockw{\ensuremath{\forall}} (\coqdef{SECF.Equiv.b1:47}{b1}{\coqdocbinder{b1}} \coqdef{SECF.Equiv.b2:48}{b2}{\coqdocbinder{b2}} : \coqref{SECF.Imp.bexp}{\coqdocinductive{bexp}}),\coqdoceol
\coqdocindent{1.00em}
\coqref{SECF.Equiv.bequiv}{\coqdocdefinition{bequiv}} \coqref{SECF.Equiv.b1:47}{\coqdocvariable{b1}} \coqref{SECF.Equiv.b2:48}{\coqdocvariable{b2}} \coqexternalref{::type scope:x '->' x}{http://rocq-prover.org/doc/V9.1.1/stdlib//Corelib.Init.Logic}{\coqdocnotation{\ensuremath{\rightarrow}}} \coqref{SECF.Equiv.bequiv}{\coqdocdefinition{bequiv}} \coqref{SECF.Equiv.b2:48}{\coqdocvariable{b2}} \coqref{SECF.Equiv.b1:47}{\coqdocvariable{b1}}.\coqdoceol
\coqdocnoindent
\coqdockw{Proof}.\coqdoceol
\coqdocindent{1.00em}
\coqdoctac{unfold} \coqref{SECF.Equiv.bequiv}{\coqdocdefinition{bequiv}}. \coqdoctac{intros} \coqdocvar{b1} \coqdocvar{b2} \coqdocvar{H}. \coqdoctac{intros} \coqdocvar{st}. \coqdoctac{symmetry}. \coqdoctac{apply} \coqdocvar{H}. \coqdockw{Qed}.\coqdoceol
\coqdocemptyline
\coqdocnoindent
\coqdockw{Lemma} \coqdef{SECF.Equiv.trans bequiv}{trans\_bequiv}{\coqdoclemma{trans\_bequiv}} : \coqdockw{\ensuremath{\forall}} (\coqdef{SECF.Equiv.b1:49}{b1}{\coqdocbinder{b1}} \coqdef{SECF.Equiv.b2:50}{b2}{\coqdocbinder{b2}} \coqdef{SECF.Equiv.b3:51}{b3}{\coqdocbinder{b3}} : \coqref{SECF.Imp.bexp}{\coqdocinductive{bexp}}),\coqdoceol
\coqdocindent{1.00em}
\coqref{SECF.Equiv.bequiv}{\coqdocdefinition{bequiv}} \coqref{SECF.Equiv.b1:49}{\coqdocvariable{b1}} \coqref{SECF.Equiv.b2:50}{\coqdocvariable{b2}} \coqexternalref{::type scope:x '->' x}{http://rocq-prover.org/doc/V9.1.1/stdlib//Corelib.Init.Logic}{\coqdocnotation{\ensuremath{\rightarrow}}} \coqref{SECF.Equiv.bequiv}{\coqdocdefinition{bequiv}} \coqref{SECF.Equiv.b2:50}{\coqdocvariable{b2}} \coqref{SECF.Equiv.b3:51}{\coqdocvariable{b3}} \coqexternalref{::type scope:x '->' x}{http://rocq-prover.org/doc/V9.1.1/stdlib//Corelib.Init.Logic}{\coqdocnotation{\ensuremath{\rightarrow}}} \coqref{SECF.Equiv.bequiv}{\coqdocdefinition{bequiv}} \coqref{SECF.Equiv.b1:49}{\coqdocvariable{b1}} \coqref{SECF.Equiv.b3:51}{\coqdocvariable{b3}}.\coqdoceol
\coqdocnoindent
\coqdockw{Proof}.\coqdoceol
\coqdocindent{1.00em}
\coqdoctac{unfold} \coqref{SECF.Equiv.bequiv}{\coqdocdefinition{bequiv}}. \coqdoctac{intros} \coqdocvar{b1} \coqdocvar{b2} \coqdocvar{b3} \coqdocvar{H12} \coqdocvar{H23} \coqdocvar{st}.\coqdoceol
\coqdocindent{1.00em}
\coqdoctac{rewrite} (\coqdocvar{H12} \coqdocvar{st}). \coqdoctac{rewrite} (\coqdocvar{H23} \coqdocvar{st}). \coqdoctac{reflexivity}. \coqdockw{Qed}.\coqdoceol
\coqdocemptyline
\coqdocnoindent
\coqdockw{Lemma} \coqdef{SECF.Equiv.refl cequiv}{refl\_cequiv}{\coqdoclemma{refl\_cequiv}} : \coqdockw{\ensuremath{\forall}} (\coqdef{SECF.Equiv.c:52}{c}{\coqdocbinder{c}} : \coqref{SECF.Imp.com}{\coqdocinductive{com}}),\coqdoceol
\coqdocindent{1.00em}
\coqref{SECF.Equiv.cequiv}{\coqdocdefinition{cequiv}} \coqref{SECF.Equiv.c:52}{\coqdocvariable{c}} \coqref{SECF.Equiv.c:52}{\coqdocvariable{c}}.\coqdoceol
\coqdocnoindent
\coqdockw{Proof}.\coqdoceol
\coqdocindent{1.00em}
\coqdoctac{unfold} \coqref{SECF.Equiv.cequiv}{\coqdocdefinition{cequiv}}. \coqdoctac{intros} \coqdocvar{c} \coqdocvar{st} \coqdocvar{st'}. \coqdoctac{reflexivity}. \coqdockw{Qed}.\coqdoceol
\coqdocemptyline
\coqdocnoindent
\coqdockw{Lemma} \coqdef{SECF.Equiv.sym cequiv}{sym\_cequiv}{\coqdoclemma{sym\_cequiv}} : \coqdockw{\ensuremath{\forall}} (\coqdef{SECF.Equiv.c1:53}{c1}{\coqdocbinder{c1}} \coqdef{SECF.Equiv.c2:54}{c2}{\coqdocbinder{c2}} : \coqref{SECF.Imp.com}{\coqdocinductive{com}}),\coqdoceol
\coqdocindent{1.00em}
\coqref{SECF.Equiv.cequiv}{\coqdocdefinition{cequiv}} \coqref{SECF.Equiv.c1:53}{\coqdocvariable{c1}} \coqref{SECF.Equiv.c2:54}{\coqdocvariable{c2}} \coqexternalref{::type scope:x '->' x}{http://rocq-prover.org/doc/V9.1.1/stdlib//Corelib.Init.Logic}{\coqdocnotation{\ensuremath{\rightarrow}}} \coqref{SECF.Equiv.cequiv}{\coqdocdefinition{cequiv}} \coqref{SECF.Equiv.c2:54}{\coqdocvariable{c2}} \coqref{SECF.Equiv.c1:53}{\coqdocvariable{c1}}.\coqdoceol
\coqdocnoindent
\coqdockw{Proof}.\coqdoceol
\coqdocindent{1.00em}
\coqdoctac{unfold} \coqref{SECF.Equiv.cequiv}{\coqdocdefinition{cequiv}}. \coqdoctac{intros} \coqdocvar{c1} \coqdocvar{c2} \coqdocvar{H} \coqdocvar{st} \coqdocvar{st'}.\coqdoceol
\coqdocindent{1.00em}
\coqdoctac{rewrite} \coqdocvar{H}. \coqdoctac{reflexivity}.\coqdoceol
\coqdocnoindent
\coqdockw{Qed}.\coqdoceol
\coqdocemptyline
\coqdocnoindent
\coqdockw{Lemma} \coqdef{SECF.Equiv.trans cequiv}{trans\_cequiv}{\coqdoclemma{trans\_cequiv}} : \coqdockw{\ensuremath{\forall}} (\coqdef{SECF.Equiv.c1:55}{c1}{\coqdocbinder{c1}} \coqdef{SECF.Equiv.c2:56}{c2}{\coqdocbinder{c2}} \coqdef{SECF.Equiv.c3:57}{c3}{\coqdocbinder{c3}} : \coqref{SECF.Imp.com}{\coqdocinductive{com}}),\coqdoceol
\coqdocindent{1.00em}
\coqref{SECF.Equiv.cequiv}{\coqdocdefinition{cequiv}} \coqref{SECF.Equiv.c1:55}{\coqdocvariable{c1}} \coqref{SECF.Equiv.c2:56}{\coqdocvariable{c2}} \coqexternalref{::type scope:x '->' x}{http://rocq-prover.org/doc/V9.1.1/stdlib//Corelib.Init.Logic}{\coqdocnotation{\ensuremath{\rightarrow}}} \coqref{SECF.Equiv.cequiv}{\coqdocdefinition{cequiv}} \coqref{SECF.Equiv.c2:56}{\coqdocvariable{c2}} \coqref{SECF.Equiv.c3:57}{\coqdocvariable{c3}} \coqexternalref{::type scope:x '->' x}{http://rocq-prover.org/doc/V9.1.1/stdlib//Corelib.Init.Logic}{\coqdocnotation{\ensuremath{\rightarrow}}} \coqref{SECF.Equiv.cequiv}{\coqdocdefinition{cequiv}} \coqref{SECF.Equiv.c1:55}{\coqdocvariable{c1}} \coqref{SECF.Equiv.c3:57}{\coqdocvariable{c3}}.\coqdoceol
\coqdocnoindent
\coqdockw{Proof}.\coqdoceol
\coqdocindent{1.00em}
\coqdoctac{unfold} \coqref{SECF.Equiv.cequiv}{\coqdocdefinition{cequiv}}. \coqdoctac{intros} \coqdocvar{c1} \coqdocvar{c2} \coqdocvar{c3} \coqdocvar{H12} \coqdocvar{H23} \coqdocvar{st} \coqdocvar{st'}.\coqdoceol
\coqdocindent{1.00em}
\coqdoctac{rewrite} \coqdocvar{H12}. \coqdoctac{apply} \coqdocvar{H23}.\coqdoceol
\coqdocnoindent
\coqdockw{Qed}.\coqdoceol
\coqdocemptyline
\end{coqdoccode}
\subsection{Behavioral Equivalence Is a Congruence}



 Less obviously, behavioral equivalence is also a \textit{congruence}.
    That is, the equivalence of two subprograms implies the
    equivalence of the larger programs in which they are embedded:


                 aequiv a a'
\par
\noindent\hrulefill\par
\noindent{}
         cequiv (x := a) (x := a')


              cequiv c1 c1'
              cequiv c2 c2'
\par
\noindent\hrulefill\par
\noindent{}
         cequiv (c1;c2) (c1';c2')


    ... and so on for the other forms of commands. 

 (Note that we are using the inference rule notation here not
    as part of an inductive definition, but simply to write down some
    valid implications in a readable format. We prove these
    implications below.) 

 We will see a concrete example of why these congruence
    properties are important in the following section (in the proof of
    \coqref{SECF.Equiv.fold constants com sound}{\coqdoclemma{fold\_constants\_com\_sound}}), but the main idea is that they allow
    us to replace a small part of a large program with an equivalent
    small part and know that the whole large programs are equivalent
    \textit{without} doing an explicit proof about the parts that didn't
    change -- i.e., the ``proof burden'' of a small change to a large
    program is proportional to the size of the change, not the
    program! 
\begin{coqdoccode}
\coqdocemptyline
\coqdocnoindent
\coqdockw{Theorem} \coqdef{SECF.Equiv.CAsgn congruence}{CAsgn\_congruence}{\coqdoclemma{CAsgn\_congruence}} : \coqdockw{\ensuremath{\forall}} \coqdef{SECF.Equiv.x:58}{x}{\coqdocbinder{x}} \coqdef{SECF.Equiv.a:59}{a}{\coqdocbinder{a}} \coqdef{SECF.Equiv.a':60}{a'}{\coqdocbinder{a'}},\coqdoceol
\coqdocindent{1.00em}
\coqref{SECF.Equiv.aequiv}{\coqdocdefinition{aequiv}} \coqref{SECF.Equiv.a:59}{\coqdocvariable{a}} \coqref{SECF.Equiv.a':60}{\coqdocvariable{a'}} \coqexternalref{::type scope:x '->' x}{http://rocq-prover.org/doc/V9.1.1/stdlib//Corelib.Init.Logic}{\coqdocnotation{\ensuremath{\rightarrow}}}\coqdoceol
\coqdocindent{1.00em}
\coqref{SECF.Equiv.cequiv}{\coqdocdefinition{cequiv}} \coqref{SECF.Imp.::com scope:'<x7B' x 'x7D>'}{\coqdocnotation{<\{}}\coqref{SECF.Equiv.x:58}{\coqdocvariable{x}} \coqref{SECF.Imp.:com:com scope:x ':=' x}{\coqdocnotation{:=}} \coqref{SECF.Equiv.a:59}{\coqdocvariable{a}}\coqref{SECF.Imp.::com scope:'<x7B' x 'x7D>'}{\coqdocnotation{\}>}} \coqref{SECF.Imp.::com scope:'<x7B' x 'x7D>'}{\coqdocnotation{<\{}}\coqref{SECF.Equiv.x:58}{\coqdocvariable{x}} \coqref{SECF.Imp.:com:com scope:x ':=' x}{\coqdocnotation{:=}} \coqref{SECF.Equiv.a':60}{\coqdocvariable{a'}}\coqref{SECF.Imp.::com scope:'<x7B' x 'x7D>'}{\coqdocnotation{\}>}}.\coqdoceol
\coqdocnoindent
\coqdockw{Proof}.\coqdoceol
\coqdocindent{1.00em}
\coqdoctac{intros} \coqdocvar{x} \coqdocvar{a} \coqdocvar{a'} \coqdocvar{Heqv} \coqdocvar{st} \coqdocvar{st'}.\coqdoceol
\coqdocindent{1.00em}
\coqdoctac{split}; \coqdoctac{intros} \coqdocvar{Hceval}.\coqdoceol
\coqdocindent{1.00em}
- \coqdoceol
\coqdocindent{2.00em}
\coqdoctac{inversion} \coqdocvar{Hceval}. \coqdoctac{subst}. \coqdoctac{apply} \coqref{SECF.Imp.E Asgn}{\coqdocconstructor{E\_Asgn}}.\coqdoceol
\coqdocindent{2.00em}
\coqdoctac{rewrite} \coqdocvar{Heqv}. \coqdoctac{reflexivity}.\coqdoceol
\coqdocindent{1.00em}
- \coqdoceol
\coqdocindent{2.00em}
\coqdoctac{inversion} \coqdocvar{Hceval}. \coqdoctac{subst}. \coqdoctac{apply} \coqref{SECF.Imp.E Asgn}{\coqdocconstructor{E\_Asgn}}.\coqdoceol
\coqdocindent{2.00em}
\coqdoctac{rewrite} \coqdocvar{Heqv}. \coqdoctac{reflexivity}. \coqdockw{Qed}.\coqdoceol
\coqdocemptyline
\end{coqdoccode}
The congruence property for loops is a little more interesting,
    since it requires induction.


    \textit{Theorem}: Equivalence is a congruence for \coqdocvar{while} -- that is, if
    \coqdocvar{b} is equivalent to \coqdocvar{b'} and \coqdocvar{c} is equivalent to \coqdocvar{c'}, then
    \coqdocvar{while} \coqdocvar{b} \coqdoctac{do} \coqdocvar{c} \coqdockw{end} is equivalent to \coqdocvar{while} \coqdocvar{b'} \coqdoctac{do} \coqdocvar{c'} \coqdockw{end}.


    \textit{Proof}: Suppose \coqdocvar{b} is equivalent to \coqdocvar{b'} and \coqdocvar{c} is
    equivalent to \coqdocvar{c'}.  We must show, for every \coqdocvar{st} and \coqdocvar{st'}, that
    \coqdocvar{st} =[ \coqdocvar{while} \coqdocvar{b} \coqdoctac{do} \coqdocvar{c} \coqdockw{end} ]=> \coqdocvar{st'} iff \coqdocvar{st} =[ \coqdocvar{while} \coqdocvar{b'} \coqdoctac{do} \coqdocvar{c'}
    \coqdockw{end} ]=> \coqdocvar{st'}.  We consider the two directions separately.



\begin{itemize}
\item  (\ensuremath{\rightarrow}) We show that \coqdocvar{st} =[ \coqdocvar{while} \coqdocvar{b} \coqdoctac{do} \coqdocvar{c} \coqdockw{end} ]=> \coqdocvar{st'} implies
        \coqdocvar{st} =[ \coqdocvar{while} \coqdocvar{b'} \coqdoctac{do} \coqdocvar{c'} \coqdockw{end} ]=> \coqdocvar{st'}, by induction on a
        derivation of \coqdocvar{st} =[ \coqdocvar{while} \coqdocvar{b} \coqdoctac{do} \coqdocvar{c} \coqdockw{end} ]=> \coqdocvar{st'}.  The only
        nontrivial cases are when the final rule in the derivation is
        \coqref{SECF.Hoare.HoareAssertAssume.E WhileFalse}{\coqdocconstructor{E\_WhileFalse}} or \coqref{SECF.Hoare.Himp.E WhileTrue}{\coqdocconstructor{E\_WhileTrue}}.



\begin{itemize}
\item  \coqref{SECF.Hoare.HoareAssertAssume.E WhileFalse}{\coqdocconstructor{E\_WhileFalse}}: In this case, the form of the rule gives us
            \coqref{SECF.Imp.beval}{\coqdocdefinition{beval}} \coqdocvar{st} \coqdocvar{b} = \coqexternalref{false}{http://rocq-prover.org/doc/V9.1.1/stdlib//Corelib.Init.Datatypes}{\coqdocconstructor{false}} and \coqdocvar{st} = \coqdocvar{st'}.  But then, since
            \coqdocvar{b} and \coqdocvar{b'} are equivalent, we have \coqref{SECF.Imp.beval}{\coqdocdefinition{beval}} \coqdocvar{st} \coqdocvar{b'} =
            \coqexternalref{false}{http://rocq-prover.org/doc/V9.1.1/stdlib//Corelib.Init.Datatypes}{\coqdocconstructor{false}}, and \coqref{SECF.Hoare.HoareAssertAssume.E WhileFalse}{\coqdocconstructor{E\_WhileFalse}} applies, giving us
            \coqdocvar{st} =[ \coqdocvar{while} \coqdocvar{b'} \coqdoctac{do} \coqdocvar{c'} \coqdockw{end} ]=> \coqdocvar{st'}, as required.



\item  \coqref{SECF.Hoare.Himp.E WhileTrue}{\coqdocconstructor{E\_WhileTrue}}: The form of the rule now gives us \coqref{SECF.Imp.beval}{\coqdocdefinition{beval}} \coqdocvar{st}
            \coqdocvar{b} = \coqexternalref{true}{http://rocq-prover.org/doc/V9.1.1/stdlib//Corelib.Init.Datatypes}{\coqdocconstructor{true}}, with \coqdocvar{st} =[ \coqdocvar{c} ]=> \coqdocvar{st'0} and \coqdocvar{st'0} =[ \coqdocvar{while}
            \coqdocvar{b} \coqdoctac{do} \coqdocvar{c} \coqdockw{end} ]=> \coqdocvar{st'} for some state \coqdocvar{st'0}, with the
            induction hypothesis \coqdocvar{st'0} =[ \coqdocvar{while} \coqdocvar{b'} \coqdoctac{do} \coqdocvar{c'} \coqdockw{end} ]=>
            \coqdocvar{st'}.


            Since \coqdocvar{c} and \coqdocvar{c'} are equivalent, we know that \coqdocvar{st} =[
            \coqdocvar{c'} ]=> \coqdocvar{st'0}.  And since \coqdocvar{b} and \coqdocvar{b'} are equivalent,
            we have \coqref{SECF.Imp.beval}{\coqdocdefinition{beval}} \coqdocvar{st} \coqdocvar{b'} = \coqexternalref{true}{http://rocq-prover.org/doc/V9.1.1/stdlib//Corelib.Init.Datatypes}{\coqdocconstructor{true}}.  Now \coqref{SECF.Hoare.Himp.E WhileTrue}{\coqdocconstructor{E\_WhileTrue}} applies,
            giving us \coqdocvar{st} =[ \coqdocvar{while} \coqdocvar{b'} \coqdoctac{do} \coqdocvar{c'} \coqdockw{end} ]=> \coqdocvar{st'}, as
            required.



\end{itemize}

\item  (\ensuremath{\leftarrow}) Similar. \ensuremath{\Box} 

\end{itemize}
\begin{coqdoccode}
\coqdocemptyline
\coqdocnoindent
\coqdockw{Theorem} \coqdef{SECF.Equiv.CWhile congruence}{CWhile\_congruence}{\coqdoclemma{CWhile\_congruence}} : \coqdockw{\ensuremath{\forall}} \coqdef{SECF.Equiv.b:61}{b}{\coqdocbinder{b}} \coqdef{SECF.Equiv.b':62}{b'}{\coqdocbinder{b'}} \coqdef{SECF.Equiv.c:63}{c}{\coqdocbinder{c}} \coqdef{SECF.Equiv.c':64}{c'}{\coqdocbinder{c'}},\coqdoceol
\coqdocindent{1.00em}
\coqref{SECF.Equiv.bequiv}{\coqdocdefinition{bequiv}} \coqref{SECF.Equiv.b:61}{\coqdocvariable{b}} \coqref{SECF.Equiv.b':62}{\coqdocvariable{b'}} \coqexternalref{::type scope:x '->' x}{http://rocq-prover.org/doc/V9.1.1/stdlib//Corelib.Init.Logic}{\coqdocnotation{\ensuremath{\rightarrow}}} \coqref{SECF.Equiv.cequiv}{\coqdocdefinition{cequiv}} \coqref{SECF.Equiv.c:63}{\coqdocvariable{c}} \coqref{SECF.Equiv.c':64}{\coqdocvariable{c'}} \coqexternalref{::type scope:x '->' x}{http://rocq-prover.org/doc/V9.1.1/stdlib//Corelib.Init.Logic}{\coqdocnotation{\ensuremath{\rightarrow}}}\coqdoceol
\coqdocindent{1.00em}
\coqref{SECF.Equiv.cequiv}{\coqdocdefinition{cequiv}} \coqref{SECF.Imp.::com scope:'<x7B' x 'x7D>'}{\coqdocnotation{<\{}} \coqref{SECF.Imp.:com:com scope:'while' x 'do' x 'end'}{\coqdocnotation{while}} \coqref{SECF.Equiv.b:61}{\coqdocvariable{b}} \coqref{SECF.Imp.:com:com scope:'while' x 'do' x 'end'}{\coqdocnotation{do}} \coqref{SECF.Equiv.c:63}{\coqdocvariable{c}} \coqref{SECF.Imp.:com:com scope:'while' x 'do' x 'end'}{\coqdocnotation{end}} \coqref{SECF.Imp.::com scope:'<x7B' x 'x7D>'}{\coqdocnotation{\}>}} \coqref{SECF.Imp.::com scope:'<x7B' x 'x7D>'}{\coqdocnotation{<\{}} \coqref{SECF.Imp.:com:com scope:'while' x 'do' x 'end'}{\coqdocnotation{while}} \coqref{SECF.Equiv.b':62}{\coqdocvariable{b'}} \coqref{SECF.Imp.:com:com scope:'while' x 'do' x 'end'}{\coqdocnotation{do}} \coqref{SECF.Equiv.c':64}{\coqdocvariable{c'}} \coqref{SECF.Imp.:com:com scope:'while' x 'do' x 'end'}{\coqdocnotation{end}} \coqref{SECF.Imp.::com scope:'<x7B' x 'x7D>'}{\coqdocnotation{\}>}}.\coqdoceol
\coqdocnoindent
\coqdockw{Proof}.\coqdoceol
\coqdocemptyline
\coqdocindent{1.00em}
\coqdoctac{assert} (\coqdocvar{A}: \coqdockw{\ensuremath{\forall}} (\coqdef{SECF.Equiv.b:65}{b}{\coqdocbinder{b}} \coqdef{SECF.Equiv.b':66}{b'}{\coqdocbinder{b'}} : \coqref{SECF.Imp.bexp}{\coqdocinductive{bexp}}) (\coqdef{SECF.Equiv.c:67}{c}{\coqdocbinder{c}} \coqdef{SECF.Equiv.c':68}{c'}{\coqdocbinder{c'}} : \coqref{SECF.Imp.com}{\coqdocinductive{com}}) (\coqdef{SECF.Equiv.st:69}{st}{\coqdocbinder{st}} \coqdef{SECF.Equiv.st':70}{st'}{\coqdocbinder{st'}} : \coqref{SECF.Imp.state}{\coqdocdefinition{state}}),\coqdoceol
\coqdocindent{6.50em}
\coqref{SECF.Equiv.bequiv}{\coqdocdefinition{bequiv}} \coqref{SECF.Equiv.b:65}{\coqdocvariable{b}} \coqref{SECF.Equiv.b':66}{\coqdocvariable{b'}} \coqexternalref{::type scope:x '->' x}{http://rocq-prover.org/doc/V9.1.1/stdlib//Corelib.Init.Logic}{\coqdocnotation{\ensuremath{\rightarrow}}} \coqref{SECF.Equiv.cequiv}{\coqdocdefinition{cequiv}} \coqref{SECF.Equiv.c:67}{\coqdocvariable{c}} \coqref{SECF.Equiv.c':68}{\coqdocvariable{c'}} \coqexternalref{::type scope:x '->' x}{http://rocq-prover.org/doc/V9.1.1/stdlib//Corelib.Init.Logic}{\coqdocnotation{\ensuremath{\rightarrow}}}\coqdoceol
\coqdocindent{6.50em}
\coqref{SECF.Equiv.st:69}{\coqdocvariable{st}} \coqref{SECF.Imp.:::x '=[' x ']=>' x}{\coqdocnotation{=[}} \coqref{SECF.Imp.:com:com scope:'while' x 'do' x 'end'}{\coqdocnotation{while}} \coqref{SECF.Equiv.b:65}{\coqdocvariable{b}} \coqref{SECF.Imp.:com:com scope:'while' x 'do' x 'end'}{\coqdocnotation{do}} \coqref{SECF.Equiv.c:67}{\coqdocvariable{c}} \coqref{SECF.Imp.:com:com scope:'while' x 'do' x 'end'}{\coqdocnotation{end}} \coqref{SECF.Imp.:::x '=[' x ']=>' x}{\coqdocnotation{]=>}} \coqref{SECF.Equiv.st':70}{\coqdocvariable{st'}} \coqexternalref{::type scope:x '->' x}{http://rocq-prover.org/doc/V9.1.1/stdlib//Corelib.Init.Logic}{\coqdocnotation{\ensuremath{\rightarrow}}}\coqdoceol
\coqdocindent{6.50em}
\coqref{SECF.Equiv.st:69}{\coqdocvariable{st}} \coqref{SECF.Imp.:::x '=[' x ']=>' x}{\coqdocnotation{=[}} \coqref{SECF.Imp.:com:com scope:'while' x 'do' x 'end'}{\coqdocnotation{while}} \coqref{SECF.Equiv.b':66}{\coqdocvariable{b'}} \coqref{SECF.Imp.:com:com scope:'while' x 'do' x 'end'}{\coqdocnotation{do}} \coqref{SECF.Equiv.c':68}{\coqdocvariable{c'}} \coqref{SECF.Imp.:com:com scope:'while' x 'do' x 'end'}{\coqdocnotation{end}} \coqref{SECF.Imp.:::x '=[' x ']=>' x}{\coqdocnotation{]=>}} \coqref{SECF.Equiv.st':70}{\coqdocvariable{st'}}).\coqdoceol
\coqdocindent{1.00em}
\{ \coqdoctac{unfold} \coqref{SECF.Equiv.bequiv}{\coqdocdefinition{bequiv}},\coqref{SECF.Equiv.cequiv}{\coqdocdefinition{cequiv}}.\coqdoceol
\coqdocindent{2.00em}
\coqdoctac{intros} \coqdocvar{b} \coqdocvar{b'} \coqdocvar{c} \coqdocvar{c'} \coqdocvar{st} \coqdocvar{st'} \coqdocvar{Hbe} \coqdocvar{Hc1e} \coqdocvar{Hce}.\coqdoceol
\coqdocindent{2.00em}
\coqdocvar{remember} \coqref{SECF.Imp.::com scope:'<x7B' x 'x7D>'}{\coqdocnotation{<\{}} \coqref{SECF.Imp.:com:com scope:'while' x 'do' x 'end'}{\coqdocnotation{while}} \coqdocvar{b} \coqref{SECF.Imp.:com:com scope:'while' x 'do' x 'end'}{\coqdocnotation{do}} \coqdocvar{c} \coqref{SECF.Imp.:com:com scope:'while' x 'do' x 'end'}{\coqdocnotation{end}} \coqref{SECF.Imp.::com scope:'<x7B' x 'x7D>'}{\coqdocnotation{\}>}} \coqdockw{as} \coqdocvar{cwhile}\coqdoceol
\coqdocindent{3.00em}
\coqdocvar{eqn}:\coqdocvar{Heqcwhile}.\coqdoceol
\coqdocindent{2.00em}
\coqdoctac{induction} \coqdocvar{Hce}; \coqdoctac{inversion} \coqdocvar{Heqcwhile}; \coqdoctac{subst}.\coqdoceol
\coqdocindent{2.00em}
+ \coqdoceol
\coqdocindent{3.00em}
\coqdoctac{apply} \coqref{SECF.Imp.E WhileFalse}{\coqdocconstructor{E\_WhileFalse}}. \coqdoctac{rewrite} \ensuremath{\leftarrow} \coqdocvar{Hbe}. \coqdoctac{apply} \coqdocvar{H}.\coqdoceol
\coqdocindent{2.00em}
+ \coqdoceol
\coqdocindent{3.00em}
\coqdoctac{apply} \coqref{SECF.Imp.E WhileTrue}{\coqdocconstructor{E\_WhileTrue}} \coqdockw{with} (\coqdocvar{st'} := \coqdocvar{st'}).\coqdoceol
\coqdocindent{3.00em}
\ensuremath{\times}  \coqdoctac{rewrite} \ensuremath{\leftarrow} \coqdocvar{Hbe}. \coqdoctac{apply} \coqdocvar{H}.\coqdoceol
\coqdocindent{3.00em}
\ensuremath{\times} \coqdoceol
\coqdocindent{4.00em}
\coqdoctac{apply} (\coqdocvar{Hc1e} \coqdocvar{st} \coqdocvar{st'}). \coqdoctac{apply} \coqdocvar{Hce1}.\coqdoceol
\coqdocindent{3.00em}
\ensuremath{\times} \coqdoceol
\coqdocindent{4.00em}
\coqdoctac{apply} \coqdocvar{IHHce2}. \coqdoctac{reflexivity}. \}\coqdoceol
\coqdocnoindent
\coqdoceol
\coqdocindent{1.00em}
\coqdoctac{intros}. \coqdoctac{split}.\coqdoceol
\coqdocindent{1.00em}
- \coqdoctac{apply} \coqdocvar{A}; \coqdoctac{assumption}.\coqdoceol
\coqdocindent{1.00em}
- \coqdoctac{apply} \coqdocvar{A}.\coqdoceol
\coqdocindent{2.00em}
+ \coqdoctac{apply} \coqref{SECF.Equiv.sym bequiv}{\coqdoclemma{sym\_bequiv}}. \coqdoctac{assumption}.\coqdoceol
\coqdocindent{2.00em}
+ \coqdoctac{apply} \coqref{SECF.Equiv.sym cequiv}{\coqdoclemma{sym\_cequiv}}. \coqdoctac{assumption}.\coqdoceol
\coqdocnoindent
\coqdockw{Qed}.\coqdoceol
\coqdocemptyline
\end{coqdoccode}
\paragraph{Exercise: 3 stars, standard, optional (CSeq\_congruence)}


\begin{coqdoccode}
\coqdocnoindent
\coqdockw{Theorem} \coqdef{SECF.Equiv.CSeq congruence}{CSeq\_congruence}{\coqdoclemma{CSeq\_congruence}} : \coqdockw{\ensuremath{\forall}} \coqdef{SECF.Equiv.c1:77}{c1}{\coqdocbinder{c1}} \coqdef{SECF.Equiv.c1':78}{c1'}{\coqdocbinder{c1'}} \coqdef{SECF.Equiv.c2:79}{c2}{\coqdocbinder{c2}} \coqdef{SECF.Equiv.c2':80}{c2'}{\coqdocbinder{c2'}},\coqdoceol
\coqdocindent{1.00em}
\coqref{SECF.Equiv.cequiv}{\coqdocdefinition{cequiv}} \coqref{SECF.Equiv.c1:77}{\coqdocvariable{c1}} \coqref{SECF.Equiv.c1':78}{\coqdocvariable{c1'}} \coqexternalref{::type scope:x '->' x}{http://rocq-prover.org/doc/V9.1.1/stdlib//Corelib.Init.Logic}{\coqdocnotation{\ensuremath{\rightarrow}}} \coqref{SECF.Equiv.cequiv}{\coqdocdefinition{cequiv}} \coqref{SECF.Equiv.c2:79}{\coqdocvariable{c2}} \coqref{SECF.Equiv.c2':80}{\coqdocvariable{c2'}} \coqexternalref{::type scope:x '->' x}{http://rocq-prover.org/doc/V9.1.1/stdlib//Corelib.Init.Logic}{\coqdocnotation{\ensuremath{\rightarrow}}}\coqdoceol
\coqdocindent{1.00em}
\coqref{SECF.Equiv.cequiv}{\coqdocdefinition{cequiv}} \coqref{SECF.Imp.::com scope:'<x7B' x 'x7D>'}{\coqdocnotation{<\{}} \coqref{SECF.Equiv.c1:77}{\coqdocvariable{c1}}\coqref{SECF.Imp.:com:com scope:x ';' x}{\coqdocnotation{;}}\coqref{SECF.Equiv.c2:79}{\coqdocvariable{c2}} \coqref{SECF.Imp.::com scope:'<x7B' x 'x7D>'}{\coqdocnotation{\}>}} \coqref{SECF.Imp.::com scope:'<x7B' x 'x7D>'}{\coqdocnotation{<\{}} \coqref{SECF.Equiv.c1':78}{\coqdocvariable{c1'}}\coqref{SECF.Imp.:com:com scope:x ';' x}{\coqdocnotation{;}}\coqref{SECF.Equiv.c2':80}{\coqdocvariable{c2'}} \coqref{SECF.Imp.::com scope:'<x7B' x 'x7D>'}{\coqdocnotation{\}>}}.\coqdoceol
\coqdocnoindent
\coqdockw{Proof}.\coqdoceol
 \coqdocvar{Admitted}.\coqdoceol
\end{coqdoccode}
\ensuremath{\Box} 

\paragraph{Exercise: 3 stars, standard (CIf\_congruence)}


\begin{coqdoccode}
\coqdocnoindent
\coqdockw{Theorem} \coqdef{SECF.Equiv.CIf congruence}{CIf\_congruence}{\coqdoclemma{CIf\_congruence}} : \coqdockw{\ensuremath{\forall}} \coqdef{SECF.Equiv.b:81}{b}{\coqdocbinder{b}} \coqdef{SECF.Equiv.b':82}{b'}{\coqdocbinder{b'}} \coqdef{SECF.Equiv.c1:83}{c1}{\coqdocbinder{c1}} \coqdef{SECF.Equiv.c1':84}{c1'}{\coqdocbinder{c1'}} \coqdef{SECF.Equiv.c2:85}{c2}{\coqdocbinder{c2}} \coqdef{SECF.Equiv.c2':86}{c2'}{\coqdocbinder{c2'}},\coqdoceol
\coqdocindent{1.00em}
\coqref{SECF.Equiv.bequiv}{\coqdocdefinition{bequiv}} \coqref{SECF.Equiv.b:81}{\coqdocvariable{b}} \coqref{SECF.Equiv.b':82}{\coqdocvariable{b'}} \coqexternalref{::type scope:x '->' x}{http://rocq-prover.org/doc/V9.1.1/stdlib//Corelib.Init.Logic}{\coqdocnotation{\ensuremath{\rightarrow}}} \coqref{SECF.Equiv.cequiv}{\coqdocdefinition{cequiv}} \coqref{SECF.Equiv.c1:83}{\coqdocvariable{c1}} \coqref{SECF.Equiv.c1':84}{\coqdocvariable{c1'}} \coqexternalref{::type scope:x '->' x}{http://rocq-prover.org/doc/V9.1.1/stdlib//Corelib.Init.Logic}{\coqdocnotation{\ensuremath{\rightarrow}}} \coqref{SECF.Equiv.cequiv}{\coqdocdefinition{cequiv}} \coqref{SECF.Equiv.c2:85}{\coqdocvariable{c2}} \coqref{SECF.Equiv.c2':86}{\coqdocvariable{c2'}} \coqexternalref{::type scope:x '->' x}{http://rocq-prover.org/doc/V9.1.1/stdlib//Corelib.Init.Logic}{\coqdocnotation{\ensuremath{\rightarrow}}}\coqdoceol
\coqdocindent{1.00em}
\coqref{SECF.Equiv.cequiv}{\coqdocdefinition{cequiv}} \coqref{SECF.Imp.::com scope:'<x7B' x 'x7D>'}{\coqdocnotation{<\{}} \coqref{SECF.Imp.:com:com scope:'if' x 'then' x 'else' x 'end'}{\coqdocnotation{if}} \coqref{SECF.Equiv.b:81}{\coqdocvariable{b}} \coqref{SECF.Imp.:com:com scope:'if' x 'then' x 'else' x 'end'}{\coqdocnotation{then}} \coqref{SECF.Equiv.c1:83}{\coqdocvariable{c1}} \coqref{SECF.Imp.:com:com scope:'if' x 'then' x 'else' x 'end'}{\coqdocnotation{else}} \coqref{SECF.Equiv.c2:85}{\coqdocvariable{c2}} \coqref{SECF.Imp.:com:com scope:'if' x 'then' x 'else' x 'end'}{\coqdocnotation{end}} \coqref{SECF.Imp.::com scope:'<x7B' x 'x7D>'}{\coqdocnotation{\}>}}\coqdoceol
\coqdocindent{4.50em}
\coqref{SECF.Imp.::com scope:'<x7B' x 'x7D>'}{\coqdocnotation{<\{}} \coqref{SECF.Imp.:com:com scope:'if' x 'then' x 'else' x 'end'}{\coqdocnotation{if}} \coqref{SECF.Equiv.b':82}{\coqdocvariable{b'}} \coqref{SECF.Imp.:com:com scope:'if' x 'then' x 'else' x 'end'}{\coqdocnotation{then}} \coqref{SECF.Equiv.c1':84}{\coqdocvariable{c1'}} \coqref{SECF.Imp.:com:com scope:'if' x 'then' x 'else' x 'end'}{\coqdocnotation{else}} \coqref{SECF.Equiv.c2':86}{\coqdocvariable{c2'}} \coqref{SECF.Imp.:com:com scope:'if' x 'then' x 'else' x 'end'}{\coqdocnotation{end}} \coqref{SECF.Imp.::com scope:'<x7B' x 'x7D>'}{\coqdocnotation{\}>}}.\coqdoceol
\coqdocnoindent
\coqdockw{Proof}.\coqdoceol
 \coqdocvar{Admitted}.\coqdoceol
\end{coqdoccode}
\ensuremath{\Box} 

 For example, here are two programs and a proof of their
    equivalence using these congruence theorems... 
\begin{coqdoccode}
\coqdocemptyline
\coqdocnoindent
\coqdockw{Example} \coqdef{SECF.Equiv.congruence example}{congruence\_example}{\coqdocdefinition{congruence\_example}}:\coqdoceol
\coqdocindent{1.00em}
\coqref{SECF.Equiv.cequiv}{\coqdocdefinition{cequiv}}\coqdoceol
\coqdocindent{2.00em}
\coqdoceol
\coqdocindent{2.00em}
\coqref{SECF.Imp.::com scope:'<x7B' x 'x7D>'}{\coqdocnotation{<\{}} \coqref{SECF.Imp.X}{\coqdocdefinition{X}} \coqref{SECF.Imp.:com:com scope:x ':=' x}{\coqdocnotation{:=}} 0\coqref{SECF.Imp.:com:com scope:x ';' x}{\coqdocnotation{;}}\coqdoceol
\coqdocindent{3.50em}
\coqref{SECF.Imp.:com:com scope:'if' x 'then' x 'else' x 'end'}{\coqdocnotation{if}} \coqref{SECF.Imp.X}{\coqdocdefinition{X}} \coqref{SECF.Imp.:com::x '=' x}{\coqdocnotation{=}} 0 \coqref{SECF.Imp.:com:com scope:'if' x 'then' x 'else' x 'end'}{\coqdocnotation{then}} \coqref{SECF.Imp.Y}{\coqdocdefinition{Y}} \coqref{SECF.Imp.:com:com scope:x ':=' x}{\coqdocnotation{:=}} 0\coqdoceol
\coqdocindent{3.50em}
\coqref{SECF.Imp.:com:com scope:'if' x 'then' x 'else' x 'end'}{\coqdocnotation{else}} \coqref{SECF.Imp.Y}{\coqdocdefinition{Y}} \coqref{SECF.Imp.:com:com scope:x ':=' x}{\coqdocnotation{:=}} 42 \coqref{SECF.Imp.:com:com scope:'if' x 'then' x 'else' x 'end'}{\coqdocnotation{end}} \coqref{SECF.Imp.::com scope:'<x7B' x 'x7D>'}{\coqdocnotation{\}>}}\coqdoceol
\coqdocindent{2.00em}
\coqdoceol
\coqdocindent{2.00em}
\coqref{SECF.Imp.::com scope:'<x7B' x 'x7D>'}{\coqdocnotation{<\{}} \coqref{SECF.Imp.X}{\coqdocdefinition{X}} \coqref{SECF.Imp.:com:com scope:x ':=' x}{\coqdocnotation{:=}} 0\coqref{SECF.Imp.:com:com scope:x ';' x}{\coqdocnotation{;}}\coqdoceol
\coqdocindent{3.50em}
\coqref{SECF.Imp.:com:com scope:'if' x 'then' x 'else' x 'end'}{\coqdocnotation{if}} \coqref{SECF.Imp.X}{\coqdocdefinition{X}} \coqref{SECF.Imp.:com::x '=' x}{\coqdocnotation{=}} 0 \coqref{SECF.Imp.:com:com scope:'if' x 'then' x 'else' x 'end'}{\coqdocnotation{then}} \coqref{SECF.Imp.Y}{\coqdocdefinition{Y}} \coqref{SECF.Imp.:com:com scope:x ':=' x}{\coqdocnotation{:=}} \coqref{SECF.Imp.X}{\coqdocdefinition{X}} \coqref{SECF.Imp.:com::x '-' x}{\coqdocnotation{-}} \coqref{SECF.Imp.X}{\coqdocdefinition{X}}   \coqdoceol
\coqdocindent{3.50em}
\coqref{SECF.Imp.:com:com scope:'if' x 'then' x 'else' x 'end'}{\coqdocnotation{else}} \coqref{SECF.Imp.Y}{\coqdocdefinition{Y}} \coqref{SECF.Imp.:com:com scope:x ':=' x}{\coqdocnotation{:=}} 42 \coqref{SECF.Imp.:com:com scope:'if' x 'then' x 'else' x 'end'}{\coqdocnotation{end}} \coqref{SECF.Imp.::com scope:'<x7B' x 'x7D>'}{\coqdocnotation{\}>}}.\coqdoceol
\coqdocnoindent
\coqdockw{Proof}.\coqdoceol
\coqdocindent{1.00em}
\coqdoctac{apply} \coqref{SECF.Equiv.CSeq congruence}{\coqdocaxiom{CSeq\_congruence}}.\coqdoceol
\coqdocindent{1.00em}
- \coqdoctac{apply} \coqref{SECF.Equiv.refl cequiv}{\coqdoclemma{refl\_cequiv}}.\coqdoceol
\coqdocindent{1.00em}
- \coqdoctac{apply} \coqref{SECF.Equiv.CIf congruence}{\coqdocaxiom{CIf\_congruence}}.\coqdoceol
\coqdocindent{2.00em}
+ \coqdoctac{apply} \coqref{SECF.Equiv.refl bequiv}{\coqdoclemma{refl\_bequiv}}.\coqdoceol
\coqdocindent{2.00em}
+ \coqdoctac{apply} \coqref{SECF.Equiv.CAsgn congruence}{\coqdoclemma{CAsgn\_congruence}}. \coqdoctac{unfold} \coqref{SECF.Equiv.aequiv}{\coqdocdefinition{aequiv}}. \coqdoctac{simpl}.\coqdoceol
\coqdocindent{3.00em}
\coqdoctac{symmetry}. \coqdoctac{apply} \coqdoclemma{sub\_diag}.\coqdoceol
\coqdocindent{2.00em}
+ \coqdoctac{apply} \coqref{SECF.Equiv.refl cequiv}{\coqdoclemma{refl\_cequiv}}.\coqdoceol
\coqdocnoindent
\coqdockw{Qed}.\coqdoceol
\coqdocemptyline
\end{coqdoccode}
\paragraph{Exercise: 3 stars, advanced (not\_congr)}




    We've shown that the \coqref{SECF.Equiv.Himp.cequiv}{\coqdocdefinition{cequiv}} relation is both an equivalence and
    a congruence on commands.  Can you think of a relation on commands
    that is an equivalence but \textit{not} a congruence?  Write down the
    relation (formally), together with an informal sketch of a proof
    that it is an equivalence and a counterexample showing it is not a
    congruence. 
\begin{coqdoccode}
\coqdocemptyline
\coqdocnoindent
\coqdockw{Definition} \coqdef{SECF.Equiv.manual grade for not congr}{manual\_grade\_for\_not\_congr}{\coqdocdefinition{manual\_grade\_for\_not\_congr}} : \coqexternalref{option}{http://rocq-prover.org/doc/V9.1.1/stdlib//Corelib.Init.Datatypes}{\coqdocinductive{option}} (\coqexternalref{nat}{http://rocq-prover.org/doc/V9.1.1/stdlib//Corelib.Init.Datatypes}{\coqdocinductive{nat}}\coqexternalref{::type scope:x '*' x}{http://rocq-prover.org/doc/V9.1.1/stdlib//Corelib.Init.Datatypes}{\coqdocnotation{\ensuremath{\times}}}\coqdocinductive{string}) := \coqexternalref{None}{http://rocq-prover.org/doc/V9.1.1/stdlib//Corelib.Init.Datatypes}{\coqdocconstructor{None}}.\coqdoceol
\end{coqdoccode}
\ensuremath{\Box} 
\begin{coqdoccode}
\coqdocemptyline
\end{coqdoccode}
\section{Program Transformations}



 A \textit{program transformation} is a function that takes a program as
    input and produces a modified program as output.  Compiler
    optimizations such as constant folding are canonical examples,
    but there are many others. 

 A program transformation is \textit{sound} if it preserves the
    behavior of the original program. 
\begin{coqdoccode}
\coqdocemptyline
\coqdocnoindent
\coqdockw{Definition} \coqdef{SECF.Equiv.atrans sound}{atrans\_sound}{\coqdocdefinition{atrans\_sound}} (\coqdef{SECF.Equiv.atrans:87}{atrans}{\coqdocbinder{atrans}} : \coqref{SECF.Imp.aexp}{\coqdocinductive{aexp}} \coqexternalref{::type scope:x '->' x}{http://rocq-prover.org/doc/V9.1.1/stdlib//Corelib.Init.Logic}{\coqdocnotation{\ensuremath{\rightarrow}}} \coqref{SECF.Imp.aexp}{\coqdocinductive{aexp}}) : \coqdockw{Prop} :=\coqdoceol
\coqdocindent{1.00em}
\coqdockw{\ensuremath{\forall}} (\coqdef{SECF.Equiv.a:88}{a}{\coqdocbinder{a}} : \coqref{SECF.Imp.aexp}{\coqdocinductive{aexp}}),\coqdoceol
\coqdocindent{2.00em}
\coqref{SECF.Equiv.aequiv}{\coqdocdefinition{aequiv}} \coqref{SECF.Equiv.a:88}{\coqdocvariable{a}} (\coqref{SECF.Equiv.atrans:87}{\coqdocvariable{atrans}} \coqref{SECF.Equiv.a:88}{\coqdocvariable{a}}).\coqdoceol
\coqdocemptyline
\coqdocnoindent
\coqdockw{Definition} \coqdef{SECF.Equiv.btrans sound}{btrans\_sound}{\coqdocdefinition{btrans\_sound}} (\coqdef{SECF.Equiv.btrans:89}{btrans}{\coqdocbinder{btrans}} : \coqref{SECF.Imp.bexp}{\coqdocinductive{bexp}} \coqexternalref{::type scope:x '->' x}{http://rocq-prover.org/doc/V9.1.1/stdlib//Corelib.Init.Logic}{\coqdocnotation{\ensuremath{\rightarrow}}} \coqref{SECF.Imp.bexp}{\coqdocinductive{bexp}}) : \coqdockw{Prop} :=\coqdoceol
\coqdocindent{1.00em}
\coqdockw{\ensuremath{\forall}} (\coqdef{SECF.Equiv.b:90}{b}{\coqdocbinder{b}} : \coqref{SECF.Imp.bexp}{\coqdocinductive{bexp}}),\coqdoceol
\coqdocindent{2.00em}
\coqref{SECF.Equiv.bequiv}{\coqdocdefinition{bequiv}} \coqref{SECF.Equiv.b:90}{\coqdocvariable{b}} (\coqref{SECF.Equiv.btrans:89}{\coqdocvariable{btrans}} \coqref{SECF.Equiv.b:90}{\coqdocvariable{b}}).\coqdoceol
\coqdocemptyline
\coqdocnoindent
\coqdockw{Definition} \coqdef{SECF.Equiv.ctrans sound}{ctrans\_sound}{\coqdocdefinition{ctrans\_sound}} (\coqdef{SECF.Equiv.ctrans:91}{ctrans}{\coqdocbinder{ctrans}} : \coqref{SECF.Imp.com}{\coqdocinductive{com}} \coqexternalref{::type scope:x '->' x}{http://rocq-prover.org/doc/V9.1.1/stdlib//Corelib.Init.Logic}{\coqdocnotation{\ensuremath{\rightarrow}}} \coqref{SECF.Imp.com}{\coqdocinductive{com}}) : \coqdockw{Prop} :=\coqdoceol
\coqdocindent{1.00em}
\coqdockw{\ensuremath{\forall}} (\coqdef{SECF.Equiv.c:92}{c}{\coqdocbinder{c}} : \coqref{SECF.Imp.com}{\coqdocinductive{com}}),\coqdoceol
\coqdocindent{2.00em}
\coqref{SECF.Equiv.cequiv}{\coqdocdefinition{cequiv}} \coqref{SECF.Equiv.c:92}{\coqdocvariable{c}} (\coqref{SECF.Equiv.ctrans:91}{\coqdocvariable{ctrans}} \coqref{SECF.Equiv.c:92}{\coqdocvariable{c}}).\coqdoceol
\coqdocemptyline
\end{coqdoccode}
\subsection{The Constant-Folding Transformation}



 An expression is \textit{constant} if it contains no variable references.


    Constant folding is an optimization that finds constant
    expressions and replaces them by their values. 
\begin{coqdoccode}
\coqdocemptyline
\coqdocnoindent
\coqdockw{Fixpoint} \coqdef{SECF.Equiv.fold constants aexp}{fold\_constants\_aexp}{\coqdocdefinition{fold\_constants\_aexp}} (\coqdef{SECF.Equiv.a:93}{a}{\coqdocbinder{a}} : \coqref{SECF.Imp.aexp}{\coqdocinductive{aexp}}) : \coqref{SECF.Imp.aexp}{\coqdocinductive{aexp}} :=\coqdoceol
\coqdocindent{1.00em}
\coqdockw{match} \coqref{SECF.Equiv.a:93}{\coqdocvariable{a}} \coqdockw{with}\coqdoceol
\coqdocindent{1.00em}
\ensuremath{|} \coqref{SECF.Imp.ANum}{\coqdocconstructor{ANum}} \coqdocvar{n}       \ensuremath{\Rightarrow} \coqref{SECF.Imp.ANum}{\coqdocconstructor{ANum}} \coqdocvar{n}\coqdoceol
\coqdocindent{1.00em}
\ensuremath{|} \coqref{SECF.Imp.AId}{\coqdocconstructor{AId}} \coqdocvar{x}        \ensuremath{\Rightarrow} \coqref{SECF.Imp.AId}{\coqdocconstructor{AId}} \coqdocvar{x}\coqdoceol
\coqdocindent{1.00em}
\ensuremath{|} \coqref{SECF.Imp.::com scope:'<x7B' x 'x7D>'}{\coqdocnotation{<\{}} \coqdocvar{a1} \coqref{SECF.Imp.:com::x '+' x}{\coqdocnotation{+}} \coqdocvar{a2} \coqref{SECF.Imp.::com scope:'<x7B' x 'x7D>'}{\coqdocnotation{\}>}}  \ensuremath{\Rightarrow}\coqdoceol
\coqdocindent{2.00em}
\coqdockw{match} \coqexternalref{::core scope:'(' x ',' x ',' '..' ',' x ')'}{http://rocq-prover.org/doc/V9.1.1/stdlib//Corelib.Init.Datatypes}{\coqdocnotation{(}}\coqref{SECF.Equiv.fold constants aexp:94}{\coqdocdefinition{fold\_constants\_aexp}} \coqdocvar{a1}\coqexternalref{::core scope:'(' x ',' x ',' '..' ',' x ')'}{http://rocq-prover.org/doc/V9.1.1/stdlib//Corelib.Init.Datatypes}{\coqdocnotation{,}}\coqdoceol
\coqdocindent{5.50em}
\coqref{SECF.Equiv.fold constants aexp:94}{\coqdocdefinition{fold\_constants\_aexp}} \coqdocvar{a2}\coqexternalref{::core scope:'(' x ',' x ',' '..' ',' x ')'}{http://rocq-prover.org/doc/V9.1.1/stdlib//Corelib.Init.Datatypes}{\coqdocnotation{)}}\coqdoceol
\coqdocindent{2.00em}
\coqdockw{with}\coqdoceol
\coqdocindent{2.00em}
\ensuremath{|} \coqexternalref{::core scope:'(' x ',' x ',' '..' ',' x ')'}{http://rocq-prover.org/doc/V9.1.1/stdlib//Corelib.Init.Datatypes}{\coqdocnotation{(}}\coqref{SECF.Imp.ANum}{\coqdocconstructor{ANum}} \coqdocvar{n1}\coqexternalref{::core scope:'(' x ',' x ',' '..' ',' x ')'}{http://rocq-prover.org/doc/V9.1.1/stdlib//Corelib.Init.Datatypes}{\coqdocnotation{,}} \coqref{SECF.Imp.ANum}{\coqdocconstructor{ANum}} \coqdocvar{n2}\coqexternalref{::core scope:'(' x ',' x ',' '..' ',' x ')'}{http://rocq-prover.org/doc/V9.1.1/stdlib//Corelib.Init.Datatypes}{\coqdocnotation{)}} \ensuremath{\Rightarrow} \coqref{SECF.Imp.ANum}{\coqdocconstructor{ANum}} (\coqdocvar{n1} \coqexternalref{::nat scope:x '+' x}{http://rocq-prover.org/doc/V9.1.1/stdlib//Corelib.Init.Nat}{\coqdocnotation{+}} \coqdocvar{n2})\coqdoceol
\coqdocindent{2.00em}
\ensuremath{|} \coqexternalref{::core scope:'(' x ',' x ',' '..' ',' x ')'}{http://rocq-prover.org/doc/V9.1.1/stdlib//Corelib.Init.Datatypes}{\coqdocnotation{(}}\coqdocvar{a1'}\coqexternalref{::core scope:'(' x ',' x ',' '..' ',' x ')'}{http://rocq-prover.org/doc/V9.1.1/stdlib//Corelib.Init.Datatypes}{\coqdocnotation{,}} \coqdocvar{a2'}\coqexternalref{::core scope:'(' x ',' x ',' '..' ',' x ')'}{http://rocq-prover.org/doc/V9.1.1/stdlib//Corelib.Init.Datatypes}{\coqdocnotation{)}} \ensuremath{\Rightarrow} \coqref{SECF.Imp.::com scope:'<x7B' x 'x7D>'}{\coqdocnotation{<\{}} \coqdocvar{a1'} \coqref{SECF.Imp.:com::x '+' x}{\coqdocnotation{+}} \coqdocvar{a2'} \coqref{SECF.Imp.::com scope:'<x7B' x 'x7D>'}{\coqdocnotation{\}>}}\coqdoceol
\coqdocindent{2.00em}
\coqdockw{end}\coqdoceol
\coqdocindent{1.00em}
\ensuremath{|} \coqref{SECF.Imp.::com scope:'<x7B' x 'x7D>'}{\coqdocnotation{<\{}} \coqdocvar{a1} \coqref{SECF.Imp.:com::x '-' x}{\coqdocnotation{-}} \coqdocvar{a2} \coqref{SECF.Imp.::com scope:'<x7B' x 'x7D>'}{\coqdocnotation{\}>}} \ensuremath{\Rightarrow}\coqdoceol
\coqdocindent{2.00em}
\coqdockw{match} \coqexternalref{::core scope:'(' x ',' x ',' '..' ',' x ')'}{http://rocq-prover.org/doc/V9.1.1/stdlib//Corelib.Init.Datatypes}{\coqdocnotation{(}}\coqref{SECF.Equiv.fold constants aexp:94}{\coqdocdefinition{fold\_constants\_aexp}} \coqdocvar{a1}\coqexternalref{::core scope:'(' x ',' x ',' '..' ',' x ')'}{http://rocq-prover.org/doc/V9.1.1/stdlib//Corelib.Init.Datatypes}{\coqdocnotation{,}}\coqdoceol
\coqdocindent{5.50em}
\coqref{SECF.Equiv.fold constants aexp:94}{\coqdocdefinition{fold\_constants\_aexp}} \coqdocvar{a2}\coqexternalref{::core scope:'(' x ',' x ',' '..' ',' x ')'}{http://rocq-prover.org/doc/V9.1.1/stdlib//Corelib.Init.Datatypes}{\coqdocnotation{)}}\coqdoceol
\coqdocindent{2.00em}
\coqdockw{with}\coqdoceol
\coqdocindent{2.00em}
\ensuremath{|} \coqexternalref{::core scope:'(' x ',' x ',' '..' ',' x ')'}{http://rocq-prover.org/doc/V9.1.1/stdlib//Corelib.Init.Datatypes}{\coqdocnotation{(}}\coqref{SECF.Imp.ANum}{\coqdocconstructor{ANum}} \coqdocvar{n1}\coqexternalref{::core scope:'(' x ',' x ',' '..' ',' x ')'}{http://rocq-prover.org/doc/V9.1.1/stdlib//Corelib.Init.Datatypes}{\coqdocnotation{,}} \coqref{SECF.Imp.ANum}{\coqdocconstructor{ANum}} \coqdocvar{n2}\coqexternalref{::core scope:'(' x ',' x ',' '..' ',' x ')'}{http://rocq-prover.org/doc/V9.1.1/stdlib//Corelib.Init.Datatypes}{\coqdocnotation{)}} \ensuremath{\Rightarrow} \coqref{SECF.Imp.ANum}{\coqdocconstructor{ANum}} (\coqdocvar{n1} \coqexternalref{::nat scope:x '-' x}{http://rocq-prover.org/doc/V9.1.1/stdlib//Corelib.Init.Nat}{\coqdocnotation{-}} \coqdocvar{n2})\coqdoceol
\coqdocindent{2.00em}
\ensuremath{|} \coqexternalref{::core scope:'(' x ',' x ',' '..' ',' x ')'}{http://rocq-prover.org/doc/V9.1.1/stdlib//Corelib.Init.Datatypes}{\coqdocnotation{(}}\coqdocvar{a1'}\coqexternalref{::core scope:'(' x ',' x ',' '..' ',' x ')'}{http://rocq-prover.org/doc/V9.1.1/stdlib//Corelib.Init.Datatypes}{\coqdocnotation{,}} \coqdocvar{a2'}\coqexternalref{::core scope:'(' x ',' x ',' '..' ',' x ')'}{http://rocq-prover.org/doc/V9.1.1/stdlib//Corelib.Init.Datatypes}{\coqdocnotation{)}} \ensuremath{\Rightarrow} \coqref{SECF.Imp.::com scope:'<x7B' x 'x7D>'}{\coqdocnotation{<\{}} \coqdocvar{a1'} \coqref{SECF.Imp.:com::x '-' x}{\coqdocnotation{-}} \coqdocvar{a2'} \coqref{SECF.Imp.::com scope:'<x7B' x 'x7D>'}{\coqdocnotation{\}>}}\coqdoceol
\coqdocindent{2.00em}
\coqdockw{end}\coqdoceol
\coqdocindent{1.00em}
\ensuremath{|} \coqref{SECF.Imp.::com scope:'<x7B' x 'x7D>'}{\coqdocnotation{<\{}} \coqdocvar{a1} \coqref{SECF.Imp.:com::x '*' x}{\coqdocnotation{\ensuremath{\times}}} \coqdocvar{a2} \coqref{SECF.Imp.::com scope:'<x7B' x 'x7D>'}{\coqdocnotation{\}>}}  \ensuremath{\Rightarrow}\coqdoceol
\coqdocindent{2.00em}
\coqdockw{match} \coqexternalref{::core scope:'(' x ',' x ',' '..' ',' x ')'}{http://rocq-prover.org/doc/V9.1.1/stdlib//Corelib.Init.Datatypes}{\coqdocnotation{(}}\coqref{SECF.Equiv.fold constants aexp:94}{\coqdocdefinition{fold\_constants\_aexp}} \coqdocvar{a1}\coqexternalref{::core scope:'(' x ',' x ',' '..' ',' x ')'}{http://rocq-prover.org/doc/V9.1.1/stdlib//Corelib.Init.Datatypes}{\coqdocnotation{,}}\coqdoceol
\coqdocindent{5.50em}
\coqref{SECF.Equiv.fold constants aexp:94}{\coqdocdefinition{fold\_constants\_aexp}} \coqdocvar{a2}\coqexternalref{::core scope:'(' x ',' x ',' '..' ',' x ')'}{http://rocq-prover.org/doc/V9.1.1/stdlib//Corelib.Init.Datatypes}{\coqdocnotation{)}}\coqdoceol
\coqdocindent{2.00em}
\coqdockw{with}\coqdoceol
\coqdocindent{2.00em}
\ensuremath{|} \coqexternalref{::core scope:'(' x ',' x ',' '..' ',' x ')'}{http://rocq-prover.org/doc/V9.1.1/stdlib//Corelib.Init.Datatypes}{\coqdocnotation{(}}\coqref{SECF.Imp.ANum}{\coqdocconstructor{ANum}} \coqdocvar{n1}\coqexternalref{::core scope:'(' x ',' x ',' '..' ',' x ')'}{http://rocq-prover.org/doc/V9.1.1/stdlib//Corelib.Init.Datatypes}{\coqdocnotation{,}} \coqref{SECF.Imp.ANum}{\coqdocconstructor{ANum}} \coqdocvar{n2}\coqexternalref{::core scope:'(' x ',' x ',' '..' ',' x ')'}{http://rocq-prover.org/doc/V9.1.1/stdlib//Corelib.Init.Datatypes}{\coqdocnotation{)}} \ensuremath{\Rightarrow} \coqref{SECF.Imp.ANum}{\coqdocconstructor{ANum}} (\coqdocvar{n1} \coqexternalref{::nat scope:x '*' x}{http://rocq-prover.org/doc/V9.1.1/stdlib//Corelib.Init.Nat}{\coqdocnotation{\ensuremath{\times}}} \coqdocvar{n2})\coqdoceol
\coqdocindent{2.00em}
\ensuremath{|} \coqexternalref{::core scope:'(' x ',' x ',' '..' ',' x ')'}{http://rocq-prover.org/doc/V9.1.1/stdlib//Corelib.Init.Datatypes}{\coqdocnotation{(}}\coqdocvar{a1'}\coqexternalref{::core scope:'(' x ',' x ',' '..' ',' x ')'}{http://rocq-prover.org/doc/V9.1.1/stdlib//Corelib.Init.Datatypes}{\coqdocnotation{,}} \coqdocvar{a2'}\coqexternalref{::core scope:'(' x ',' x ',' '..' ',' x ')'}{http://rocq-prover.org/doc/V9.1.1/stdlib//Corelib.Init.Datatypes}{\coqdocnotation{)}} \ensuremath{\Rightarrow} \coqref{SECF.Imp.::com scope:'<x7B' x 'x7D>'}{\coqdocnotation{<\{}} \coqdocvar{a1'} \coqref{SECF.Imp.:com::x '*' x}{\coqdocnotation{\ensuremath{\times}}} \coqdocvar{a2'} \coqref{SECF.Imp.::com scope:'<x7B' x 'x7D>'}{\coqdocnotation{\}>}}\coqdoceol
\coqdocindent{2.00em}
\coqdockw{end}\coqdoceol
\coqdocindent{1.00em}
\coqdockw{end}.\coqdoceol
\coqdocemptyline
\coqdocnoindent
\coqdockw{Example} \coqdef{SECF.Equiv.fold aexp ex1}{fold\_aexp\_ex1}{\coqdocdefinition{fold\_aexp\_ex1}} :\coqdoceol
\coqdocindent{2.00em}
\coqref{SECF.Equiv.fold constants aexp}{\coqdocdefinition{fold\_constants\_aexp}} \coqref{SECF.Imp.::com scope:'<x7B' x 'x7D>'}{\coqdocnotation{<\{}} \coqref{SECF.Imp.:com::'(' x ')'}{\coqdocnotation{(}}1 \coqref{SECF.Imp.:com::x '+' x}{\coqdocnotation{+}} 2\coqref{SECF.Imp.:com::'(' x ')'}{\coqdocnotation{)}} \coqref{SECF.Imp.:com::x '*' x}{\coqdocnotation{\ensuremath{\times}}} \coqref{SECF.Imp.X}{\coqdocdefinition{X}} \coqref{SECF.Imp.::com scope:'<x7B' x 'x7D>'}{\coqdocnotation{\}>}}\coqdoceol
\coqdocindent{1.00em}
\coqexternalref{::type scope:x '=' x}{http://rocq-prover.org/doc/V9.1.1/stdlib//Corelib.Init.Logic}{\coqdocnotation{=}} \coqref{SECF.Imp.::com scope:'<x7B' x 'x7D>'}{\coqdocnotation{<\{}} 3 \coqref{SECF.Imp.:com::x '*' x}{\coqdocnotation{\ensuremath{\times}}} \coqref{SECF.Imp.X}{\coqdocdefinition{X}} \coqref{SECF.Imp.::com scope:'<x7B' x 'x7D>'}{\coqdocnotation{\}>}}.\coqdoceol
\coqdocnoindent
\coqdockw{Proof}. \coqdoctac{reflexivity}. \coqdockw{Qed}.\coqdoceol
\coqdocemptyline
\end{coqdoccode}
Note that this version of constant folding doesn't do other
    ``obvious'' things like eliminating trivial additions (e.g.,
    rewriting 0 + \coqref{SECF.SpecCT.X}{\coqdocdefinition{X}} to just \coqref{SECF.SpecCT.X}{\coqdocdefinition{X}}).: we are focusing on a single
    optimization for the sake of simplicity.


    It is not hard to incorporate other ways of simplifying
    expressions -- the definitions and proofs just get longer.  We'll
    consider some in the exercises. 
\begin{coqdoccode}
\coqdocemptyline
\coqdocnoindent
\coqdockw{Example} \coqdef{SECF.Equiv.fold aexp ex2}{fold\_aexp\_ex2}{\coqdocdefinition{fold\_aexp\_ex2}} :\coqdoceol
\coqdocindent{1.00em}
\coqref{SECF.Equiv.fold constants aexp}{\coqdocdefinition{fold\_constants\_aexp}} \coqref{SECF.Imp.::com scope:'<x7B' x 'x7D>'}{\coqdocnotation{<\{}} \coqref{SECF.Imp.X}{\coqdocdefinition{X}} \coqref{SECF.Imp.:com::x '-' x}{\coqdocnotation{-}} \coqref{SECF.Imp.:com::'(' x ')'}{\coqdocnotation{((}}0 \coqref{SECF.Imp.:com::x '*' x}{\coqdocnotation{\ensuremath{\times}}} 6\coqref{SECF.Imp.:com::'(' x ')'}{\coqdocnotation{)}} \coqref{SECF.Imp.:com::x '+' x}{\coqdocnotation{+}} \coqref{SECF.Imp.Y}{\coqdocdefinition{Y}}\coqref{SECF.Imp.:com::'(' x ')'}{\coqdocnotation{)}} \coqref{SECF.Imp.::com scope:'<x7B' x 'x7D>'}{\coqdocnotation{\}>}} \coqexternalref{::type scope:x '=' x}{http://rocq-prover.org/doc/V9.1.1/stdlib//Corelib.Init.Logic}{\coqdocnotation{=}} \coqref{SECF.Imp.::com scope:'<x7B' x 'x7D>'}{\coqdocnotation{<\{}} \coqref{SECF.Imp.X}{\coqdocdefinition{X}} \coqref{SECF.Imp.:com::x '-' x}{\coqdocnotation{-}} \coqref{SECF.Imp.:com::'(' x ')'}{\coqdocnotation{(}}0 \coqref{SECF.Imp.:com::x '+' x}{\coqdocnotation{+}} \coqref{SECF.Imp.Y}{\coqdocdefinition{Y}}\coqref{SECF.Imp.:com::'(' x ')'}{\coqdocnotation{)}} \coqref{SECF.Imp.::com scope:'<x7B' x 'x7D>'}{\coqdocnotation{\}>}}.\coqdoceol
\coqdocnoindent
\coqdockw{Proof}. \coqdoctac{reflexivity}. \coqdockw{Qed}.\coqdoceol
\coqdocemptyline
\end{coqdoccode}
Not only can we lift \coqref{SECF.Equiv.fold constants aexp}{\coqdocdefinition{fold\_constants\_aexp}} to \coqref{SECF.Imp.bexp}{\coqdocinductive{bexp}}s (in the
    \coqref{SECF.SpecCT.BEq}{\coqdocdefinition{BEq}}, \coqref{SECF.SpecCT.BNeq}{\coqdocdefinition{BNeq}}, and \coqref{SECF.SpecCT.BLe}{\coqdocdefinition{BLe}} cases); we can also look for constant
    \textit{boolean} expressions and evaluate them in-place as well. 
\begin{coqdoccode}
\coqdocemptyline
\coqdocnoindent
\coqdockw{Fixpoint} \coqdef{SECF.Equiv.fold constants bexp}{fold\_constants\_bexp}{\coqdocdefinition{fold\_constants\_bexp}} (\coqdef{SECF.Equiv.b:96}{b}{\coqdocbinder{b}} : \coqref{SECF.Imp.bexp}{\coqdocinductive{bexp}}) : \coqref{SECF.Imp.bexp}{\coqdocinductive{bexp}} :=\coqdoceol
\coqdocindent{1.00em}
\coqdockw{match} \coqref{SECF.Equiv.b:96}{\coqdocvariable{b}} \coqdockw{with}\coqdoceol
\coqdocindent{1.00em}
\ensuremath{|} \coqref{SECF.Imp.::com scope:'<x7B' x 'x7D>'}{\coqdocnotation{<\{}}\coqref{SECF.Imp.:com::'true'}{\coqdocnotation{true}}\coqref{SECF.Imp.::com scope:'<x7B' x 'x7D>'}{\coqdocnotation{\}>}}        \ensuremath{\Rightarrow} \coqref{SECF.Imp.::com scope:'<x7B' x 'x7D>'}{\coqdocnotation{<\{}}\coqref{SECF.Imp.:com::'true'}{\coqdocnotation{true}}\coqref{SECF.Imp.::com scope:'<x7B' x 'x7D>'}{\coqdocnotation{\}>}}\coqdoceol
\coqdocindent{1.00em}
\ensuremath{|} \coqref{SECF.Imp.::com scope:'<x7B' x 'x7D>'}{\coqdocnotation{<\{}}\coqref{SECF.Imp.:com::'false'}{\coqdocnotation{false}}\coqref{SECF.Imp.::com scope:'<x7B' x 'x7D>'}{\coqdocnotation{\}>}}       \ensuremath{\Rightarrow} \coqref{SECF.Imp.::com scope:'<x7B' x 'x7D>'}{\coqdocnotation{<\{}}\coqref{SECF.Imp.:com::'false'}{\coqdocnotation{false}}\coqref{SECF.Imp.::com scope:'<x7B' x 'x7D>'}{\coqdocnotation{\}>}}\coqdoceol
\coqdocindent{1.00em}
\ensuremath{|} \coqref{SECF.Imp.::com scope:'<x7B' x 'x7D>'}{\coqdocnotation{<\{}} \coqdocvar{a1} \coqref{SECF.Imp.:com::x '=' x}{\coqdocnotation{=}} \coqdocvar{a2} \coqref{SECF.Imp.::com scope:'<x7B' x 'x7D>'}{\coqdocnotation{\}>}}  \ensuremath{\Rightarrow}\coqdoceol
\coqdocindent{3.00em}
\coqdockw{match} \coqexternalref{::core scope:'(' x ',' x ',' '..' ',' x ')'}{http://rocq-prover.org/doc/V9.1.1/stdlib//Corelib.Init.Datatypes}{\coqdocnotation{(}}\coqref{SECF.Equiv.fold constants aexp}{\coqdocdefinition{fold\_constants\_aexp}} \coqdocvar{a1}\coqexternalref{::core scope:'(' x ',' x ',' '..' ',' x ')'}{http://rocq-prover.org/doc/V9.1.1/stdlib//Corelib.Init.Datatypes}{\coqdocnotation{,}}\coqdoceol
\coqdocindent{6.50em}
\coqref{SECF.Equiv.fold constants aexp}{\coqdocdefinition{fold\_constants\_aexp}} \coqdocvar{a2}\coqexternalref{::core scope:'(' x ',' x ',' '..' ',' x ')'}{http://rocq-prover.org/doc/V9.1.1/stdlib//Corelib.Init.Datatypes}{\coqdocnotation{)}} \coqdockw{with}\coqdoceol
\coqdocindent{3.00em}
\ensuremath{|} \coqexternalref{::core scope:'(' x ',' x ',' '..' ',' x ')'}{http://rocq-prover.org/doc/V9.1.1/stdlib//Corelib.Init.Datatypes}{\coqdocnotation{(}}\coqref{SECF.Imp.ANum}{\coqdocconstructor{ANum}} \coqdocvar{n1}\coqexternalref{::core scope:'(' x ',' x ',' '..' ',' x ')'}{http://rocq-prover.org/doc/V9.1.1/stdlib//Corelib.Init.Datatypes}{\coqdocnotation{,}} \coqref{SECF.Imp.ANum}{\coqdocconstructor{ANum}} \coqdocvar{n2}\coqexternalref{::core scope:'(' x ',' x ',' '..' ',' x ')'}{http://rocq-prover.org/doc/V9.1.1/stdlib//Corelib.Init.Datatypes}{\coqdocnotation{)}} \ensuremath{\Rightarrow}\coqdoceol
\coqdocindent{5.00em}
\coqdockw{if} \coqdocvar{n1} \coqexternalref{::nat scope:x '=?' x}{http://rocq-prover.org/doc/V9.1.1/stdlib//Corelib.Init.Nat}{\coqdocnotation{=?}} \coqdocvar{n2} \coqdockw{then} \coqref{SECF.Imp.::com scope:'<x7B' x 'x7D>'}{\coqdocnotation{<\{}}\coqref{SECF.Imp.:com::'true'}{\coqdocnotation{true}}\coqref{SECF.Imp.::com scope:'<x7B' x 'x7D>'}{\coqdocnotation{\}>}} \coqdockw{else} \coqref{SECF.Imp.::com scope:'<x7B' x 'x7D>'}{\coqdocnotation{<\{}}\coqref{SECF.Imp.:com::'false'}{\coqdocnotation{false}}\coqref{SECF.Imp.::com scope:'<x7B' x 'x7D>'}{\coqdocnotation{\}>}}\coqdoceol
\coqdocindent{3.00em}
\ensuremath{|} \coqexternalref{::core scope:'(' x ',' x ',' '..' ',' x ')'}{http://rocq-prover.org/doc/V9.1.1/stdlib//Corelib.Init.Datatypes}{\coqdocnotation{(}}\coqdocvar{a1'}\coqexternalref{::core scope:'(' x ',' x ',' '..' ',' x ')'}{http://rocq-prover.org/doc/V9.1.1/stdlib//Corelib.Init.Datatypes}{\coqdocnotation{,}} \coqdocvar{a2'}\coqexternalref{::core scope:'(' x ',' x ',' '..' ',' x ')'}{http://rocq-prover.org/doc/V9.1.1/stdlib//Corelib.Init.Datatypes}{\coqdocnotation{)}} \ensuremath{\Rightarrow}\coqdoceol
\coqdocindent{5.00em}
\coqref{SECF.Imp.::com scope:'<x7B' x 'x7D>'}{\coqdocnotation{<\{}} \coqdocvar{a1'} \coqref{SECF.Imp.:com::x '=' x}{\coqdocnotation{=}} \coqdocvar{a2'} \coqref{SECF.Imp.::com scope:'<x7B' x 'x7D>'}{\coqdocnotation{\}>}}\coqdoceol
\coqdocindent{3.00em}
\coqdockw{end}\coqdoceol
\coqdocindent{1.00em}
\ensuremath{|} \coqref{SECF.Imp.::com scope:'<x7B' x 'x7D>'}{\coqdocnotation{<\{}} \coqdocvar{a1} \coqref{SECF.Imp.:com::x '<>' x}{\coqdocnotation{\ensuremath{\not=}}} \coqdocvar{a2} \coqref{SECF.Imp.::com scope:'<x7B' x 'x7D>'}{\coqdocnotation{\}>}}  \ensuremath{\Rightarrow}\coqdoceol
\coqdocindent{3.00em}
\coqdockw{match} \coqexternalref{::core scope:'(' x ',' x ',' '..' ',' x ')'}{http://rocq-prover.org/doc/V9.1.1/stdlib//Corelib.Init.Datatypes}{\coqdocnotation{(}}\coqref{SECF.Equiv.fold constants aexp}{\coqdocdefinition{fold\_constants\_aexp}} \coqdocvar{a1}\coqexternalref{::core scope:'(' x ',' x ',' '..' ',' x ')'}{http://rocq-prover.org/doc/V9.1.1/stdlib//Corelib.Init.Datatypes}{\coqdocnotation{,}}\coqdoceol
\coqdocindent{6.50em}
\coqref{SECF.Equiv.fold constants aexp}{\coqdocdefinition{fold\_constants\_aexp}} \coqdocvar{a2}\coqexternalref{::core scope:'(' x ',' x ',' '..' ',' x ')'}{http://rocq-prover.org/doc/V9.1.1/stdlib//Corelib.Init.Datatypes}{\coqdocnotation{)}} \coqdockw{with}\coqdoceol
\coqdocindent{3.00em}
\ensuremath{|} \coqexternalref{::core scope:'(' x ',' x ',' '..' ',' x ')'}{http://rocq-prover.org/doc/V9.1.1/stdlib//Corelib.Init.Datatypes}{\coqdocnotation{(}}\coqref{SECF.Imp.ANum}{\coqdocconstructor{ANum}} \coqdocvar{n1}\coqexternalref{::core scope:'(' x ',' x ',' '..' ',' x ')'}{http://rocq-prover.org/doc/V9.1.1/stdlib//Corelib.Init.Datatypes}{\coqdocnotation{,}} \coqref{SECF.Imp.ANum}{\coqdocconstructor{ANum}} \coqdocvar{n2}\coqexternalref{::core scope:'(' x ',' x ',' '..' ',' x ')'}{http://rocq-prover.org/doc/V9.1.1/stdlib//Corelib.Init.Datatypes}{\coqdocnotation{)}} \ensuremath{\Rightarrow}\coqdoceol
\coqdocindent{5.00em}
\coqdockw{if} \coqexternalref{negb}{http://rocq-prover.org/doc/V9.1.1/stdlib//Corelib.Init.Datatypes}{\coqdocdefinition{negb}} (\coqdocvar{n1} \coqexternalref{::nat scope:x '=?' x}{http://rocq-prover.org/doc/V9.1.1/stdlib//Corelib.Init.Nat}{\coqdocnotation{=?}} \coqdocvar{n2}) \coqdockw{then} \coqref{SECF.Imp.::com scope:'<x7B' x 'x7D>'}{\coqdocnotation{<\{}}\coqref{SECF.Imp.:com::'true'}{\coqdocnotation{true}}\coqref{SECF.Imp.::com scope:'<x7B' x 'x7D>'}{\coqdocnotation{\}>}} \coqdockw{else} \coqref{SECF.Imp.::com scope:'<x7B' x 'x7D>'}{\coqdocnotation{<\{}}\coqref{SECF.Imp.:com::'false'}{\coqdocnotation{false}}\coqref{SECF.Imp.::com scope:'<x7B' x 'x7D>'}{\coqdocnotation{\}>}}\coqdoceol
\coqdocindent{3.00em}
\ensuremath{|} \coqexternalref{::core scope:'(' x ',' x ',' '..' ',' x ')'}{http://rocq-prover.org/doc/V9.1.1/stdlib//Corelib.Init.Datatypes}{\coqdocnotation{(}}\coqdocvar{a1'}\coqexternalref{::core scope:'(' x ',' x ',' '..' ',' x ')'}{http://rocq-prover.org/doc/V9.1.1/stdlib//Corelib.Init.Datatypes}{\coqdocnotation{,}} \coqdocvar{a2'}\coqexternalref{::core scope:'(' x ',' x ',' '..' ',' x ')'}{http://rocq-prover.org/doc/V9.1.1/stdlib//Corelib.Init.Datatypes}{\coqdocnotation{)}} \ensuremath{\Rightarrow}\coqdoceol
\coqdocindent{5.00em}
\coqref{SECF.Imp.::com scope:'<x7B' x 'x7D>'}{\coqdocnotation{<\{}} \coqdocvar{a1'} \coqref{SECF.Imp.:com::x '<>' x}{\coqdocnotation{\ensuremath{\not=}}} \coqdocvar{a2'} \coqref{SECF.Imp.::com scope:'<x7B' x 'x7D>'}{\coqdocnotation{\}>}}\coqdoceol
\coqdocindent{3.00em}
\coqdockw{end}\coqdoceol
\coqdocindent{1.00em}
\ensuremath{|} \coqref{SECF.Imp.::com scope:'<x7B' x 'x7D>'}{\coqdocnotation{<\{}} \coqdocvar{a1} \coqref{SECF.Imp.:com::x '<=' x}{\coqdocnotation{\ensuremath{\le}}} \coqdocvar{a2} \coqref{SECF.Imp.::com scope:'<x7B' x 'x7D>'}{\coqdocnotation{\}>}}  \ensuremath{\Rightarrow}\coqdoceol
\coqdocindent{3.00em}
\coqdockw{match} \coqexternalref{::core scope:'(' x ',' x ',' '..' ',' x ')'}{http://rocq-prover.org/doc/V9.1.1/stdlib//Corelib.Init.Datatypes}{\coqdocnotation{(}}\coqref{SECF.Equiv.fold constants aexp}{\coqdocdefinition{fold\_constants\_aexp}} \coqdocvar{a1}\coqexternalref{::core scope:'(' x ',' x ',' '..' ',' x ')'}{http://rocq-prover.org/doc/V9.1.1/stdlib//Corelib.Init.Datatypes}{\coqdocnotation{,}}\coqdoceol
\coqdocindent{6.50em}
\coqref{SECF.Equiv.fold constants aexp}{\coqdocdefinition{fold\_constants\_aexp}} \coqdocvar{a2}\coqexternalref{::core scope:'(' x ',' x ',' '..' ',' x ')'}{http://rocq-prover.org/doc/V9.1.1/stdlib//Corelib.Init.Datatypes}{\coqdocnotation{)}} \coqdockw{with}\coqdoceol
\coqdocindent{3.00em}
\ensuremath{|} \coqexternalref{::core scope:'(' x ',' x ',' '..' ',' x ')'}{http://rocq-prover.org/doc/V9.1.1/stdlib//Corelib.Init.Datatypes}{\coqdocnotation{(}}\coqref{SECF.Imp.ANum}{\coqdocconstructor{ANum}} \coqdocvar{n1}\coqexternalref{::core scope:'(' x ',' x ',' '..' ',' x ')'}{http://rocq-prover.org/doc/V9.1.1/stdlib//Corelib.Init.Datatypes}{\coqdocnotation{,}} \coqref{SECF.Imp.ANum}{\coqdocconstructor{ANum}} \coqdocvar{n2}\coqexternalref{::core scope:'(' x ',' x ',' '..' ',' x ')'}{http://rocq-prover.org/doc/V9.1.1/stdlib//Corelib.Init.Datatypes}{\coqdocnotation{)}} \ensuremath{\Rightarrow}\coqdoceol
\coqdocindent{5.00em}
\coqdockw{if} \coqdocvar{n1} \coqexternalref{::nat scope:x '<=?' x}{http://rocq-prover.org/doc/V9.1.1/stdlib//Corelib.Init.Nat}{\coqdocnotation{<=?}} \coqdocvar{n2} \coqdockw{then} \coqref{SECF.Imp.::com scope:'<x7B' x 'x7D>'}{\coqdocnotation{<\{}}\coqref{SECF.Imp.:com::'true'}{\coqdocnotation{true}}\coqref{SECF.Imp.::com scope:'<x7B' x 'x7D>'}{\coqdocnotation{\}>}} \coqdockw{else} \coqref{SECF.Imp.::com scope:'<x7B' x 'x7D>'}{\coqdocnotation{<\{}}\coqref{SECF.Imp.:com::'false'}{\coqdocnotation{false}}\coqref{SECF.Imp.::com scope:'<x7B' x 'x7D>'}{\coqdocnotation{\}>}}\coqdoceol
\coqdocindent{3.00em}
\ensuremath{|} \coqexternalref{::core scope:'(' x ',' x ',' '..' ',' x ')'}{http://rocq-prover.org/doc/V9.1.1/stdlib//Corelib.Init.Datatypes}{\coqdocnotation{(}}\coqdocvar{a1'}\coqexternalref{::core scope:'(' x ',' x ',' '..' ',' x ')'}{http://rocq-prover.org/doc/V9.1.1/stdlib//Corelib.Init.Datatypes}{\coqdocnotation{,}} \coqdocvar{a2'}\coqexternalref{::core scope:'(' x ',' x ',' '..' ',' x ')'}{http://rocq-prover.org/doc/V9.1.1/stdlib//Corelib.Init.Datatypes}{\coqdocnotation{)}} \ensuremath{\Rightarrow}\coqdoceol
\coqdocindent{5.00em}
\coqref{SECF.Imp.::com scope:'<x7B' x 'x7D>'}{\coqdocnotation{<\{}} \coqdocvar{a1'} \coqref{SECF.Imp.:com::x '<=' x}{\coqdocnotation{\ensuremath{\le}}} \coqdocvar{a2'} \coqref{SECF.Imp.::com scope:'<x7B' x 'x7D>'}{\coqdocnotation{\}>}}\coqdoceol
\coqdocindent{3.00em}
\coqdockw{end}\coqdoceol
\coqdocindent{1.00em}
\ensuremath{|} \coqref{SECF.Imp.::com scope:'<x7B' x 'x7D>'}{\coqdocnotation{<\{}} \coqdocvar{a1} \coqref{SECF.Imp.:com::x '>' x}{\coqdocnotation{>}} \coqdocvar{a2} \coqref{SECF.Imp.::com scope:'<x7B' x 'x7D>'}{\coqdocnotation{\}>}}  \ensuremath{\Rightarrow}\coqdoceol
\coqdocindent{3.00em}
\coqdockw{match} \coqexternalref{::core scope:'(' x ',' x ',' '..' ',' x ')'}{http://rocq-prover.org/doc/V9.1.1/stdlib//Corelib.Init.Datatypes}{\coqdocnotation{(}}\coqref{SECF.Equiv.fold constants aexp}{\coqdocdefinition{fold\_constants\_aexp}} \coqdocvar{a1}\coqexternalref{::core scope:'(' x ',' x ',' '..' ',' x ')'}{http://rocq-prover.org/doc/V9.1.1/stdlib//Corelib.Init.Datatypes}{\coqdocnotation{,}}\coqdoceol
\coqdocindent{6.50em}
\coqref{SECF.Equiv.fold constants aexp}{\coqdocdefinition{fold\_constants\_aexp}} \coqdocvar{a2}\coqexternalref{::core scope:'(' x ',' x ',' '..' ',' x ')'}{http://rocq-prover.org/doc/V9.1.1/stdlib//Corelib.Init.Datatypes}{\coqdocnotation{)}} \coqdockw{with}\coqdoceol
\coqdocindent{3.00em}
\ensuremath{|} \coqexternalref{::core scope:'(' x ',' x ',' '..' ',' x ')'}{http://rocq-prover.org/doc/V9.1.1/stdlib//Corelib.Init.Datatypes}{\coqdocnotation{(}}\coqref{SECF.Imp.ANum}{\coqdocconstructor{ANum}} \coqdocvar{n1}\coqexternalref{::core scope:'(' x ',' x ',' '..' ',' x ')'}{http://rocq-prover.org/doc/V9.1.1/stdlib//Corelib.Init.Datatypes}{\coqdocnotation{,}} \coqref{SECF.Imp.ANum}{\coqdocconstructor{ANum}} \coqdocvar{n2}\coqexternalref{::core scope:'(' x ',' x ',' '..' ',' x ')'}{http://rocq-prover.org/doc/V9.1.1/stdlib//Corelib.Init.Datatypes}{\coqdocnotation{)}} \ensuremath{\Rightarrow}\coqdoceol
\coqdocindent{5.00em}
\coqdockw{if} \coqdocvar{n1} \coqexternalref{::nat scope:x '<=?' x}{http://rocq-prover.org/doc/V9.1.1/stdlib//Corelib.Init.Nat}{\coqdocnotation{<=?}} \coqdocvar{n2} \coqdockw{then} \coqref{SECF.Imp.::com scope:'<x7B' x 'x7D>'}{\coqdocnotation{<\{}}\coqref{SECF.Imp.:com::'false'}{\coqdocnotation{false}}\coqref{SECF.Imp.::com scope:'<x7B' x 'x7D>'}{\coqdocnotation{\}>}} \coqdockw{else} \coqref{SECF.Imp.::com scope:'<x7B' x 'x7D>'}{\coqdocnotation{<\{}}\coqref{SECF.Imp.:com::'true'}{\coqdocnotation{true}}\coqref{SECF.Imp.::com scope:'<x7B' x 'x7D>'}{\coqdocnotation{\}>}}\coqdoceol
\coqdocindent{3.00em}
\ensuremath{|} \coqexternalref{::core scope:'(' x ',' x ',' '..' ',' x ')'}{http://rocq-prover.org/doc/V9.1.1/stdlib//Corelib.Init.Datatypes}{\coqdocnotation{(}}\coqdocvar{a1'}\coqexternalref{::core scope:'(' x ',' x ',' '..' ',' x ')'}{http://rocq-prover.org/doc/V9.1.1/stdlib//Corelib.Init.Datatypes}{\coqdocnotation{,}} \coqdocvar{a2'}\coqexternalref{::core scope:'(' x ',' x ',' '..' ',' x ')'}{http://rocq-prover.org/doc/V9.1.1/stdlib//Corelib.Init.Datatypes}{\coqdocnotation{)}} \ensuremath{\Rightarrow}\coqdoceol
\coqdocindent{5.00em}
\coqref{SECF.Imp.::com scope:'<x7B' x 'x7D>'}{\coqdocnotation{<\{}} \coqdocvar{a1'} \coqref{SECF.Imp.:com::x '>' x}{\coqdocnotation{>}} \coqdocvar{a2'} \coqref{SECF.Imp.::com scope:'<x7B' x 'x7D>'}{\coqdocnotation{\}>}}\coqdoceol
\coqdocindent{3.00em}
\coqdockw{end}\coqdoceol
\coqdocindent{1.00em}
\ensuremath{|} \coqref{SECF.Imp.::com scope:'<x7B' x 'x7D>'}{\coqdocnotation{<\{}} \coqref{SECF.Imp.:com::'x7E' x}{\coqdocnotation{\ensuremath{\lnot}}} \coqdocvar{b1} \coqref{SECF.Imp.::com scope:'<x7B' x 'x7D>'}{\coqdocnotation{\}>}}  \ensuremath{\Rightarrow}\coqdoceol
\coqdocindent{3.00em}
\coqdockw{match} (\coqref{SECF.Equiv.fold constants bexp:97}{\coqdocdefinition{fold\_constants\_bexp}} \coqdocvar{b1}) \coqdockw{with}\coqdoceol
\coqdocindent{3.00em}
\ensuremath{|} \coqref{SECF.Imp.::com scope:'<x7B' x 'x7D>'}{\coqdocnotation{<\{}}\coqref{SECF.Imp.:com::'true'}{\coqdocnotation{true}}\coqref{SECF.Imp.::com scope:'<x7B' x 'x7D>'}{\coqdocnotation{\}>}} \ensuremath{\Rightarrow} \coqref{SECF.Imp.::com scope:'<x7B' x 'x7D>'}{\coqdocnotation{<\{}}\coqref{SECF.Imp.:com::'false'}{\coqdocnotation{false}}\coqref{SECF.Imp.::com scope:'<x7B' x 'x7D>'}{\coqdocnotation{\}>}}\coqdoceol
\coqdocindent{3.00em}
\ensuremath{|} \coqref{SECF.Imp.::com scope:'<x7B' x 'x7D>'}{\coqdocnotation{<\{}}\coqref{SECF.Imp.:com::'false'}{\coqdocnotation{false}}\coqref{SECF.Imp.::com scope:'<x7B' x 'x7D>'}{\coqdocnotation{\}>}} \ensuremath{\Rightarrow} \coqref{SECF.Imp.::com scope:'<x7B' x 'x7D>'}{\coqdocnotation{<\{}}\coqref{SECF.Imp.:com::'true'}{\coqdocnotation{true}}\coqref{SECF.Imp.::com scope:'<x7B' x 'x7D>'}{\coqdocnotation{\}>}}\coqdoceol
\coqdocindent{3.00em}
\ensuremath{|} \coqdocvar{b1'} \ensuremath{\Rightarrow} \coqref{SECF.Imp.::com scope:'<x7B' x 'x7D>'}{\coqdocnotation{<\{}} \coqref{SECF.Imp.:com::'x7E' x}{\coqdocnotation{\ensuremath{\lnot}}} \coqdocvar{b1'} \coqref{SECF.Imp.::com scope:'<x7B' x 'x7D>'}{\coqdocnotation{\}>}}\coqdoceol
\coqdocindent{3.00em}
\coqdockw{end}\coqdoceol
\coqdocindent{1.00em}
\ensuremath{|} \coqref{SECF.Imp.::com scope:'<x7B' x 'x7D>'}{\coqdocnotation{<\{}} \coqdocvar{b1} \coqref{SECF.Imp.:com::x 'x26x26' x}{\coqdocnotation{\&\&}} \coqdocvar{b2} \coqref{SECF.Imp.::com scope:'<x7B' x 'x7D>'}{\coqdocnotation{\}>}}  \ensuremath{\Rightarrow}\coqdoceol
\coqdocindent{3.00em}
\coqdockw{match} \coqexternalref{::core scope:'(' x ',' x ',' '..' ',' x ')'}{http://rocq-prover.org/doc/V9.1.1/stdlib//Corelib.Init.Datatypes}{\coqdocnotation{(}}\coqref{SECF.Equiv.fold constants bexp:97}{\coqdocdefinition{fold\_constants\_bexp}} \coqdocvar{b1}\coqexternalref{::core scope:'(' x ',' x ',' '..' ',' x ')'}{http://rocq-prover.org/doc/V9.1.1/stdlib//Corelib.Init.Datatypes}{\coqdocnotation{,}}\coqdoceol
\coqdocindent{6.50em}
\coqref{SECF.Equiv.fold constants bexp:97}{\coqdocdefinition{fold\_constants\_bexp}} \coqdocvar{b2}\coqexternalref{::core scope:'(' x ',' x ',' '..' ',' x ')'}{http://rocq-prover.org/doc/V9.1.1/stdlib//Corelib.Init.Datatypes}{\coqdocnotation{)}} \coqdockw{with}\coqdoceol
\coqdocindent{3.00em}
\ensuremath{|} \coqexternalref{::core scope:'(' x ',' x ',' '..' ',' x ')'}{http://rocq-prover.org/doc/V9.1.1/stdlib//Corelib.Init.Datatypes}{\coqdocnotation{(}}\coqref{SECF.Imp.::com scope:'<x7B' x 'x7D>'}{\coqdocnotation{<\{}}\coqref{SECF.Imp.:com::'true'}{\coqdocnotation{true}}\coqref{SECF.Imp.::com scope:'<x7B' x 'x7D>'}{\coqdocnotation{\}>}}\coqexternalref{::core scope:'(' x ',' x ',' '..' ',' x ')'}{http://rocq-prover.org/doc/V9.1.1/stdlib//Corelib.Init.Datatypes}{\coqdocnotation{,}} \coqref{SECF.Imp.::com scope:'<x7B' x 'x7D>'}{\coqdocnotation{<\{}}\coqref{SECF.Imp.:com::'true'}{\coqdocnotation{true}}\coqref{SECF.Imp.::com scope:'<x7B' x 'x7D>'}{\coqdocnotation{\}>}}\coqexternalref{::core scope:'(' x ',' x ',' '..' ',' x ')'}{http://rocq-prover.org/doc/V9.1.1/stdlib//Corelib.Init.Datatypes}{\coqdocnotation{)}} \ensuremath{\Rightarrow} \coqref{SECF.Imp.::com scope:'<x7B' x 'x7D>'}{\coqdocnotation{<\{}}\coqref{SECF.Imp.:com::'true'}{\coqdocnotation{true}}\coqref{SECF.Imp.::com scope:'<x7B' x 'x7D>'}{\coqdocnotation{\}>}}\coqdoceol
\coqdocindent{3.00em}
\ensuremath{|} \coqexternalref{::core scope:'(' x ',' x ',' '..' ',' x ')'}{http://rocq-prover.org/doc/V9.1.1/stdlib//Corelib.Init.Datatypes}{\coqdocnotation{(}}\coqref{SECF.Imp.::com scope:'<x7B' x 'x7D>'}{\coqdocnotation{<\{}}\coqref{SECF.Imp.:com::'true'}{\coqdocnotation{true}}\coqref{SECF.Imp.::com scope:'<x7B' x 'x7D>'}{\coqdocnotation{\}>}}\coqexternalref{::core scope:'(' x ',' x ',' '..' ',' x ')'}{http://rocq-prover.org/doc/V9.1.1/stdlib//Corelib.Init.Datatypes}{\coqdocnotation{,}} \coqref{SECF.Imp.::com scope:'<x7B' x 'x7D>'}{\coqdocnotation{<\{}}\coqref{SECF.Imp.:com::'false'}{\coqdocnotation{false}}\coqref{SECF.Imp.::com scope:'<x7B' x 'x7D>'}{\coqdocnotation{\}>}}\coqexternalref{::core scope:'(' x ',' x ',' '..' ',' x ')'}{http://rocq-prover.org/doc/V9.1.1/stdlib//Corelib.Init.Datatypes}{\coqdocnotation{)}} \ensuremath{\Rightarrow} \coqref{SECF.Imp.::com scope:'<x7B' x 'x7D>'}{\coqdocnotation{<\{}}\coqref{SECF.Imp.:com::'false'}{\coqdocnotation{false}}\coqref{SECF.Imp.::com scope:'<x7B' x 'x7D>'}{\coqdocnotation{\}>}}\coqdoceol
\coqdocindent{3.00em}
\ensuremath{|} \coqexternalref{::core scope:'(' x ',' x ',' '..' ',' x ')'}{http://rocq-prover.org/doc/V9.1.1/stdlib//Corelib.Init.Datatypes}{\coqdocnotation{(}}\coqref{SECF.Imp.::com scope:'<x7B' x 'x7D>'}{\coqdocnotation{<\{}}\coqref{SECF.Imp.:com::'false'}{\coqdocnotation{false}}\coqref{SECF.Imp.::com scope:'<x7B' x 'x7D>'}{\coqdocnotation{\}>}}\coqexternalref{::core scope:'(' x ',' x ',' '..' ',' x ')'}{http://rocq-prover.org/doc/V9.1.1/stdlib//Corelib.Init.Datatypes}{\coqdocnotation{,}} \coqref{SECF.Imp.::com scope:'<x7B' x 'x7D>'}{\coqdocnotation{<\{}}\coqref{SECF.Imp.:com::'true'}{\coqdocnotation{true}}\coqref{SECF.Imp.::com scope:'<x7B' x 'x7D>'}{\coqdocnotation{\}>}}\coqexternalref{::core scope:'(' x ',' x ',' '..' ',' x ')'}{http://rocq-prover.org/doc/V9.1.1/stdlib//Corelib.Init.Datatypes}{\coqdocnotation{)}} \ensuremath{\Rightarrow} \coqref{SECF.Imp.::com scope:'<x7B' x 'x7D>'}{\coqdocnotation{<\{}}\coqref{SECF.Imp.:com::'false'}{\coqdocnotation{false}}\coqref{SECF.Imp.::com scope:'<x7B' x 'x7D>'}{\coqdocnotation{\}>}}\coqdoceol
\coqdocindent{3.00em}
\ensuremath{|} \coqexternalref{::core scope:'(' x ',' x ',' '..' ',' x ')'}{http://rocq-prover.org/doc/V9.1.1/stdlib//Corelib.Init.Datatypes}{\coqdocnotation{(}}\coqref{SECF.Imp.::com scope:'<x7B' x 'x7D>'}{\coqdocnotation{<\{}}\coqref{SECF.Imp.:com::'false'}{\coqdocnotation{false}}\coqref{SECF.Imp.::com scope:'<x7B' x 'x7D>'}{\coqdocnotation{\}>}}\coqexternalref{::core scope:'(' x ',' x ',' '..' ',' x ')'}{http://rocq-prover.org/doc/V9.1.1/stdlib//Corelib.Init.Datatypes}{\coqdocnotation{,}} \coqref{SECF.Imp.::com scope:'<x7B' x 'x7D>'}{\coqdocnotation{<\{}}\coqref{SECF.Imp.:com::'false'}{\coqdocnotation{false}}\coqref{SECF.Imp.::com scope:'<x7B' x 'x7D>'}{\coqdocnotation{\}>}}\coqexternalref{::core scope:'(' x ',' x ',' '..' ',' x ')'}{http://rocq-prover.org/doc/V9.1.1/stdlib//Corelib.Init.Datatypes}{\coqdocnotation{)}} \ensuremath{\Rightarrow} \coqref{SECF.Imp.::com scope:'<x7B' x 'x7D>'}{\coqdocnotation{<\{}}\coqref{SECF.Imp.:com::'false'}{\coqdocnotation{false}}\coqref{SECF.Imp.::com scope:'<x7B' x 'x7D>'}{\coqdocnotation{\}>}}\coqdoceol
\coqdocindent{3.00em}
\ensuremath{|} \coqexternalref{::core scope:'(' x ',' x ',' '..' ',' x ')'}{http://rocq-prover.org/doc/V9.1.1/stdlib//Corelib.Init.Datatypes}{\coqdocnotation{(}}\coqdocvar{b1'}\coqexternalref{::core scope:'(' x ',' x ',' '..' ',' x ')'}{http://rocq-prover.org/doc/V9.1.1/stdlib//Corelib.Init.Datatypes}{\coqdocnotation{,}} \coqdocvar{b2'}\coqexternalref{::core scope:'(' x ',' x ',' '..' ',' x ')'}{http://rocq-prover.org/doc/V9.1.1/stdlib//Corelib.Init.Datatypes}{\coqdocnotation{)}} \ensuremath{\Rightarrow} \coqref{SECF.Imp.::com scope:'<x7B' x 'x7D>'}{\coqdocnotation{<\{}} \coqdocvar{b1'} \coqref{SECF.Imp.:com::x 'x26x26' x}{\coqdocnotation{\&\&}} \coqdocvar{b2'} \coqref{SECF.Imp.::com scope:'<x7B' x 'x7D>'}{\coqdocnotation{\}>}}\coqdoceol
\coqdocindent{3.00em}
\coqdockw{end}\coqdoceol
\coqdocindent{1.00em}
\coqdockw{end}.\coqdoceol
\coqdocemptyline
\coqdocnoindent
\coqdockw{Example} \coqdef{SECF.Equiv.fold bexp ex1}{fold\_bexp\_ex1}{\coqdocdefinition{fold\_bexp\_ex1}} :\coqdoceol
\coqdocindent{1.00em}
\coqref{SECF.Equiv.fold constants bexp}{\coqdocdefinition{fold\_constants\_bexp}} \coqref{SECF.Imp.::com scope:'<x7B' x 'x7D>'}{\coqdocnotation{<\{}} \coqref{SECF.Imp.:com::'true'}{\coqdocnotation{true}} \coqref{SECF.Imp.:com::x 'x26x26' x}{\coqdocnotation{\&\&}} \coqref{SECF.Imp.:com::'x7E' x}{\coqdocnotation{\ensuremath{\lnot}}}\coqref{SECF.Imp.:com::'(' x ')'}{\coqdocnotation{(}}\coqref{SECF.Imp.:com::'false'}{\coqdocnotation{false}} \coqref{SECF.Imp.:com::x 'x26x26' x}{\coqdocnotation{\&\&}} \coqref{SECF.Imp.:com::'true'}{\coqdocnotation{true}}\coqref{SECF.Imp.:com::'(' x ')'}{\coqdocnotation{)}} \coqref{SECF.Imp.::com scope:'<x7B' x 'x7D>'}{\coqdocnotation{\}>}}\coqdoceol
\coqdocindent{1.00em}
\coqexternalref{::type scope:x '=' x}{http://rocq-prover.org/doc/V9.1.1/stdlib//Corelib.Init.Logic}{\coqdocnotation{=}} \coqref{SECF.Imp.::com scope:'<x7B' x 'x7D>'}{\coqdocnotation{<\{}} \coqref{SECF.Imp.:com::'true'}{\coqdocnotation{true}} \coqref{SECF.Imp.::com scope:'<x7B' x 'x7D>'}{\coqdocnotation{\}>}}.\coqdoceol
\coqdocnoindent
\coqdockw{Proof}. \coqdoctac{reflexivity}. \coqdockw{Qed}.\coqdoceol
\coqdocemptyline
\coqdocnoindent
\coqdockw{Example} \coqdef{SECF.Equiv.fold bexp ex2}{fold\_bexp\_ex2}{\coqdocdefinition{fold\_bexp\_ex2}} :\coqdoceol
\coqdocindent{1.00em}
\coqref{SECF.Equiv.fold constants bexp}{\coqdocdefinition{fold\_constants\_bexp}} \coqref{SECF.Imp.::com scope:'<x7B' x 'x7D>'}{\coqdocnotation{<\{}} \coqref{SECF.Imp.:com::'(' x ')'}{\coqdocnotation{(}}\coqref{SECF.Imp.X}{\coqdocdefinition{X}} \coqref{SECF.Imp.:com::x '=' x}{\coqdocnotation{=}} \coqref{SECF.Imp.Y}{\coqdocdefinition{Y}}\coqref{SECF.Imp.:com::'(' x ')'}{\coqdocnotation{)}} \coqref{SECF.Imp.:com::x 'x26x26' x}{\coqdocnotation{\&\&}} \coqref{SECF.Imp.:com::'(' x ')'}{\coqdocnotation{(}}0 \coqref{SECF.Imp.:com::x '=' x}{\coqdocnotation{=}} \coqref{SECF.Imp.:com::'(' x ')'}{\coqdocnotation{(}}2 \coqref{SECF.Imp.:com::x '-' x}{\coqdocnotation{-}} \coqref{SECF.Imp.:com::'(' x ')'}{\coqdocnotation{(}}1 \coqref{SECF.Imp.:com::x '+' x}{\coqdocnotation{+}} 1\coqref{SECF.Imp.:com::'(' x ')'}{\coqdocnotation{)))}} \coqref{SECF.Imp.::com scope:'<x7B' x 'x7D>'}{\coqdocnotation{\}>}}\coqdoceol
\coqdocindent{1.00em}
\coqexternalref{::type scope:x '=' x}{http://rocq-prover.org/doc/V9.1.1/stdlib//Corelib.Init.Logic}{\coqdocnotation{=}} \coqref{SECF.Imp.::com scope:'<x7B' x 'x7D>'}{\coqdocnotation{<\{}} \coqref{SECF.Imp.:com::'(' x ')'}{\coqdocnotation{(}}\coqref{SECF.Imp.X}{\coqdocdefinition{X}} \coqref{SECF.Imp.:com::x '=' x}{\coqdocnotation{=}} \coqref{SECF.Imp.Y}{\coqdocdefinition{Y}}\coqref{SECF.Imp.:com::'(' x ')'}{\coqdocnotation{)}} \coqref{SECF.Imp.:com::x 'x26x26' x}{\coqdocnotation{\&\&}} \coqref{SECF.Imp.:com::'true'}{\coqdocnotation{true}} \coqref{SECF.Imp.::com scope:'<x7B' x 'x7D>'}{\coqdocnotation{\}>}}.\coqdoceol
\coqdocnoindent
\coqdockw{Proof}. \coqdoctac{reflexivity}. \coqdockw{Qed}.\coqdoceol
\coqdocemptyline
\end{coqdoccode}
To fold constants in a command, we simply apply the
    appropriate folding functions on all embedded expressions. 
\begin{coqdoccode}
\coqdocemptyline
\coqdocnoindent
\coqdockw{Fixpoint} \coqdef{SECF.Equiv.fold constants com}{fold\_constants\_com}{\coqdocdefinition{fold\_constants\_com}} (\coqdef{SECF.Equiv.c:99}{c}{\coqdocbinder{c}} : \coqref{SECF.Imp.com}{\coqdocinductive{com}}) : \coqref{SECF.Imp.com}{\coqdocinductive{com}} :=\coqdoceol
\coqdocindent{1.00em}
\coqdockw{match} \coqref{SECF.Equiv.c:99}{\coqdocvariable{c}} \coqdockw{with}\coqdoceol
\coqdocindent{1.00em}
\ensuremath{|} \coqref{SECF.Imp.::com scope:'<x7B' x 'x7D>'}{\coqdocnotation{<\{}} \coqref{SECF.Imp.:com:com scope:'skip'}{\coqdocnotation{skip}} \coqref{SECF.Imp.::com scope:'<x7B' x 'x7D>'}{\coqdocnotation{\}>}} \ensuremath{\Rightarrow}\coqdoceol
\coqdocindent{3.00em}
\coqref{SECF.Imp.::com scope:'<x7B' x 'x7D>'}{\coqdocnotation{<\{}} \coqref{SECF.Imp.:com:com scope:'skip'}{\coqdocnotation{skip}} \coqref{SECF.Imp.::com scope:'<x7B' x 'x7D>'}{\coqdocnotation{\}>}}\coqdoceol
\coqdocindent{1.00em}
\ensuremath{|} \coqref{SECF.Imp.::com scope:'<x7B' x 'x7D>'}{\coqdocnotation{<\{}} \coqdocvar{x} \coqref{SECF.Imp.:com:com scope:x ':=' x}{\coqdocnotation{:=}} \coqdocvar{a} \coqref{SECF.Imp.::com scope:'<x7B' x 'x7D>'}{\coqdocnotation{\}>}} \ensuremath{\Rightarrow}\coqdoceol
\coqdocindent{3.00em}
\coqref{SECF.Imp.::com scope:'<x7B' x 'x7D>'}{\coqdocnotation{<\{}} \coqdocvar{x} \coqref{SECF.Imp.:com:com scope:x ':=' x}{\coqdocnotation{:=}} \coqref{SECF.Imp.:com::'(' x ')'}{\coqdocnotation{(}}\coqref{SECF.Equiv.fold constants aexp}{\coqdocdefinition{fold\_constants\_aexp}} \coqdocvar{a}\coqref{SECF.Imp.:com::'(' x ')'}{\coqdocnotation{)}} \coqref{SECF.Imp.::com scope:'<x7B' x 'x7D>'}{\coqdocnotation{\}>}}\coqdoceol
\coqdocindent{1.00em}
\ensuremath{|} \coqref{SECF.Imp.::com scope:'<x7B' x 'x7D>'}{\coqdocnotation{<\{}} \coqdocvar{c1} \coqref{SECF.Imp.:com:com scope:x ';' x}{\coqdocnotation{;}} \coqdocvar{c2} \coqref{SECF.Imp.::com scope:'<x7B' x 'x7D>'}{\coqdocnotation{\}>}}  \ensuremath{\Rightarrow}\coqdoceol
\coqdocindent{3.00em}
\coqref{SECF.Imp.::com scope:'<x7B' x 'x7D>'}{\coqdocnotation{<\{}} \coqref{SECF.Equiv.fold constants com:100}{\coqdocdefinition{fold\_constants\_com}} \coqdocvar{c1} \coqref{SECF.Imp.:com:com scope:x ';' x}{\coqdocnotation{;}} \coqref{SECF.Equiv.fold constants com:100}{\coqdocdefinition{fold\_constants\_com}} \coqdocvar{c2} \coqref{SECF.Imp.::com scope:'<x7B' x 'x7D>'}{\coqdocnotation{\}>}}\coqdoceol
\coqdocindent{1.00em}
\ensuremath{|} \coqref{SECF.Imp.::com scope:'<x7B' x 'x7D>'}{\coqdocnotation{<\{}} \coqref{SECF.Imp.:com:com scope:'if' x 'then' x 'else' x 'end'}{\coqdocnotation{if}} \coqdocvar{b} \coqref{SECF.Imp.:com:com scope:'if' x 'then' x 'else' x 'end'}{\coqdocnotation{then}} \coqdocvar{c1} \coqref{SECF.Imp.:com:com scope:'if' x 'then' x 'else' x 'end'}{\coqdocnotation{else}} \coqdocvar{c2} \coqref{SECF.Imp.:com:com scope:'if' x 'then' x 'else' x 'end'}{\coqdocnotation{end}} \coqref{SECF.Imp.::com scope:'<x7B' x 'x7D>'}{\coqdocnotation{\}>}} \ensuremath{\Rightarrow}\coqdoceol
\coqdocindent{3.00em}
\coqdockw{match} \coqref{SECF.Equiv.fold constants bexp}{\coqdocdefinition{fold\_constants\_bexp}} \coqdocvar{b} \coqdockw{with}\coqdoceol
\coqdocindent{3.00em}
\ensuremath{|} \coqref{SECF.Imp.::com scope:'<x7B' x 'x7D>'}{\coqdocnotation{<\{}}\coqref{SECF.Imp.:com::'true'}{\coqdocnotation{true}}\coqref{SECF.Imp.::com scope:'<x7B' x 'x7D>'}{\coqdocnotation{\}>}}  \ensuremath{\Rightarrow} \coqref{SECF.Equiv.fold constants com:100}{\coqdocdefinition{fold\_constants\_com}} \coqdocvar{c1}\coqdoceol
\coqdocindent{3.00em}
\ensuremath{|} \coqref{SECF.Imp.::com scope:'<x7B' x 'x7D>'}{\coqdocnotation{<\{}}\coqref{SECF.Imp.:com::'false'}{\coqdocnotation{false}}\coqref{SECF.Imp.::com scope:'<x7B' x 'x7D>'}{\coqdocnotation{\}>}} \ensuremath{\Rightarrow} \coqref{SECF.Equiv.fold constants com:100}{\coqdocdefinition{fold\_constants\_com}} \coqdocvar{c2}\coqdoceol
\coqdocindent{3.00em}
\ensuremath{|} \coqdocvar{b'} \ensuremath{\Rightarrow} \coqref{SECF.Imp.::com scope:'<x7B' x 'x7D>'}{\coqdocnotation{<\{}} \coqref{SECF.Imp.:com:com scope:'if' x 'then' x 'else' x 'end'}{\coqdocnotation{if}} \coqdocvar{b'} \coqref{SECF.Imp.:com:com scope:'if' x 'then' x 'else' x 'end'}{\coqdocnotation{then}} \coqref{SECF.Equiv.fold constants com:100}{\coqdocdefinition{fold\_constants\_com}} \coqdocvar{c1}\coqdoceol
\coqdocindent{11.50em}
\coqref{SECF.Imp.:com:com scope:'if' x 'then' x 'else' x 'end'}{\coqdocnotation{else}} \coqref{SECF.Equiv.fold constants com:100}{\coqdocdefinition{fold\_constants\_com}} \coqdocvar{c2} \coqref{SECF.Imp.:com:com scope:'if' x 'then' x 'else' x 'end'}{\coqdocnotation{end}}\coqref{SECF.Imp.::com scope:'<x7B' x 'x7D>'}{\coqdocnotation{\}>}}\coqdoceol
\coqdocindent{3.00em}
\coqdockw{end}\coqdoceol
\coqdocindent{1.00em}
\ensuremath{|} \coqref{SECF.Imp.::com scope:'<x7B' x 'x7D>'}{\coqdocnotation{<\{}} \coqref{SECF.Imp.:com:com scope:'while' x 'do' x 'end'}{\coqdocnotation{while}} \coqdocvar{b} \coqref{SECF.Imp.:com:com scope:'while' x 'do' x 'end'}{\coqdocnotation{do}} \coqdocvar{c1} \coqref{SECF.Imp.:com:com scope:'while' x 'do' x 'end'}{\coqdocnotation{end}} \coqref{SECF.Imp.::com scope:'<x7B' x 'x7D>'}{\coqdocnotation{\}>}} \ensuremath{\Rightarrow}\coqdoceol
\coqdocindent{3.00em}
\coqdockw{match} \coqref{SECF.Equiv.fold constants bexp}{\coqdocdefinition{fold\_constants\_bexp}} \coqdocvar{b} \coqdockw{with}\coqdoceol
\coqdocindent{3.00em}
\ensuremath{|} \coqref{SECF.Imp.::com scope:'<x7B' x 'x7D>'}{\coqdocnotation{<\{}}\coqref{SECF.Imp.:com::'true'}{\coqdocnotation{true}}\coqref{SECF.Imp.::com scope:'<x7B' x 'x7D>'}{\coqdocnotation{\}>}} \ensuremath{\Rightarrow} \coqref{SECF.Imp.::com scope:'<x7B' x 'x7D>'}{\coqdocnotation{<\{}} \coqref{SECF.Imp.:com:com scope:'while' x 'do' x 'end'}{\coqdocnotation{while}} \coqref{SECF.Imp.:com::'true'}{\coqdocnotation{true}} \coqref{SECF.Imp.:com:com scope:'while' x 'do' x 'end'}{\coqdocnotation{do}} \coqref{SECF.Imp.:com:com scope:'skip'}{\coqdocnotation{skip}} \coqref{SECF.Imp.:com:com scope:'while' x 'do' x 'end'}{\coqdocnotation{end}} \coqref{SECF.Imp.::com scope:'<x7B' x 'x7D>'}{\coqdocnotation{\}>}}\coqdoceol
\coqdocindent{3.00em}
\ensuremath{|} \coqref{SECF.Imp.::com scope:'<x7B' x 'x7D>'}{\coqdocnotation{<\{}}\coqref{SECF.Imp.:com::'false'}{\coqdocnotation{false}}\coqref{SECF.Imp.::com scope:'<x7B' x 'x7D>'}{\coqdocnotation{\}>}} \ensuremath{\Rightarrow} \coqref{SECF.Imp.::com scope:'<x7B' x 'x7D>'}{\coqdocnotation{<\{}} \coqref{SECF.Imp.:com:com scope:'skip'}{\coqdocnotation{skip}} \coqref{SECF.Imp.::com scope:'<x7B' x 'x7D>'}{\coqdocnotation{\}>}}\coqdoceol
\coqdocindent{3.00em}
\ensuremath{|} \coqdocvar{b'} \ensuremath{\Rightarrow} \coqref{SECF.Imp.::com scope:'<x7B' x 'x7D>'}{\coqdocnotation{<\{}} \coqref{SECF.Imp.:com:com scope:'while' x 'do' x 'end'}{\coqdocnotation{while}} \coqdocvar{b'} \coqref{SECF.Imp.:com:com scope:'while' x 'do' x 'end'}{\coqdocnotation{do}} \coqref{SECF.Imp.:com::'(' x ')'}{\coqdocnotation{(}}\coqref{SECF.Equiv.fold constants com:100}{\coqdocdefinition{fold\_constants\_com}} \coqdocvar{c1}\coqref{SECF.Imp.:com::'(' x ')'}{\coqdocnotation{)}} \coqref{SECF.Imp.:com:com scope:'while' x 'do' x 'end'}{\coqdocnotation{end}} \coqref{SECF.Imp.::com scope:'<x7B' x 'x7D>'}{\coqdocnotation{\}>}}\coqdoceol
\coqdocindent{3.00em}
\coqdockw{end}\coqdoceol
\coqdocindent{1.00em}
\coqdockw{end}.\coqdoceol
\coqdocemptyline
\coqdocnoindent
\coqdockw{Example} \coqdef{SECF.Equiv.fold com ex1}{fold\_com\_ex1}{\coqdocdefinition{fold\_com\_ex1}} :\coqdoceol
\coqdocindent{1.00em}
\coqref{SECF.Equiv.fold constants com}{\coqdocdefinition{fold\_constants\_com}}\coqdoceol
\coqdocindent{2.00em}
\coqdoceol
\coqdocindent{2.00em}
\coqref{SECF.Imp.::com scope:'<x7B' x 'x7D>'}{\coqdocnotation{<\{}} \coqref{SECF.Imp.X}{\coqdocdefinition{X}} \coqref{SECF.Imp.:com:com scope:x ':=' x}{\coqdocnotation{:=}} 4 \coqref{SECF.Imp.:com::x '+' x}{\coqdocnotation{+}} 5\coqref{SECF.Imp.:com:com scope:x ';' x}{\coqdocnotation{;}}\coqdoceol
\coqdocindent{3.50em}
\coqref{SECF.Imp.Y}{\coqdocdefinition{Y}} \coqref{SECF.Imp.:com:com scope:x ':=' x}{\coqdocnotation{:=}} \coqref{SECF.Imp.X}{\coqdocdefinition{X}} \coqref{SECF.Imp.:com::x '-' x}{\coqdocnotation{-}} 3\coqref{SECF.Imp.:com:com scope:x ';' x}{\coqdocnotation{;}}\coqdoceol
\coqdocindent{3.50em}
\coqref{SECF.Imp.:com:com scope:'if' x 'then' x 'else' x 'end'}{\coqdocnotation{if}} \coqref{SECF.Imp.:com::'(' x ')'}{\coqdocnotation{(}}\coqref{SECF.Imp.X}{\coqdocdefinition{X}} \coqref{SECF.Imp.:com::x '-' x}{\coqdocnotation{-}} \coqref{SECF.Imp.Y}{\coqdocdefinition{Y}}\coqref{SECF.Imp.:com::'(' x ')'}{\coqdocnotation{)}} \coqref{SECF.Imp.:com::x '=' x}{\coqdocnotation{=}} \coqref{SECF.Imp.:com::'(' x ')'}{\coqdocnotation{(}}2 \coqref{SECF.Imp.:com::x '+' x}{\coqdocnotation{+}} 4\coqref{SECF.Imp.:com::'(' x ')'}{\coqdocnotation{)}} \coqref{SECF.Imp.:com:com scope:'if' x 'then' x 'else' x 'end'}{\coqdocnotation{then}} \coqref{SECF.Imp.:com:com scope:'skip'}{\coqdocnotation{skip}}\coqdoceol
\coqdocindent{3.50em}
\coqref{SECF.Imp.:com:com scope:'if' x 'then' x 'else' x 'end'}{\coqdocnotation{else}} \coqref{SECF.Imp.Y}{\coqdocdefinition{Y}} \coqref{SECF.Imp.:com:com scope:x ':=' x}{\coqdocnotation{:=}} 0 \coqref{SECF.Imp.:com:com scope:'if' x 'then' x 'else' x 'end'}{\coqdocnotation{end}}\coqref{SECF.Imp.:com:com scope:x ';' x}{\coqdocnotation{;}}\coqdoceol
\coqdocindent{3.50em}
\coqref{SECF.Imp.:com:com scope:'if' x 'then' x 'else' x 'end'}{\coqdocnotation{if}} 0 \coqref{SECF.Imp.:com::x '<=' x}{\coqdocnotation{\ensuremath{\le}}} \coqref{SECF.Imp.:com::'(' x ')'}{\coqdocnotation{(}}4 \coqref{SECF.Imp.:com::x '-' x}{\coqdocnotation{-}} \coqref{SECF.Imp.:com::'(' x ')'}{\coqdocnotation{(}}2 \coqref{SECF.Imp.:com::x '+' x}{\coqdocnotation{+}} 1\coqref{SECF.Imp.:com::'(' x ')'}{\coqdocnotation{))}} \coqref{SECF.Imp.:com:com scope:'if' x 'then' x 'else' x 'end'}{\coqdocnotation{then}} \coqref{SECF.Imp.Y}{\coqdocdefinition{Y}} \coqref{SECF.Imp.:com:com scope:x ':=' x}{\coqdocnotation{:=}} 0\coqdoceol
\coqdocindent{3.50em}
\coqref{SECF.Imp.:com:com scope:'if' x 'then' x 'else' x 'end'}{\coqdocnotation{else}} \coqref{SECF.Imp.:com:com scope:'skip'}{\coqdocnotation{skip}} \coqref{SECF.Imp.:com:com scope:'if' x 'then' x 'else' x 'end'}{\coqdocnotation{end}}\coqref{SECF.Imp.:com:com scope:x ';' x}{\coqdocnotation{;}}\coqdoceol
\coqdocindent{3.50em}
\coqref{SECF.Imp.:com:com scope:'while' x 'do' x 'end'}{\coqdocnotation{while}} \coqref{SECF.Imp.Y}{\coqdocdefinition{Y}} \coqref{SECF.Imp.:com::x '=' x}{\coqdocnotation{=}} 0 \coqref{SECF.Imp.:com:com scope:'while' x 'do' x 'end'}{\coqdocnotation{do}}\coqdoceol
\coqdocindent{4.50em}
\coqref{SECF.Imp.X}{\coqdocdefinition{X}} \coqref{SECF.Imp.:com:com scope:x ':=' x}{\coqdocnotation{:=}} \coqref{SECF.Imp.X}{\coqdocdefinition{X}} \coqref{SECF.Imp.:com::x '+' x}{\coqdocnotation{+}} 1\coqdoceol
\coqdocindent{3.50em}
\coqref{SECF.Imp.:com:com scope:'while' x 'do' x 'end'}{\coqdocnotation{end}} \coqref{SECF.Imp.::com scope:'<x7B' x 'x7D>'}{\coqdocnotation{\}>}}\coqdoceol
\coqdocindent{1.00em}
\coqexternalref{::type scope:x '=' x}{http://rocq-prover.org/doc/V9.1.1/stdlib//Corelib.Init.Logic}{\coqdocnotation{=}} \coqdoceol
\coqdocindent{2.00em}
\coqref{SECF.Imp.::com scope:'<x7B' x 'x7D>'}{\coqdocnotation{<\{}} \coqref{SECF.Imp.X}{\coqdocdefinition{X}} \coqref{SECF.Imp.:com:com scope:x ':=' x}{\coqdocnotation{:=}} 9\coqref{SECF.Imp.:com:com scope:x ';' x}{\coqdocnotation{;}}\coqdoceol
\coqdocindent{3.50em}
\coqref{SECF.Imp.Y}{\coqdocdefinition{Y}} \coqref{SECF.Imp.:com:com scope:x ':=' x}{\coqdocnotation{:=}} \coqref{SECF.Imp.X}{\coqdocdefinition{X}} \coqref{SECF.Imp.:com::x '-' x}{\coqdocnotation{-}} 3\coqref{SECF.Imp.:com:com scope:x ';' x}{\coqdocnotation{;}}\coqdoceol
\coqdocindent{3.50em}
\coqref{SECF.Imp.:com:com scope:'if' x 'then' x 'else' x 'end'}{\coqdocnotation{if}} \coqref{SECF.Imp.:com::'(' x ')'}{\coqdocnotation{(}}\coqref{SECF.Imp.X}{\coqdocdefinition{X}} \coqref{SECF.Imp.:com::x '-' x}{\coqdocnotation{-}} \coqref{SECF.Imp.Y}{\coqdocdefinition{Y}}\coqref{SECF.Imp.:com::'(' x ')'}{\coqdocnotation{)}} \coqref{SECF.Imp.:com::x '=' x}{\coqdocnotation{=}} 6 \coqref{SECF.Imp.:com:com scope:'if' x 'then' x 'else' x 'end'}{\coqdocnotation{then}} \coqref{SECF.Imp.:com:com scope:'skip'}{\coqdocnotation{skip}}\coqdoceol
\coqdocindent{3.50em}
\coqref{SECF.Imp.:com:com scope:'if' x 'then' x 'else' x 'end'}{\coqdocnotation{else}} \coqref{SECF.Imp.Y}{\coqdocdefinition{Y}} \coqref{SECF.Imp.:com:com scope:x ':=' x}{\coqdocnotation{:=}} 0 \coqref{SECF.Imp.:com:com scope:'if' x 'then' x 'else' x 'end'}{\coqdocnotation{end}}\coqref{SECF.Imp.:com:com scope:x ';' x}{\coqdocnotation{;}}\coqdoceol
\coqdocindent{3.50em}
\coqref{SECF.Imp.Y}{\coqdocdefinition{Y}} \coqref{SECF.Imp.:com:com scope:x ':=' x}{\coqdocnotation{:=}} 0\coqref{SECF.Imp.:com:com scope:x ';' x}{\coqdocnotation{;}}\coqdoceol
\coqdocindent{3.50em}
\coqref{SECF.Imp.:com:com scope:'while' x 'do' x 'end'}{\coqdocnotation{while}} \coqref{SECF.Imp.Y}{\coqdocdefinition{Y}} \coqref{SECF.Imp.:com::x '=' x}{\coqdocnotation{=}} 0 \coqref{SECF.Imp.:com:com scope:'while' x 'do' x 'end'}{\coqdocnotation{do}}\coqdoceol
\coqdocindent{4.50em}
\coqref{SECF.Imp.X}{\coqdocdefinition{X}} \coqref{SECF.Imp.:com:com scope:x ':=' x}{\coqdocnotation{:=}} \coqref{SECF.Imp.X}{\coqdocdefinition{X}} \coqref{SECF.Imp.:com::x '+' x}{\coqdocnotation{+}} 1\coqdoceol
\coqdocindent{3.50em}
\coqref{SECF.Imp.:com:com scope:'while' x 'do' x 'end'}{\coqdocnotation{end}} \coqref{SECF.Imp.::com scope:'<x7B' x 'x7D>'}{\coqdocnotation{\}>}}.\coqdoceol
\coqdocnoindent
\coqdockw{Proof}. \coqdoctac{reflexivity}. \coqdockw{Qed}.\coqdoceol
\coqdocemptyline
\end{coqdoccode}
\subsection{Soundness of Constant Folding}



 Now we need to show that what we've done is correct. 

 Here's the proof for arithmetic expressions: 
\begin{coqdoccode}
\coqdocemptyline
\coqdocnoindent
\coqdockw{Theorem} \coqdef{SECF.Equiv.fold constants aexp sound}{fold\_constants\_aexp\_sound}{\coqdoclemma{fold\_constants\_aexp\_sound}} :\coqdoceol
\coqdocindent{1.00em}
\coqref{SECF.Equiv.atrans sound}{\coqdocdefinition{atrans\_sound}} \coqref{SECF.Equiv.fold constants aexp}{\coqdocdefinition{fold\_constants\_aexp}}.\coqdoceol
\coqdocnoindent
\coqdockw{Proof}.\coqdoceol
\coqdocindent{1.00em}
\coqdoctac{unfold} \coqref{SECF.Equiv.atrans sound}{\coqdocdefinition{atrans\_sound}}. \coqdoctac{intros} \coqdocvar{a}. \coqdoctac{unfold} \coqref{SECF.Equiv.aequiv}{\coqdocdefinition{aequiv}}. \coqdoctac{intros} \coqdocvar{st}.\coqdoceol
\coqdocindent{1.00em}
\coqdoctac{induction} \coqdocvar{a}; \coqdoctac{simpl};\coqdoceol
\coqdocindent{2.00em}
\coqdoceol
\coqdocindent{2.00em}
\coqdoctac{try} \coqdoctac{reflexivity};\coqdoceol
\coqdocindent{2.00em}
\coqdoceol
\coqdocindent{2.00em}
\coqdoctac{try} (\coqdoctac{destruct} (\coqref{SECF.Equiv.fold constants aexp}{\coqdocdefinition{fold\_constants\_aexp}} \coqdocvar{a1});\coqdoceol
\coqdocindent{4.50em}
\coqdoctac{destruct} (\coqref{SECF.Equiv.fold constants aexp}{\coqdocdefinition{fold\_constants\_aexp}} \coqdocvar{a2});\coqdoceol
\coqdocindent{4.50em}
\coqdoctac{rewrite} \coqdocvar{IHa1}; \coqdoctac{rewrite} \coqdocvar{IHa2}; \coqdoctac{reflexivity}). \coqdockw{Qed}.\coqdoceol
\coqdocemptyline
\end{coqdoccode}
\paragraph{Exercise: 3 stars, standard, optional (fold\_bexp\_Eq\_informal)}




    Here is an informal proof of the \coqref{SECF.SpecCT.BEq}{\coqdocdefinition{BEq}} case of the soundness
    argument for boolean expression constant folding.  Read it
    carefully and compare it to the formal proof that follows.  Then
    fill in the \coqref{SECF.SpecCT.BLe}{\coqdocdefinition{BLe}} case of the formal proof (without looking at the
    \coqref{SECF.SpecCT.BEq}{\coqdocdefinition{BEq}} case, if possible).


   \textit{Theorem}: The constant folding function for booleans,
   \coqref{SECF.Equiv.fold constants bexp}{\coqdocdefinition{fold\_constants\_bexp}}, is sound.


   \textit{Proof}: We must show that \coqdocvar{b} is equivalent to \coqref{SECF.Equiv.fold constants bexp}{\coqdocdefinition{fold\_constants\_bexp}} \coqdocvar{b},
   for all boolean expressions \coqdocvar{b}.  Proceed by induction on \coqdocvar{b}.  We
   show just the case where \coqdocvar{b} has the form \coqdocvar{a1} = \coqdocvar{a2}.


   In this case, we must show


       beval st <\{ a1 = a2 \}>
     = beval st (fold\_constants\_bexp <\{ a1 = a2 \}>).


   There are two cases to consider:



\begin{itemize}
\item  First, suppose \coqref{SECF.Equiv.fold constants aexp}{\coqdocdefinition{fold\_constants\_aexp}} \coqdocvar{a1} = \coqref{SECF.SpecCT.ANum}{\coqdocconstructor{ANum}} \coqdocvar{n1} and
       \coqref{SECF.Equiv.fold constants aexp}{\coqdocdefinition{fold\_constants\_aexp}} \coqdocvar{a2} = \coqref{SECF.SpecCT.ANum}{\coqdocconstructor{ANum}} \coqdocvar{n2} for some \coqdocvar{n1} and \coqdocvar{n2}.


       In this case, we have


           fold\_constants\_bexp <\{ a1 = a2 \}>
         = if n1 =? n2 then <\{true\}> else <\{false\}>


       and


           beval st <\{a1 = a2\}>
         = (aeval st a1) =? (aeval st a2).


       By the soundness of constant folding for arithmetic
       expressions (Lemma \coqref{SECF.Equiv.fold constants aexp sound}{\coqdoclemma{fold\_constants\_aexp\_sound}}), we know


           aeval st a1
         = aeval st (fold\_constants\_aexp a1)
         = aeval st (ANum n1)
         = n1


       and


           aeval st a2
         = aeval st (fold\_constants\_aexp a2)
         = aeval st (ANum n2)
         = n2,


       so


           beval st <\{a1 = a2\}>
         = (aeval a1) =? (aeval a2)
         = n1 =? n2.


       Also, it is easy to see (by considering the cases \coqdocvar{n1} = \coqdocvar{n2} and
       \coqdocvar{n1} \ensuremath{\not=} \coqdocvar{n2} separately) that


           beval st (if n1 =? n2 then <\{true\}> else <\{false\}>)
         = if n1 =? n2 then beval st <\{true\}> else beval st <\{false\}>
         = if n1 =? n2 then true else false
         = n1 =? n2.


       So


           beval st (<\{ a1 = a2 \}>)
         = n1 =? n2.
         = beval st (if n1 =? n2 then <\{true\}> else <\{false\}>),


       as required.



\item  Otherwise, one of \coqref{SECF.Equiv.fold constants aexp}{\coqdocdefinition{fold\_constants\_aexp}} \coqdocvar{a1} and
       \coqref{SECF.Equiv.fold constants aexp}{\coqdocdefinition{fold\_constants\_aexp}} \coqdocvar{a2} is not a constant.  In this case, we
       must show


           beval st <\{a1 = a2\}>
         = beval st (<\{ (fold\_constants\_aexp a1) =
                         (fold\_constants\_aexp a2) \}>),


       which, by the definition of \coqref{SECF.Imp.beval}{\coqdocdefinition{beval}}, is the same as showing


           (aeval st a1) =? (aeval st a2)
         = (aeval st (fold\_constants\_aexp a1)) =?
                   (aeval st (fold\_constants\_aexp a2)).


       But the soundness of constant folding for arithmetic
       expressions (\coqref{SECF.Equiv.fold constants aexp sound}{\coqdoclemma{fold\_constants\_aexp\_sound}}) gives us


         aeval st a1 = aeval st (fold\_constants\_aexp a1)
         aeval st a2 = aeval st (fold\_constants\_aexp a2),


       completing the case.  \ensuremath{\Box} 

\end{itemize}
\begin{coqdoccode}
\coqdocemptyline
\coqdocnoindent
\coqdockw{Theorem} \coqdef{SECF.Equiv.fold constants bexp sound}{fold\_constants\_bexp\_sound}{\coqdoclemma{fold\_constants\_bexp\_sound}}:\coqdoceol
\coqdocindent{1.00em}
\coqref{SECF.Equiv.btrans sound}{\coqdocdefinition{btrans\_sound}} \coqref{SECF.Equiv.fold constants bexp}{\coqdocdefinition{fold\_constants\_bexp}}.\coqdoceol
\coqdocnoindent
\coqdockw{Proof}.\coqdoceol
\coqdocindent{1.00em}
\coqdoctac{unfold} \coqref{SECF.Equiv.btrans sound}{\coqdocdefinition{btrans\_sound}}. \coqdoctac{intros} \coqdocvar{b}. \coqdoctac{unfold} \coqref{SECF.Equiv.bequiv}{\coqdocdefinition{bequiv}}. \coqdoctac{intros} \coqdocvar{st}.\coqdoceol
\coqdocindent{1.00em}
\coqdoctac{induction} \coqdocvar{b};\coqdoceol
\coqdocindent{2.00em}
\coqdoceol
\coqdocindent{2.00em}
\coqdoctac{try} \coqdoctac{reflexivity}.\coqdoceol
\coqdocindent{1.00em}
- \coqdoceol
\coqdocindent{2.00em}
\coqdoctac{simpl}.\coqdoceol
\coqdocindent{2.00em}
\coqdocvar{remember} (\coqref{SECF.Equiv.fold constants aexp}{\coqdocdefinition{fold\_constants\_aexp}} \coqdocvar{a1}) \coqdockw{as} \coqdocvar{a1'} \coqdocvar{eqn}:\coqdocvar{Heqa1'}.\coqdoceol
\coqdocindent{2.00em}
\coqdocvar{remember} (\coqref{SECF.Equiv.fold constants aexp}{\coqdocdefinition{fold\_constants\_aexp}} \coqdocvar{a2}) \coqdockw{as} \coqdocvar{a2'} \coqdocvar{eqn}:\coqdocvar{Heqa2'}.\coqdoceol
\coqdocindent{2.00em}
\coqdoctac{replace} (\coqref{SECF.Imp.aeval}{\coqdocdefinition{aeval}} \coqdocvar{st} \coqdocvar{a1}) \coqdockw{with} (\coqref{SECF.Imp.aeval}{\coqdocdefinition{aeval}} \coqdocvar{st} \coqdocvar{a1'}) \coqdoctac{by}\coqdoceol
\coqdocindent{3.50em}
(\coqdoctac{subst} \coqdocvar{a1'}; \coqdoctac{rewrite} \ensuremath{\leftarrow} \coqref{SECF.Equiv.fold constants aexp sound}{\coqdoclemma{fold\_constants\_aexp\_sound}}; \coqdoctac{reflexivity}).\coqdoceol
\coqdocindent{2.00em}
\coqdoctac{replace} (\coqref{SECF.Imp.aeval}{\coqdocdefinition{aeval}} \coqdocvar{st} \coqdocvar{a2}) \coqdockw{with} (\coqref{SECF.Imp.aeval}{\coqdocdefinition{aeval}} \coqdocvar{st} \coqdocvar{a2'}) \coqdoctac{by}\coqdoceol
\coqdocindent{3.50em}
(\coqdoctac{subst} \coqdocvar{a2'}; \coqdoctac{rewrite} \ensuremath{\leftarrow} \coqref{SECF.Equiv.fold constants aexp sound}{\coqdoclemma{fold\_constants\_aexp\_sound}}; \coqdoctac{reflexivity}).\coqdoceol
\coqdocindent{2.00em}
\coqdoctac{destruct} \coqdocvar{a1'}; \coqdoctac{destruct} \coqdocvar{a2'}; \coqdoctac{try} \coqdoctac{reflexivity}.\coqdoceol
\coqdocindent{3.00em}
\coqdoctac{simpl}. \coqdoctac{destruct} (\coqdocvar{n} \coqexternalref{::nat scope:x '=?' x}{http://rocq-prover.org/doc/V9.1.1/stdlib//Corelib.Init.Nat}{\coqdocnotation{=?}} \coqdocvar{n0}); \coqdoctac{reflexivity}.\coqdoceol
\coqdocindent{1.00em}
- \coqdoceol
\coqdocindent{2.00em}
\coqdoctac{simpl}.\coqdoceol
\coqdocindent{2.00em}
\coqdocvar{remember} (\coqref{SECF.Equiv.fold constants aexp}{\coqdocdefinition{fold\_constants\_aexp}} \coqdocvar{a1}) \coqdockw{as} \coqdocvar{a1'} \coqdocvar{eqn}:\coqdocvar{Heqa1'}.\coqdoceol
\coqdocindent{2.00em}
\coqdocvar{remember} (\coqref{SECF.Equiv.fold constants aexp}{\coqdocdefinition{fold\_constants\_aexp}} \coqdocvar{a2}) \coqdockw{as} \coqdocvar{a2'} \coqdocvar{eqn}:\coqdocvar{Heqa2'}.\coqdoceol
\coqdocindent{2.00em}
\coqdoctac{replace} (\coqref{SECF.Imp.aeval}{\coqdocdefinition{aeval}} \coqdocvar{st} \coqdocvar{a1}) \coqdockw{with} (\coqref{SECF.Imp.aeval}{\coqdocdefinition{aeval}} \coqdocvar{st} \coqdocvar{a1'}) \coqdoctac{by}\coqdoceol
\coqdocindent{3.50em}
(\coqdoctac{subst} \coqdocvar{a1'}; \coqdoctac{rewrite} \ensuremath{\leftarrow} \coqref{SECF.Equiv.fold constants aexp sound}{\coqdoclemma{fold\_constants\_aexp\_sound}}; \coqdoctac{reflexivity}).\coqdoceol
\coqdocindent{2.00em}
\coqdoctac{replace} (\coqref{SECF.Imp.aeval}{\coqdocdefinition{aeval}} \coqdocvar{st} \coqdocvar{a2}) \coqdockw{with} (\coqref{SECF.Imp.aeval}{\coqdocdefinition{aeval}} \coqdocvar{st} \coqdocvar{a2'}) \coqdoctac{by}\coqdoceol
\coqdocindent{3.50em}
(\coqdoctac{subst} \coqdocvar{a2'}; \coqdoctac{rewrite} \ensuremath{\leftarrow} \coqref{SECF.Equiv.fold constants aexp sound}{\coqdoclemma{fold\_constants\_aexp\_sound}}; \coqdoctac{reflexivity}).\coqdoceol
\coqdocindent{2.00em}
\coqdoctac{destruct} \coqdocvar{a1'}; \coqdoctac{destruct} \coqdocvar{a2'}; \coqdoctac{try} \coqdoctac{reflexivity}.\coqdoceol
\coqdocindent{3.00em}
\coqdoctac{simpl}. \coqdoctac{destruct} (\coqdocvar{n} \coqexternalref{::nat scope:x '=?' x}{http://rocq-prover.org/doc/V9.1.1/stdlib//Corelib.Init.Nat}{\coqdocnotation{=?}} \coqdocvar{n0}); \coqdoctac{reflexivity}.\coqdoceol
\coqdocindent{1.00em}
- \coqdoceol
\coqdocindent{2.00em}
 \coqdocvar{admit}.\coqdoceol
\coqdocindent{1.00em}
- \coqdoceol
\coqdocindent{2.00em}
 \coqdocvar{admit}.\coqdoceol
\coqdocindent{1.00em}
- \coqdoceol
\coqdocindent{2.00em}
\coqdoctac{simpl}. \coqdocvar{remember} (\coqref{SECF.Equiv.fold constants bexp}{\coqdocdefinition{fold\_constants\_bexp}} \coqdocvar{b}) \coqdockw{as} \coqdocvar{b'} \coqdocvar{eqn}:\coqdocvar{Heqb'}.\coqdoceol
\coqdocindent{2.00em}
\coqdoctac{rewrite} \coqdocvar{IHb}.\coqdoceol
\coqdocindent{2.00em}
\coqdoctac{destruct} \coqdocvar{b'}; \coqdoctac{reflexivity}.\coqdoceol
\coqdocindent{1.00em}
- \coqdoceol
\coqdocindent{2.00em}
\coqdoctac{simpl}.\coqdoceol
\coqdocindent{2.00em}
\coqdocvar{remember} (\coqref{SECF.Equiv.fold constants bexp}{\coqdocdefinition{fold\_constants\_bexp}} \coqdocvar{b1}) \coqdockw{as} \coqdocvar{b1'} \coqdocvar{eqn}:\coqdocvar{Heqb1'}.\coqdoceol
\coqdocindent{2.00em}
\coqdocvar{remember} (\coqref{SECF.Equiv.fold constants bexp}{\coqdocdefinition{fold\_constants\_bexp}} \coqdocvar{b2}) \coqdockw{as} \coqdocvar{b2'} \coqdocvar{eqn}:\coqdocvar{Heqb2'}.\coqdoceol
\coqdocindent{2.00em}
\coqdoctac{rewrite} \coqdocvar{IHb1}. \coqdoctac{rewrite} \coqdocvar{IHb2}.\coqdoceol
\coqdocindent{2.00em}
\coqdoctac{destruct} \coqdocvar{b1'}; \coqdoctac{destruct} \coqdocvar{b2'}; \coqdoctac{reflexivity}.\coqdoceol
 \coqdocvar{Admitted}.\coqdoceol
\end{coqdoccode}
\ensuremath{\Box} 

\paragraph{Exercise: 3 stars, standard (fold\_constants\_com\_sound)}




    Complete the \coqdocvar{while} case of the following proof. 
\begin{coqdoccode}
\coqdocemptyline
\coqdocnoindent
\coqdockw{Theorem} \coqdef{SECF.Equiv.fold constants com sound}{fold\_constants\_com\_sound}{\coqdoclemma{fold\_constants\_com\_sound}} :\coqdoceol
\coqdocindent{1.00em}
\coqref{SECF.Equiv.ctrans sound}{\coqdocdefinition{ctrans\_sound}} \coqref{SECF.Equiv.fold constants com}{\coqdocdefinition{fold\_constants\_com}}.\coqdoceol
\coqdocnoindent
\coqdockw{Proof}.\coqdoceol
\coqdocindent{1.00em}
\coqdoctac{unfold} \coqref{SECF.Equiv.ctrans sound}{\coqdocdefinition{ctrans\_sound}}. \coqdoctac{intros} \coqdocvar{c}.\coqdoceol
\coqdocindent{1.00em}
\coqdoctac{induction} \coqdocvar{c}; \coqdoctac{simpl}.\coqdoceol
\coqdocindent{1.00em}
-  \coqdoctac{apply} \coqref{SECF.Equiv.refl cequiv}{\coqdoclemma{refl\_cequiv}}.\coqdoceol
\coqdocindent{1.00em}
-  \coqdoctac{apply} \coqref{SECF.Equiv.CAsgn congruence}{\coqdoclemma{CAsgn\_congruence}}.\coqdoceol
\coqdocindent{7.00em}
\coqdoctac{apply} \coqref{SECF.Equiv.fold constants aexp sound}{\coqdoclemma{fold\_constants\_aexp\_sound}}.\coqdoceol
\coqdocindent{1.00em}
-  \coqdoctac{apply} \coqref{SECF.Equiv.CSeq congruence}{\coqdocaxiom{CSeq\_congruence}}; \coqdoctac{assumption}.\coqdoceol
\coqdocindent{1.00em}
- \coqdoceol
\coqdocindent{2.00em}
\coqdoctac{assert} (\coqref{SECF.Equiv.bequiv}{\coqdocdefinition{bequiv}} \coqdocvar{b} (\coqref{SECF.Equiv.fold constants bexp}{\coqdocdefinition{fold\_constants\_bexp}} \coqdocvar{b})). \{\coqdoceol
\coqdocindent{3.00em}
\coqdoctac{apply} \coqref{SECF.Equiv.fold constants bexp sound}{\coqdocaxiom{fold\_constants\_bexp\_sound}}. \}\coqdoceol
\coqdocindent{2.00em}
\coqdoctac{destruct} (\coqref{SECF.Equiv.fold constants bexp}{\coqdocdefinition{fold\_constants\_bexp}} \coqdocvar{b}) \coqdocvar{eqn}:\coqdocvar{Heqb};\coqdoceol
\coqdocindent{3.00em}
\coqdoctac{try} (\coqdoctac{apply} \coqref{SECF.Equiv.CIf congruence}{\coqdocaxiom{CIf\_congruence}}; \coqdoctac{assumption}).\coqdoceol
\coqdocindent{2.00em}
+ \coqdoceol
\coqdocindent{3.00em}
\coqdoctac{apply} \coqref{SECF.Equiv.trans cequiv}{\coqdoclemma{trans\_cequiv}} \coqdockw{with} \coqdocvar{c1}; \coqdoctac{try} \coqdoctac{assumption}.\coqdoceol
\coqdocindent{3.00em}
\coqdoctac{apply} \coqref{SECF.Equiv.if true}{\coqdoclemma{if\_true}}; \coqdoctac{assumption}.\coqdoceol
\coqdocindent{2.00em}
+ \coqdoceol
\coqdocindent{3.00em}
\coqdoctac{apply} \coqref{SECF.Equiv.trans cequiv}{\coqdoclemma{trans\_cequiv}} \coqdockw{with} \coqdocvar{c2}; \coqdoctac{try} \coqdoctac{assumption}.\coqdoceol
\coqdocindent{3.00em}
\coqdoctac{apply} \coqref{SECF.Equiv.if false}{\coqdocaxiom{if\_false}}; \coqdoctac{assumption}.\coqdoceol
\coqdocindent{1.00em}
- \coqdoceol
\coqdocindent{2.00em}
 \coqdocvar{Admitted}.\coqdoceol
\end{coqdoccode}
\ensuremath{\Box} 
\begin{coqdoccode}
\coqdocemptyline
\end{coqdoccode}
\subsection{Soundness of (0 + n) Elimination, Redux}



\paragraph{Exercise: 4 stars, standard, optional (optimize\_0plus\_var)}




    Recall the definition \coqref{SECF.Imp.AExp.optimize 0plus}{\coqdocdefinition{optimize\_0plus}} from the \coqdoclibrary{Imp} chapter
    of \textit{Logical Foundations}:


    Fixpoint optimize\_0plus (a:aexp) : aexp :=
      match a with
      | ANum n =>
          ANum n
      | <\{ 0 + a2 \}> =>
          optimize\_0plus a2
      | <\{ a1 + a2 \}> =>
          <\{ (optimize\_0plus a1) + (optimize\_0plus a2) \}>
      | <\{ a1 - a2 \}> =>
          <\{ (optimize\_0plus a1) - (optimize\_0plus a2) \}>
      | <\{ a1 * a2 \}> =>
          <\{ (optimize\_0plus a1) * (optimize\_0plus a2) \}>
      end.


   Note that this function is defined over the old version of \coqref{SECF.Imp.aexp}{\coqdocinductive{aexp}}s,
   without states.


   Write a new version of this function that deals with variables (by
   leaving them alone), plus analogous ones for \coqref{SECF.Imp.bexp}{\coqdocinductive{bexp}}s and commands:


     optimize\_0plus\_aexp
     optimize\_0plus\_bexp
     optimize\_0plus\_com

\begin{coqdoccode}
\coqdocemptyline
\coqdocnoindent
\coqdockw{Fixpoint} \coqdef{SECF.Equiv.optimize 0plus aexp}{optimize\_0plus\_aexp}{\coqdocdefinition{optimize\_0plus\_aexp}} (\coqdef{SECF.Equiv.a:102}{a}{\coqdocbinder{a}} : \coqref{SECF.Imp.aexp}{\coqdocinductive{aexp}}) : \coqref{SECF.Imp.aexp}{\coqdocinductive{aexp}}\coqdoceol
\coqdocindent{1.00em}
. \coqdocvar{Admitted}.\coqdoceol
\coqdocemptyline
\coqdocnoindent
\coqdockw{Fixpoint} \coqdef{SECF.Equiv.optimize 0plus bexp}{optimize\_0plus\_bexp}{\coqdocdefinition{optimize\_0plus\_bexp}} (\coqdef{SECF.Equiv.b:104}{b}{\coqdocbinder{b}} : \coqref{SECF.Imp.bexp}{\coqdocinductive{bexp}}) : \coqref{SECF.Imp.bexp}{\coqdocinductive{bexp}}\coqdoceol
\coqdocindent{1.00em}
. \coqdocvar{Admitted}.\coqdoceol
\coqdocemptyline
\coqdocnoindent
\coqdockw{Fixpoint} \coqdef{SECF.Equiv.optimize 0plus com}{optimize\_0plus\_com}{\coqdocdefinition{optimize\_0plus\_com}} (\coqdef{SECF.Equiv.c:106}{c}{\coqdocbinder{c}} : \coqref{SECF.Imp.com}{\coqdocinductive{com}}) : \coqref{SECF.Imp.com}{\coqdocinductive{com}}\coqdoceol
\coqdocindent{1.00em}
. \coqdocvar{Admitted}.\coqdoceol
\coqdocemptyline
\coqdocnoindent
\coqdockw{Example} \coqdef{SECF.Equiv.test optimize 0plus}{test\_optimize\_0plus}{\coqdocdefinition{test\_optimize\_0plus}}:\coqdoceol
\coqdocindent{2.00em}
\coqref{SECF.Equiv.optimize 0plus com}{\coqdocaxiom{optimize\_0plus\_com}}\coqdoceol
\coqdocindent{3.50em}
\coqref{SECF.Imp.::com scope:'<x7B' x 'x7D>'}{\coqdocnotation{<\{}} \coqref{SECF.Imp.:com:com scope:'while' x 'do' x 'end'}{\coqdocnotation{while}} \coqref{SECF.Imp.X}{\coqdocdefinition{X}} \coqref{SECF.Imp.:com::x '<>' x}{\coqdocnotation{\ensuremath{\not=}}} 0 \coqref{SECF.Imp.:com:com scope:'while' x 'do' x 'end'}{\coqdocnotation{do}} \coqref{SECF.Imp.X}{\coqdocdefinition{X}} \coqref{SECF.Imp.:com:com scope:x ':=' x}{\coqdocnotation{:=}} 0 \coqref{SECF.Imp.:com::x '+' x}{\coqdocnotation{+}} \coqref{SECF.Imp.X}{\coqdocdefinition{X}} \coqref{SECF.Imp.:com::x '-' x}{\coqdocnotation{-}} 1 \coqref{SECF.Imp.:com:com scope:'while' x 'do' x 'end'}{\coqdocnotation{end}} \coqref{SECF.Imp.::com scope:'<x7B' x 'x7D>'}{\coqdocnotation{\}>}}\coqdoceol
\coqdocindent{1.00em}
\coqexternalref{::type scope:x '=' x}{http://rocq-prover.org/doc/V9.1.1/stdlib//Corelib.Init.Logic}{\coqdocnotation{=}}    \coqref{SECF.Imp.::com scope:'<x7B' x 'x7D>'}{\coqdocnotation{<\{}} \coqref{SECF.Imp.:com:com scope:'while' x 'do' x 'end'}{\coqdocnotation{while}} \coqref{SECF.Imp.X}{\coqdocdefinition{X}} \coqref{SECF.Imp.:com::x '<>' x}{\coqdocnotation{\ensuremath{\not=}}} 0 \coqref{SECF.Imp.:com:com scope:'while' x 'do' x 'end'}{\coqdocnotation{do}} \coqref{SECF.Imp.X}{\coqdocdefinition{X}} \coqref{SECF.Imp.:com:com scope:x ':=' x}{\coqdocnotation{:=}} \coqref{SECF.Imp.X}{\coqdocdefinition{X}} \coqref{SECF.Imp.:com::x '-' x}{\coqdocnotation{-}} 1 \coqref{SECF.Imp.:com:com scope:'while' x 'do' x 'end'}{\coqdocnotation{end}} \coqref{SECF.Imp.::com scope:'<x7B' x 'x7D>'}{\coqdocnotation{\}>}}.\coqdoceol
\coqdocnoindent
\coqdockw{Proof}.\coqdoceol
 \coqdocvar{Admitted}.\coqdoceol
\coqdocemptyline
\end{coqdoccode}
Prove that these three functions are sound, as we did for
    \coqdocvar{fold\_constants\_}\ensuremath{\times}.  Make sure you use the congruence lemmas in the
    proof of \coqref{SECF.Equiv.optimize 0plus com}{\coqdocdefinition{optimize\_0plus\_com}} -- otherwise it will be \textit{long}! 
\begin{coqdoccode}
\coqdocemptyline
\coqdocnoindent
\coqdockw{Theorem} \coqdef{SECF.Equiv.optimize 0plus aexp sound}{optimize\_0plus\_aexp\_sound}{\coqdoclemma{optimize\_0plus\_aexp\_sound}}:\coqdoceol
\coqdocindent{1.00em}
\coqref{SECF.Equiv.atrans sound}{\coqdocdefinition{atrans\_sound}} \coqref{SECF.Equiv.optimize 0plus aexp}{\coqdocaxiom{optimize\_0plus\_aexp}}.\coqdoceol
\coqdocnoindent
\coqdockw{Proof}.\coqdoceol
 \coqdocvar{Admitted}.\coqdoceol
\coqdocemptyline
\coqdocnoindent
\coqdockw{Theorem} \coqdef{SECF.Equiv.optimize 0plus bexp sound}{optimize\_0plus\_bexp\_sound}{\coqdoclemma{optimize\_0plus\_bexp\_sound}} :\coqdoceol
\coqdocindent{1.00em}
\coqref{SECF.Equiv.btrans sound}{\coqdocdefinition{btrans\_sound}} \coqref{SECF.Equiv.optimize 0plus bexp}{\coqdocaxiom{optimize\_0plus\_bexp}}.\coqdoceol
\coqdocnoindent
\coqdockw{Proof}.\coqdoceol
 \coqdocvar{Admitted}.\coqdoceol
\coqdocemptyline
\coqdocnoindent
\coqdockw{Theorem} \coqdef{SECF.Equiv.optimize 0plus com sound}{optimize\_0plus\_com\_sound}{\coqdoclemma{optimize\_0plus\_com\_sound}} :\coqdoceol
\coqdocindent{1.00em}
\coqref{SECF.Equiv.ctrans sound}{\coqdocdefinition{ctrans\_sound}} \coqref{SECF.Equiv.optimize 0plus com}{\coqdocaxiom{optimize\_0plus\_com}}.\coqdoceol
\coqdocnoindent
\coqdockw{Proof}.\coqdoceol
 \coqdocvar{Admitted}.\coqdoceol
\coqdocemptyline
\end{coqdoccode}
Finally, let's define a compound optimizer on commands that first
    folds constants (using \coqref{SECF.Equiv.fold constants com}{\coqdocdefinition{fold\_constants\_com}}) and then eliminates
    0 + \coqdocvar{n} terms (using \coqref{SECF.Equiv.optimize 0plus com}{\coqdocdefinition{optimize\_0plus\_com}}). 
\begin{coqdoccode}
\coqdocemptyline
\coqdocnoindent
\coqdockw{Definition} \coqdef{SECF.Equiv.optimizer}{optimizer}{\coqdocdefinition{optimizer}} (\coqdef{SECF.Equiv.c:108}{c}{\coqdocbinder{c}} : \coqref{SECF.Imp.com}{\coqdocinductive{com}}) := \coqref{SECF.Equiv.optimize 0plus com}{\coqdocaxiom{optimize\_0plus\_com}} (\coqref{SECF.Equiv.fold constants com}{\coqdocdefinition{fold\_constants\_com}} \coqref{SECF.Equiv.c:108}{\coqdocvariable{c}}).\coqdoceol
\coqdocemptyline
\end{coqdoccode}
Prove that this optimizer is sound. 
\begin{coqdoccode}
\coqdocemptyline
\coqdocnoindent
\coqdockw{Theorem} \coqdef{SECF.Equiv.optimizer sound}{optimizer\_sound}{\coqdoclemma{optimizer\_sound}} :\coqdoceol
\coqdocindent{1.00em}
\coqref{SECF.Equiv.ctrans sound}{\coqdocdefinition{ctrans\_sound}} \coqref{SECF.Equiv.optimizer}{\coqdocdefinition{optimizer}}.\coqdoceol
\coqdocnoindent
\coqdockw{Proof}.\coqdoceol
 \coqdocvar{Admitted}.\coqdoceol
\end{coqdoccode}
\ensuremath{\Box} 
\begin{coqdoccode}
\coqdocemptyline
\end{coqdoccode}
\section{Proving Inequivalence}



 Next, let's look at some programs that are \textit{not} equivalent. 

 Suppose that \coqdocvar{c1} is a command of the form


       X := a1; Y := a2


    and \coqdocvar{c2} is the command


       X := a1; Y := a2'


    where \coqdocvar{a2'} is formed by substituting \coqdocvar{a1} for all occurrences
    of \coqref{SECF.SpecCT.X}{\coqdocdefinition{X}} in \coqdocvar{a2}.


    For example, \coqdocvar{c1} and \coqdocvar{c2} might be:


       c1  =  (X := 42 + 53;
               Y := Y + X)
       c2  =  (X := 42 + 53;
               Y := Y + (42 + 53))


    Clearly, this \textit{particular} \coqdocvar{c1} and \coqdocvar{c2} are equivalent.  Is this
    true in general? 

 We will see in a moment that it is not, but it is worthwhile
    to pause, now, and see if you can find a counter-example on your
    own. 

 More formally, here is the function that substitutes an arithmetic
    expression \coqdocvar{u} for each occurrence of a given variable \coqdocvar{x} in
    another expression \coqdocvar{a}: 
\begin{coqdoccode}
\coqdocemptyline
\coqdocnoindent
\coqdockw{Fixpoint} \coqdef{SECF.Equiv.subst aexp}{subst\_aexp}{\coqdocdefinition{subst\_aexp}} (\coqdef{SECF.Equiv.x:109}{x}{\coqdocbinder{x}} : \coqdocinductive{string}) (\coqdef{SECF.Equiv.u:110}{u}{\coqdocbinder{u}} : \coqref{SECF.Imp.aexp}{\coqdocinductive{aexp}}) (\coqdef{SECF.Equiv.a:111}{a}{\coqdocbinder{a}} : \coqref{SECF.Imp.aexp}{\coqdocinductive{aexp}}) : \coqref{SECF.Imp.aexp}{\coqdocinductive{aexp}} :=\coqdoceol
\coqdocindent{1.00em}
\coqdockw{match} \coqref{SECF.Equiv.a:111}{\coqdocvariable{a}} \coqdockw{with}\coqdoceol
\coqdocindent{1.00em}
\ensuremath{|} \coqref{SECF.Imp.ANum}{\coqdocconstructor{ANum}} \coqdocvar{n}       \ensuremath{\Rightarrow}\coqdoceol
\coqdocindent{3.00em}
\coqref{SECF.Imp.ANum}{\coqdocconstructor{ANum}} \coqdocvar{n}\coqdoceol
\coqdocindent{1.00em}
\ensuremath{|} \coqref{SECF.Imp.AId}{\coqdocconstructor{AId}} \coqdocvar{x'}       \ensuremath{\Rightarrow}\coqdoceol
\coqdocindent{3.00em}
\coqdockw{if} \coqdocdefinition{String.eqb} \coqref{SECF.Equiv.x:109}{\coqdocvariable{x}} \coqdocvar{x'} \coqdockw{then} \coqref{SECF.Equiv.u:110}{\coqdocvariable{u}} \coqdockw{else} \coqref{SECF.Imp.AId}{\coqdocconstructor{AId}} \coqdocvar{x'}\coqdoceol
\coqdocindent{1.00em}
\ensuremath{|} \coqref{SECF.Imp.::com scope:'<x7B' x 'x7D>'}{\coqdocnotation{<\{}} \coqdocvar{a1} \coqref{SECF.Imp.:com::x '+' x}{\coqdocnotation{+}} \coqdocvar{a2} \coqref{SECF.Imp.::com scope:'<x7B' x 'x7D>'}{\coqdocnotation{\}>}}  \ensuremath{\Rightarrow}\coqdoceol
\coqdocindent{3.00em}
\coqref{SECF.Imp.::com scope:'<x7B' x 'x7D>'}{\coqdocnotation{<\{}} \coqref{SECF.Imp.:com::'(' x ')'}{\coqdocnotation{(}}\coqref{SECF.Equiv.subst aexp:112}{\coqdocdefinition{subst\_aexp}} \coqref{SECF.Equiv.x:109}{\coqdocvariable{x}} \coqref{SECF.Equiv.u:110}{\coqdocvariable{u}} \coqdocvar{a1}\coqref{SECF.Imp.:com::'(' x ')'}{\coqdocnotation{)}} \coqref{SECF.Imp.:com::x '+' x}{\coqdocnotation{+}} \coqref{SECF.Imp.:com::'(' x ')'}{\coqdocnotation{(}}\coqref{SECF.Equiv.subst aexp:112}{\coqdocdefinition{subst\_aexp}} \coqref{SECF.Equiv.x:109}{\coqdocvariable{x}} \coqref{SECF.Equiv.u:110}{\coqdocvariable{u}} \coqdocvar{a2}\coqref{SECF.Imp.:com::'(' x ')'}{\coqdocnotation{)}} \coqref{SECF.Imp.::com scope:'<x7B' x 'x7D>'}{\coqdocnotation{\}>}}\coqdoceol
\coqdocindent{1.00em}
\ensuremath{|} \coqref{SECF.Imp.::com scope:'<x7B' x 'x7D>'}{\coqdocnotation{<\{}} \coqdocvar{a1} \coqref{SECF.Imp.:com::x '-' x}{\coqdocnotation{-}} \coqdocvar{a2} \coqref{SECF.Imp.::com scope:'<x7B' x 'x7D>'}{\coqdocnotation{\}>}} \ensuremath{\Rightarrow}\coqdoceol
\coqdocindent{3.00em}
\coqref{SECF.Imp.::com scope:'<x7B' x 'x7D>'}{\coqdocnotation{<\{}} \coqref{SECF.Imp.:com::'(' x ')'}{\coqdocnotation{(}}\coqref{SECF.Equiv.subst aexp:112}{\coqdocdefinition{subst\_aexp}} \coqref{SECF.Equiv.x:109}{\coqdocvariable{x}} \coqref{SECF.Equiv.u:110}{\coqdocvariable{u}} \coqdocvar{a1}\coqref{SECF.Imp.:com::'(' x ')'}{\coqdocnotation{)}} \coqref{SECF.Imp.:com::x '-' x}{\coqdocnotation{-}} \coqref{SECF.Imp.:com::'(' x ')'}{\coqdocnotation{(}}\coqref{SECF.Equiv.subst aexp:112}{\coqdocdefinition{subst\_aexp}} \coqref{SECF.Equiv.x:109}{\coqdocvariable{x}} \coqref{SECF.Equiv.u:110}{\coqdocvariable{u}} \coqdocvar{a2}\coqref{SECF.Imp.:com::'(' x ')'}{\coqdocnotation{)}} \coqref{SECF.Imp.::com scope:'<x7B' x 'x7D>'}{\coqdocnotation{\}>}}\coqdoceol
\coqdocindent{1.00em}
\ensuremath{|} \coqref{SECF.Imp.::com scope:'<x7B' x 'x7D>'}{\coqdocnotation{<\{}} \coqdocvar{a1} \coqref{SECF.Imp.:com::x '*' x}{\coqdocnotation{\ensuremath{\times}}} \coqdocvar{a2} \coqref{SECF.Imp.::com scope:'<x7B' x 'x7D>'}{\coqdocnotation{\}>}}  \ensuremath{\Rightarrow}\coqdoceol
\coqdocindent{3.00em}
\coqref{SECF.Imp.::com scope:'<x7B' x 'x7D>'}{\coqdocnotation{<\{}} \coqref{SECF.Imp.:com::'(' x ')'}{\coqdocnotation{(}}\coqref{SECF.Equiv.subst aexp:112}{\coqdocdefinition{subst\_aexp}} \coqref{SECF.Equiv.x:109}{\coqdocvariable{x}} \coqref{SECF.Equiv.u:110}{\coqdocvariable{u}} \coqdocvar{a1}\coqref{SECF.Imp.:com::'(' x ')'}{\coqdocnotation{)}} \coqref{SECF.Imp.:com::x '*' x}{\coqdocnotation{\ensuremath{\times}}} \coqref{SECF.Imp.:com::'(' x ')'}{\coqdocnotation{(}}\coqref{SECF.Equiv.subst aexp:112}{\coqdocdefinition{subst\_aexp}} \coqref{SECF.Equiv.x:109}{\coqdocvariable{x}} \coqref{SECF.Equiv.u:110}{\coqdocvariable{u}} \coqdocvar{a2}\coqref{SECF.Imp.:com::'(' x ')'}{\coqdocnotation{)}} \coqref{SECF.Imp.::com scope:'<x7B' x 'x7D>'}{\coqdocnotation{\}>}}\coqdoceol
\coqdocindent{1.00em}
\coqdockw{end}.\coqdoceol
\coqdocemptyline
\coqdocnoindent
\coqdockw{Example} \coqdef{SECF.Equiv.subst aexp ex}{subst\_aexp\_ex}{\coqdocdefinition{subst\_aexp\_ex}} :\coqdoceol
\coqdocindent{1.00em}
\coqref{SECF.Equiv.subst aexp}{\coqdocdefinition{subst\_aexp}} \coqref{SECF.Imp.X}{\coqdocdefinition{X}} \coqref{SECF.Imp.::com scope:'<x7B' x 'x7D>'}{\coqdocnotation{<\{}}42 \coqref{SECF.Imp.:com::x '+' x}{\coqdocnotation{+}} 53\coqref{SECF.Imp.::com scope:'<x7B' x 'x7D>'}{\coqdocnotation{\}>}} \coqref{SECF.Imp.::com scope:'<x7B' x 'x7D>'}{\coqdocnotation{<\{}}\coqref{SECF.Imp.Y}{\coqdocdefinition{Y}} \coqref{SECF.Imp.:com::x '+' x}{\coqdocnotation{+}} \coqref{SECF.Imp.X}{\coqdocdefinition{X}}\coqref{SECF.Imp.::com scope:'<x7B' x 'x7D>'}{\coqdocnotation{\}>}}\coqdoceol
\coqdocindent{1.00em}
\coqexternalref{::type scope:x '=' x}{http://rocq-prover.org/doc/V9.1.1/stdlib//Corelib.Init.Logic}{\coqdocnotation{=}} \coqref{SECF.Imp.::com scope:'<x7B' x 'x7D>'}{\coqdocnotation{<\{}} \coqref{SECF.Imp.Y}{\coqdocdefinition{Y}} \coqref{SECF.Imp.:com::x '+' x}{\coqdocnotation{+}} \coqref{SECF.Imp.:com::'(' x ')'}{\coqdocnotation{(}}42 \coqref{SECF.Imp.:com::x '+' x}{\coqdocnotation{+}} 53\coqref{SECF.Imp.:com::'(' x ')'}{\coqdocnotation{)}}\coqref{SECF.Imp.::com scope:'<x7B' x 'x7D>'}{\coqdocnotation{\}>}}.\coqdoceol
\coqdocnoindent
\coqdockw{Proof}. \coqdoctac{simpl}. \coqdoctac{reflexivity}. \coqdockw{Qed}.\coqdoceol
\coqdocemptyline
\end{coqdoccode}
And here is the property we are interested in, expressing the
    claim that commands \coqdocvar{c1} and \coqdocvar{c2} as described above are
    always equivalent.  
\begin{coqdoccode}
\coqdocemptyline
\coqdocnoindent
\coqdockw{Definition} \coqdef{SECF.Equiv.subst equiv property}{subst\_equiv\_property}{\coqdocdefinition{subst\_equiv\_property}} : \coqdockw{Prop} := \coqdockw{\ensuremath{\forall}} \coqdef{SECF.Equiv.x1:114}{x1}{\coqdocbinder{x1}} \coqdef{SECF.Equiv.x2:115}{x2}{\coqdocbinder{x2}} \coqdef{SECF.Equiv.a1:116}{a1}{\coqdocbinder{a1}} \coqdef{SECF.Equiv.a2:117}{a2}{\coqdocbinder{a2}},\coqdoceol
\coqdocindent{1.00em}
\coqref{SECF.Equiv.cequiv}{\coqdocdefinition{cequiv}} \coqref{SECF.Imp.::com scope:'<x7B' x 'x7D>'}{\coqdocnotation{<\{}} \coqref{SECF.Equiv.x1:114}{\coqdocvariable{x1}} \coqref{SECF.Imp.:com:com scope:x ':=' x}{\coqdocnotation{:=}} \coqref{SECF.Equiv.a1:116}{\coqdocvariable{a1}}\coqref{SECF.Imp.:com:com scope:x ';' x}{\coqdocnotation{;}} \coqref{SECF.Equiv.x2:115}{\coqdocvariable{x2}} \coqref{SECF.Imp.:com:com scope:x ':=' x}{\coqdocnotation{:=}} \coqref{SECF.Equiv.a2:117}{\coqdocvariable{a2}} \coqref{SECF.Imp.::com scope:'<x7B' x 'x7D>'}{\coqdocnotation{\}>}}\coqdoceol
\coqdocindent{4.50em}
\coqref{SECF.Imp.::com scope:'<x7B' x 'x7D>'}{\coqdocnotation{<\{}} \coqref{SECF.Equiv.x1:114}{\coqdocvariable{x1}} \coqref{SECF.Imp.:com:com scope:x ':=' x}{\coqdocnotation{:=}} \coqref{SECF.Equiv.a1:116}{\coqdocvariable{a1}}\coqref{SECF.Imp.:com:com scope:x ';' x}{\coqdocnotation{;}} \coqref{SECF.Equiv.x2:115}{\coqdocvariable{x2}} \coqref{SECF.Imp.:com:com scope:x ':=' x}{\coqdocnotation{:=}} \coqref{SECF.Equiv.subst aexp}{\coqdocdefinition{subst\_aexp}} \coqref{SECF.Equiv.x1:114}{\coqdocvariable{x1}} \coqref{SECF.Equiv.a1:116}{\coqdocvariable{a1}} \coqref{SECF.Equiv.a2:117}{\coqdocvariable{a2}} \coqref{SECF.Imp.::com scope:'<x7B' x 'x7D>'}{\coqdocnotation{\}>}}.\coqdoceol
\coqdocemptyline
\end{coqdoccode}
Sadly, the property does \textit{not} always hold.


    Here is a counterexample:


       X := X + 1; Y := X


    If we perform the substitution, we get


       X := X + 1; Y := X + 1


    which clearly isn't equivalent. 
\begin{coqdoccode}
\coqdocemptyline
\coqdocnoindent
\coqdockw{Theorem} \coqdef{SECF.Equiv.subst inequiv}{subst\_inequiv}{\coqdoclemma{subst\_inequiv}} :\coqdoceol
\coqdocindent{1.00em}
\coqexternalref{::type scope:'x7E' x}{http://rocq-prover.org/doc/V9.1.1/stdlib//Corelib.Init.Logic}{\coqdocnotation{\ensuremath{\lnot}}} \coqref{SECF.Equiv.subst equiv property}{\coqdocdefinition{subst\_equiv\_property}}.\coqdoceol
\coqdocnoindent
\coqdockw{Proof}.\coqdoceol
\coqdocindent{1.00em}
\coqdoctac{unfold} \coqref{SECF.Equiv.subst equiv property}{\coqdocdefinition{subst\_equiv\_property}}.\coqdoceol
\coqdocindent{1.00em}
\coqdoctac{intros} \coqdocvar{Contra}.\coqdoceol
\coqdocemptyline
\coqdocindent{1.00em}
\coqdocvar{remember} \coqref{SECF.Imp.::com scope:'<x7B' x 'x7D>'}{\coqdocnotation{<\{}} \coqref{SECF.Imp.X}{\coqdocdefinition{X}} \coqref{SECF.Imp.:com:com scope:x ':=' x}{\coqdocnotation{:=}} \coqref{SECF.Imp.X}{\coqdocdefinition{X}} \coqref{SECF.Imp.:com::x '+' x}{\coqdocnotation{+}} 1\coqref{SECF.Imp.:com:com scope:x ';' x}{\coqdocnotation{;}}\coqdoceol
\coqdocindent{7.00em}
\coqref{SECF.Imp.Y}{\coqdocdefinition{Y}} \coqref{SECF.Imp.:com:com scope:x ':=' x}{\coqdocnotation{:=}} \coqref{SECF.Imp.X}{\coqdocdefinition{X}} \coqref{SECF.Imp.::com scope:'<x7B' x 'x7D>'}{\coqdocnotation{\}>}}\coqdoceol
\coqdocindent{3.00em}
\coqdockw{as} \coqdocvar{c1}.\coqdoceol
\coqdocindent{1.00em}
\coqdocvar{remember} \coqref{SECF.Imp.::com scope:'<x7B' x 'x7D>'}{\coqdocnotation{<\{}} \coqref{SECF.Imp.X}{\coqdocdefinition{X}} \coqref{SECF.Imp.:com:com scope:x ':=' x}{\coqdocnotation{:=}} \coqref{SECF.Imp.X}{\coqdocdefinition{X}} \coqref{SECF.Imp.:com::x '+' x}{\coqdocnotation{+}} 1\coqref{SECF.Imp.:com:com scope:x ';' x}{\coqdocnotation{;}}\coqdoceol
\coqdocindent{7.00em}
\coqref{SECF.Imp.Y}{\coqdocdefinition{Y}} \coqref{SECF.Imp.:com:com scope:x ':=' x}{\coqdocnotation{:=}} \coqref{SECF.Imp.X}{\coqdocdefinition{X}} \coqref{SECF.Imp.:com::x '+' x}{\coqdocnotation{+}} 1 \coqref{SECF.Imp.::com scope:'<x7B' x 'x7D>'}{\coqdocnotation{\}>}}\coqdoceol
\coqdocindent{3.00em}
\coqdockw{as} \coqdocvar{c2}.\coqdoceol
\coqdocindent{1.00em}
\coqdoctac{assert} (\coqref{SECF.Equiv.cequiv}{\coqdocdefinition{cequiv}} \coqdocvar{c1} \coqdocvar{c2}) \coqdoctac{by} (\coqdoctac{subst}; \coqdoctac{apply} \coqdocvar{Contra}).\coqdoceol
\coqdocindent{1.00em}
\coqdoctac{clear} \coqdocvar{Contra}.\coqdoceol
\coqdocemptyline
\coqdocindent{1.00em}
\coqdocvar{remember} (\coqref{SECF.Imp.Y}{\coqdocdefinition{Y}} \coqref{SECF.Maps.:::x '!->' x ';' x}{\coqdocnotation{!->}} 1 \coqref{SECF.Maps.:::x '!->' x ';' x}{\coqdocnotation{;}} \coqref{SECF.Imp.X}{\coqdocdefinition{X}} \coqref{SECF.Imp.:::x '!->' x}{\coqdocnotation{!->}} 1) \coqdockw{as} \coqdocvar{st1}.\coqdoceol
\coqdocindent{1.00em}
\coqdocvar{remember} (\coqref{SECF.Imp.Y}{\coqdocdefinition{Y}} \coqref{SECF.Maps.:::x '!->' x ';' x}{\coqdocnotation{!->}} 2 \coqref{SECF.Maps.:::x '!->' x ';' x}{\coqdocnotation{;}} \coqref{SECF.Imp.X}{\coqdocdefinition{X}} \coqref{SECF.Imp.:::x '!->' x}{\coqdocnotation{!->}} 1) \coqdockw{as} \coqdocvar{st2}.\coqdoceol
\coqdocindent{1.00em}
\coqdoctac{assert} (\coqdocvar{H1} : \coqref{SECF.Imp.empty st}{\coqdocdefinition{empty\_st}} \coqref{SECF.Imp.:::x '=[' x ']=>' x}{\coqdocnotation{=[}} \coqdocvar{c1} \coqref{SECF.Imp.:::x '=[' x ']=>' x}{\coqdocnotation{]=>}} \coqdocvar{st1});\coqdoceol
\coqdocindent{1.00em}
\coqdoctac{assert} (\coqdocvar{H2} : \coqref{SECF.Imp.empty st}{\coqdocdefinition{empty\_st}} \coqref{SECF.Imp.:::x '=[' x ']=>' x}{\coqdocnotation{=[}} \coqdocvar{c2} \coqref{SECF.Imp.:::x '=[' x ']=>' x}{\coqdocnotation{]=>}} \coqdocvar{st2});\coqdoceol
\coqdocindent{1.00em}
\coqdoctac{try} (\coqdoctac{subst};\coqdoceol
\coqdocindent{3.50em}
\coqdoctac{apply} \coqref{SECF.Imp.E Seq}{\coqdocconstructor{E\_Seq}} \coqdockw{with} (\coqdocvar{st'} := (\coqref{SECF.Imp.X}{\coqdocdefinition{X}} \coqref{SECF.Imp.:::x '!->' x}{\coqdocnotation{!->}} 1));\coqdoceol
\coqdocindent{3.50em}
\coqdoctac{apply} \coqref{SECF.Imp.E Asgn}{\coqdocconstructor{E\_Asgn}}; \coqdoctac{reflexivity}).\coqdoceol
\coqdocindent{1.00em}
\coqdoctac{clear} \coqdocvar{Heqc1} \coqdocvar{Heqc2}.\coqdoceol
\coqdocemptyline
\coqdocindent{1.00em}
\coqdoctac{apply} \coqdocvar{H} \coqdoctac{in} \coqdocvar{H1}.\coqdoceol
\coqdocindent{1.00em}
\coqdoctac{clear} \coqdocvar{H}.\coqdoceol
\coqdocemptyline
\coqdocindent{1.00em}
\coqdoctac{assert} (\coqdocvar{Hcontra} : \coqdocvar{st1} \coqexternalref{::type scope:x '=' x}{http://rocq-prover.org/doc/V9.1.1/stdlib//Corelib.Init.Logic}{\coqdocnotation{=}} \coqdocvar{st2})\coqdoceol
\coqdocindent{2.00em}
\coqdoctac{by} (\coqdoctac{apply} (\coqref{SECF.Imp.ceval deterministic}{\coqdoclemma{ceval\_deterministic}} \coqdocvar{c2} \coqref{SECF.Imp.empty st}{\coqdocdefinition{empty\_st}}); \coqdoctac{assumption}).\coqdoceol
\coqdocindent{1.00em}
\coqdoctac{clear} \coqdocvar{H1} \coqdocvar{H2}.\coqdoceol
\coqdocemptyline
\coqdocindent{1.00em}
\coqdoctac{assert} (\coqdocvar{Hcontra'} : \coqdocvar{st1} \coqref{SECF.Imp.Y}{\coqdocdefinition{Y}} \coqexternalref{::type scope:x '=' x}{http://rocq-prover.org/doc/V9.1.1/stdlib//Corelib.Init.Logic}{\coqdocnotation{=}} \coqdocvar{st2} \coqref{SECF.Imp.Y}{\coqdocdefinition{Y}})\coqdoceol
\coqdocindent{2.00em}
\coqdoctac{by} (\coqdoctac{rewrite} \coqdocvar{Hcontra}; \coqdoctac{reflexivity}).\coqdoceol
\coqdocindent{1.00em}
\coqdoctac{subst}. \coqdoctac{discriminate}. \coqdockw{Qed}.\coqdoceol
\coqdocemptyline
\end{coqdoccode}
\paragraph{Exercise: 4 stars, standard, optional (better\_subst\_equiv)}




    The equivalence we had in mind above was not complete nonsense --
    in fact, it was actually almost right.  To make it correct, we
    just need to exclude the case where the variable \coqref{SECF.SpecCT.X}{\coqdocdefinition{X}} occurs in the
    right-hand side of the first assignment statement. 
\begin{coqdoccode}
\coqdocemptyline
\coqdocnoindent
\coqdockw{Inductive} \coqdef{SECF.Equiv.var not used in aexp}{var\_not\_used\_in\_aexp}{\coqdocinductive{var\_not\_used\_in\_aexp}} (\coqdef{SECF.Equiv.x:118}{x}{\coqdocbinder{x}} : \coqdocinductive{string}) : \coqref{SECF.Imp.aexp}{\coqdocinductive{aexp}} \coqexternalref{::type scope:x '->' x}{http://rocq-prover.org/doc/V9.1.1/stdlib//Corelib.Init.Logic}{\coqdocnotation{\ensuremath{\rightarrow}}} \coqdockw{Prop} :=\coqdoceol
\coqdocindent{1.00em}
\ensuremath{|} \coqdef{SECF.Equiv.VNUNum}{VNUNum}{\coqdocconstructor{VNUNum}} : \coqdockw{\ensuremath{\forall}} \coqdef{SECF.Equiv.n:121}{n}{\coqdocbinder{n}}, \coqref{SECF.Equiv.var not used in aexp:119}{\coqdocinductive{var\_not\_used\_in\_aexp}} \coqref{SECF.Equiv.x:118}{\coqdocvariable{x}} (\coqref{SECF.Imp.ANum}{\coqdocconstructor{ANum}} \coqref{SECF.Equiv.n:121}{\coqdocvariable{n}})\coqdoceol
\coqdocindent{1.00em}
\ensuremath{|} \coqdef{SECF.Equiv.VNUId}{VNUId}{\coqdocconstructor{VNUId}} : \coqdockw{\ensuremath{\forall}} \coqdef{SECF.Equiv.y:122}{y}{\coqdocbinder{y}}, \coqref{SECF.Equiv.x:118}{\coqdocvariable{x}} \coqexternalref{::type scope:x '<>' x}{http://rocq-prover.org/doc/V9.1.1/stdlib//Corelib.Init.Logic}{\coqdocnotation{\ensuremath{\not=}}} \coqref{SECF.Equiv.y:122}{\coqdocvariable{y}} \coqexternalref{::type scope:x '->' x}{http://rocq-prover.org/doc/V9.1.1/stdlib//Corelib.Init.Logic}{\coqdocnotation{\ensuremath{\rightarrow}}} \coqref{SECF.Equiv.var not used in aexp:119}{\coqdocinductive{var\_not\_used\_in\_aexp}} \coqref{SECF.Equiv.x:118}{\coqdocvariable{x}} (\coqref{SECF.Imp.AId}{\coqdocconstructor{AId}} \coqref{SECF.Equiv.y:122}{\coqdocvariable{y}})\coqdoceol
\coqdocindent{1.00em}
\ensuremath{|} \coqdef{SECF.Equiv.VNUPlus}{VNUPlus}{\coqdocconstructor{VNUPlus}} : \coqdockw{\ensuremath{\forall}} \coqdef{SECF.Equiv.a1:123}{a1}{\coqdocbinder{a1}} \coqdef{SECF.Equiv.a2:124}{a2}{\coqdocbinder{a2}},\coqdoceol
\coqdocindent{3.00em}
\coqref{SECF.Equiv.var not used in aexp:119}{\coqdocinductive{var\_not\_used\_in\_aexp}} \coqref{SECF.Equiv.x:118}{\coqdocvariable{x}} \coqref{SECF.Equiv.a1:123}{\coqdocvariable{a1}} \coqexternalref{::type scope:x '->' x}{http://rocq-prover.org/doc/V9.1.1/stdlib//Corelib.Init.Logic}{\coqdocnotation{\ensuremath{\rightarrow}}}\coqdoceol
\coqdocindent{3.00em}
\coqref{SECF.Equiv.var not used in aexp:119}{\coqdocinductive{var\_not\_used\_in\_aexp}} \coqref{SECF.Equiv.x:118}{\coqdocvariable{x}} \coqref{SECF.Equiv.a2:124}{\coqdocvariable{a2}} \coqexternalref{::type scope:x '->' x}{http://rocq-prover.org/doc/V9.1.1/stdlib//Corelib.Init.Logic}{\coqdocnotation{\ensuremath{\rightarrow}}}\coqdoceol
\coqdocindent{3.00em}
\coqref{SECF.Equiv.var not used in aexp:119}{\coqdocinductive{var\_not\_used\_in\_aexp}} \coqref{SECF.Equiv.x:118}{\coqdocvariable{x}} (\coqref{SECF.Imp.::com scope:'<x7B' x 'x7D>'}{\coqdocnotation{<\{}} \coqref{SECF.Equiv.a1:123}{\coqdocvariable{a1}} \coqref{SECF.Imp.:com::x '+' x}{\coqdocnotation{+}} \coqref{SECF.Equiv.a2:124}{\coqdocvariable{a2}} \coqref{SECF.Imp.::com scope:'<x7B' x 'x7D>'}{\coqdocnotation{\}>}})\coqdoceol
\coqdocindent{1.00em}
\ensuremath{|} \coqdef{SECF.Equiv.VNUMinus}{VNUMinus}{\coqdocconstructor{VNUMinus}} : \coqdockw{\ensuremath{\forall}} \coqdef{SECF.Equiv.a1:125}{a1}{\coqdocbinder{a1}} \coqdef{SECF.Equiv.a2:126}{a2}{\coqdocbinder{a2}},\coqdoceol
\coqdocindent{3.00em}
\coqref{SECF.Equiv.var not used in aexp:119}{\coqdocinductive{var\_not\_used\_in\_aexp}} \coqref{SECF.Equiv.x:118}{\coqdocvariable{x}} \coqref{SECF.Equiv.a1:125}{\coqdocvariable{a1}} \coqexternalref{::type scope:x '->' x}{http://rocq-prover.org/doc/V9.1.1/stdlib//Corelib.Init.Logic}{\coqdocnotation{\ensuremath{\rightarrow}}}\coqdoceol
\coqdocindent{3.00em}
\coqref{SECF.Equiv.var not used in aexp:119}{\coqdocinductive{var\_not\_used\_in\_aexp}} \coqref{SECF.Equiv.x:118}{\coqdocvariable{x}} \coqref{SECF.Equiv.a2:126}{\coqdocvariable{a2}} \coqexternalref{::type scope:x '->' x}{http://rocq-prover.org/doc/V9.1.1/stdlib//Corelib.Init.Logic}{\coqdocnotation{\ensuremath{\rightarrow}}}\coqdoceol
\coqdocindent{3.00em}
\coqref{SECF.Equiv.var not used in aexp:119}{\coqdocinductive{var\_not\_used\_in\_aexp}} \coqref{SECF.Equiv.x:118}{\coqdocvariable{x}} (\coqref{SECF.Imp.::com scope:'<x7B' x 'x7D>'}{\coqdocnotation{<\{}} \coqref{SECF.Equiv.a1:125}{\coqdocvariable{a1}} \coqref{SECF.Imp.:com::x '-' x}{\coqdocnotation{-}} \coqref{SECF.Equiv.a2:126}{\coqdocvariable{a2}} \coqref{SECF.Imp.::com scope:'<x7B' x 'x7D>'}{\coqdocnotation{\}>}})\coqdoceol
\coqdocindent{1.00em}
\ensuremath{|} \coqdef{SECF.Equiv.VNUMult}{VNUMult}{\coqdocconstructor{VNUMult}} : \coqdockw{\ensuremath{\forall}} \coqdef{SECF.Equiv.a1:127}{a1}{\coqdocbinder{a1}} \coqdef{SECF.Equiv.a2:128}{a2}{\coqdocbinder{a2}},\coqdoceol
\coqdocindent{3.00em}
\coqref{SECF.Equiv.var not used in aexp:119}{\coqdocinductive{var\_not\_used\_in\_aexp}} \coqref{SECF.Equiv.x:118}{\coqdocvariable{x}} \coqref{SECF.Equiv.a1:127}{\coqdocvariable{a1}} \coqexternalref{::type scope:x '->' x}{http://rocq-prover.org/doc/V9.1.1/stdlib//Corelib.Init.Logic}{\coqdocnotation{\ensuremath{\rightarrow}}}\coqdoceol
\coqdocindent{3.00em}
\coqref{SECF.Equiv.var not used in aexp:119}{\coqdocinductive{var\_not\_used\_in\_aexp}} \coqref{SECF.Equiv.x:118}{\coqdocvariable{x}} \coqref{SECF.Equiv.a2:128}{\coqdocvariable{a2}} \coqexternalref{::type scope:x '->' x}{http://rocq-prover.org/doc/V9.1.1/stdlib//Corelib.Init.Logic}{\coqdocnotation{\ensuremath{\rightarrow}}}\coqdoceol
\coqdocindent{3.00em}
\coqref{SECF.Equiv.var not used in aexp:119}{\coqdocinductive{var\_not\_used\_in\_aexp}} \coqref{SECF.Equiv.x:118}{\coqdocvariable{x}} (\coqref{SECF.Imp.::com scope:'<x7B' x 'x7D>'}{\coqdocnotation{<\{}} \coqref{SECF.Equiv.a1:127}{\coqdocvariable{a1}} \coqref{SECF.Imp.:com::x '*' x}{\coqdocnotation{\ensuremath{\times}}} \coqref{SECF.Equiv.a2:128}{\coqdocvariable{a2}} \coqref{SECF.Imp.::com scope:'<x7B' x 'x7D>'}{\coqdocnotation{\}>}}).\coqdoceol
\coqdocemptyline
\coqdocnoindent
\coqdockw{Lemma} \coqdef{SECF.Equiv.aeval weakening}{aeval\_weakening}{\coqdoclemma{aeval\_weakening}} : \coqdockw{\ensuremath{\forall}} \coqdef{SECF.Equiv.x:129}{x}{\coqdocbinder{x}} \coqdef{SECF.Equiv.st:130}{st}{\coqdocbinder{st}} \coqdef{SECF.Equiv.a:131}{a}{\coqdocbinder{a}} \coqdef{SECF.Equiv.ni:132}{ni}{\coqdocbinder{ni}},\coqdoceol
\coqdocindent{1.00em}
\coqref{SECF.Equiv.var not used in aexp}{\coqdocinductive{var\_not\_used\_in\_aexp}} \coqref{SECF.Equiv.x:129}{\coqdocvariable{x}} \coqref{SECF.Equiv.a:131}{\coqdocvariable{a}} \coqexternalref{::type scope:x '->' x}{http://rocq-prover.org/doc/V9.1.1/stdlib//Corelib.Init.Logic}{\coqdocnotation{\ensuremath{\rightarrow}}}\coqdoceol
\coqdocindent{1.00em}
\coqref{SECF.Imp.aeval}{\coqdocdefinition{aeval}} (\coqref{SECF.Equiv.x:129}{\coqdocvariable{x}} \coqref{SECF.Maps.:::x '!->' x ';' x}{\coqdocnotation{!->}} \coqref{SECF.Equiv.ni:132}{\coqdocvariable{ni}} \coqref{SECF.Maps.:::x '!->' x ';' x}{\coqdocnotation{;}} \coqref{SECF.Equiv.st:130}{\coqdocvariable{st}}) \coqref{SECF.Equiv.a:131}{\coqdocvariable{a}} \coqexternalref{::type scope:x '=' x}{http://rocq-prover.org/doc/V9.1.1/stdlib//Corelib.Init.Logic}{\coqdocnotation{=}} \coqref{SECF.Imp.aeval}{\coqdocdefinition{aeval}} \coqref{SECF.Equiv.st:130}{\coqdocvariable{st}} \coqref{SECF.Equiv.a:131}{\coqdocvariable{a}}.\coqdoceol
\coqdocnoindent
\coqdockw{Proof}.\coqdoceol
 \coqdocvar{Admitted}.\coqdoceol
\coqdocemptyline
\end{coqdoccode}
Using \coqref{SECF.Equiv.var not used in aexp}{\coqdocinductive{var\_not\_used\_in\_aexp}}, formalize and prove a correct version
    of \coqref{SECF.Equiv.subst equiv property}{\coqdocdefinition{subst\_equiv\_property}}. 
\begin{coqdoccode}
\coqdocemptyline
\coqdocemptyline
\end{coqdoccode}
\paragraph{Exercise: 3 stars, standard (inequiv\_exercise)}




    Prove that an infinite loop is not equivalent to \coqdocvar{skip} 
\begin{coqdoccode}
\coqdocemptyline
\coqdocnoindent
\coqdockw{Theorem} \coqdef{SECF.Equiv.inequiv exercise}{inequiv\_exercise}{\coqdoclemma{inequiv\_exercise}}:\coqdoceol
\coqdocindent{1.00em}
\coqexternalref{::type scope:'x7E' x}{http://rocq-prover.org/doc/V9.1.1/stdlib//Corelib.Init.Logic}{\coqdocnotation{\ensuremath{\lnot}}} \coqref{SECF.Equiv.cequiv}{\coqdocdefinition{cequiv}} \coqref{SECF.Imp.::com scope:'<x7B' x 'x7D>'}{\coqdocnotation{<\{}} \coqref{SECF.Imp.:com:com scope:'while' x 'do' x 'end'}{\coqdocnotation{while}} \coqref{SECF.Imp.:com::'true'}{\coqdocnotation{true}} \coqref{SECF.Imp.:com:com scope:'while' x 'do' x 'end'}{\coqdocnotation{do}} \coqref{SECF.Imp.:com:com scope:'skip'}{\coqdocnotation{skip}} \coqref{SECF.Imp.:com:com scope:'while' x 'do' x 'end'}{\coqdocnotation{end}} \coqref{SECF.Imp.::com scope:'<x7B' x 'x7D>'}{\coqdocnotation{\}>}} \coqref{SECF.Imp.::com scope:'<x7B' x 'x7D>'}{\coqdocnotation{<\{}} \coqref{SECF.Imp.:com:com scope:'skip'}{\coqdocnotation{skip}} \coqref{SECF.Imp.::com scope:'<x7B' x 'x7D>'}{\coqdocnotation{\}>}}.\coqdoceol
\coqdocnoindent
\coqdockw{Proof}.\coqdoceol
 \coqdocvar{Admitted}.\coqdoceol
\end{coqdoccode}
\ensuremath{\Box} 
\begin{coqdoccode}
\coqdocemptyline
\end{coqdoccode}
\section{Extended Exercise: Nondeterministic Imp}



 As we have seen (in theorem \coqref{SECF.Imp.BreakImp.ceval deterministic}{\coqdoclemma{ceval\_deterministic}} in the \coqdoclibrary{Imp}
    chapter), Imp's evaluation relation is deterministic.  However,
    \textit{non}-determinism is an important part of the definition of many
    real programming languages. For example, in many imperative
    languages (such as C and its relatives), the order in which
    function arguments are evaluated is unspecified: the program
    fragment


      x = 0;
      f(++x, x)


    might call \coqdocvar{f} with arguments (1, 0) or (1, 1), depending how
    the compiler chooses to order things.  This can be a little
    confusing for programmers, but it gives compiler writers useful
    freedom.


    In this exercise, we will extend Imp with a simple
    nondeterministic command and study how this change affects
    program equivalence.  The new command has the syntax \coqdocvar{HAVOC} \coqref{SECF.SpecCT.X}{\coqdocdefinition{X}},
    where \coqref{SECF.SpecCT.X}{\coqdocdefinition{X}} is an identifier. The effect of executing \coqdocvar{HAVOC} \coqref{SECF.SpecCT.X}{\coqdocdefinition{X}} is
    to assign an \textit{arbitrary} number to the variable \coqref{SECF.SpecCT.X}{\coqdocdefinition{X}},
    nondeterministically. For example, after executing the program:


      HAVOC Y;
      Z := Y * 2


    the value of \coqref{SECF.SpecCT.Y}{\coqdocdefinition{Y}} can be any number, while the value of \coqref{SECF.SpecCT.Z}{\coqdocdefinition{Z}} is
    twice that of \coqref{SECF.SpecCT.Y}{\coqdocdefinition{Y}} (so \coqref{SECF.SpecCT.Z}{\coqdocdefinition{Z}} is always even). Note that we are not
    saying anything about the \textit{probabilities} of the outcomes -- just
    that there are (infinitely) many different outcomes that can
    possibly happen after executing this nondeterministic code.


    In a sense, a variable on which we do \coqdocvar{HAVOC} roughly corresponds
    to an uninitialized variable in a low-level language like C.  After
    the \coqdocvar{HAVOC}, the variable holds a fixed but arbitrary number.  Most
    sources of nondeterminism in language definitions are there
    precisely because programmers don't care which choice is made (and
    so it is good to leave it open to the compiler to choose whichever
    will run faster).


    We call this new language \textit{Himp} (``Imp extended with \coqdocvar{HAVOC}''). 
\begin{coqdoccode}
\coqdocemptyline
\coqdocnoindent
\coqdockw{Module} \coqdef{SECF.Equiv.Himp}{Himp}{\coqdocmodule{Himp}}.\coqdoceol
\coqdocemptyline
\end{coqdoccode}
To formalize Himp, we first add a clause to the definition of
    commands. 
\begin{coqdoccode}
\coqdocemptyline
\coqdocnoindent
\coqdockw{Inductive} \coqdef{SECF.Equiv.Himp.com}{com}{\coqdocinductive{com}} : \coqdockw{Type} :=\coqdoceol
\coqdocindent{1.00em}
\ensuremath{|} \coqdef{SECF.Equiv.Himp.CSkip}{CSkip}{\coqdocconstructor{CSkip}} : \coqref{SECF.Equiv.com:133}{\coqdocinductive{com}}\coqdoceol
\coqdocindent{1.00em}
\ensuremath{|} \coqdef{SECF.Equiv.Himp.CAsgn}{CAsgn}{\coqdocconstructor{CAsgn}} : \coqdocinductive{string} \coqexternalref{::type scope:x '->' x}{http://rocq-prover.org/doc/V9.1.1/stdlib//Corelib.Init.Logic}{\coqdocnotation{\ensuremath{\rightarrow}}} \coqref{SECF.Imp.aexp}{\coqdocinductive{aexp}} \coqexternalref{::type scope:x '->' x}{http://rocq-prover.org/doc/V9.1.1/stdlib//Corelib.Init.Logic}{\coqdocnotation{\ensuremath{\rightarrow}}} \coqref{SECF.Equiv.com:133}{\coqdocinductive{com}}\coqdoceol
\coqdocindent{1.00em}
\ensuremath{|} \coqdef{SECF.Equiv.Himp.CSeq}{CSeq}{\coqdocconstructor{CSeq}} : \coqref{SECF.Equiv.com:133}{\coqdocinductive{com}} \coqexternalref{::type scope:x '->' x}{http://rocq-prover.org/doc/V9.1.1/stdlib//Corelib.Init.Logic}{\coqdocnotation{\ensuremath{\rightarrow}}} \coqref{SECF.Equiv.com:133}{\coqdocinductive{com}} \coqexternalref{::type scope:x '->' x}{http://rocq-prover.org/doc/V9.1.1/stdlib//Corelib.Init.Logic}{\coqdocnotation{\ensuremath{\rightarrow}}} \coqref{SECF.Equiv.com:133}{\coqdocinductive{com}}\coqdoceol
\coqdocindent{1.00em}
\ensuremath{|} \coqdef{SECF.Equiv.Himp.CIf}{CIf}{\coqdocconstructor{CIf}} : \coqref{SECF.Imp.bexp}{\coqdocinductive{bexp}} \coqexternalref{::type scope:x '->' x}{http://rocq-prover.org/doc/V9.1.1/stdlib//Corelib.Init.Logic}{\coqdocnotation{\ensuremath{\rightarrow}}} \coqref{SECF.Equiv.com:133}{\coqdocinductive{com}} \coqexternalref{::type scope:x '->' x}{http://rocq-prover.org/doc/V9.1.1/stdlib//Corelib.Init.Logic}{\coqdocnotation{\ensuremath{\rightarrow}}} \coqref{SECF.Equiv.com:133}{\coqdocinductive{com}} \coqexternalref{::type scope:x '->' x}{http://rocq-prover.org/doc/V9.1.1/stdlib//Corelib.Init.Logic}{\coqdocnotation{\ensuremath{\rightarrow}}} \coqref{SECF.Equiv.com:133}{\coqdocinductive{com}}\coqdoceol
\coqdocindent{1.00em}
\ensuremath{|} \coqdef{SECF.Equiv.Himp.CWhile}{CWhile}{\coqdocconstructor{CWhile}} : \coqref{SECF.Imp.bexp}{\coqdocinductive{bexp}} \coqexternalref{::type scope:x '->' x}{http://rocq-prover.org/doc/V9.1.1/stdlib//Corelib.Init.Logic}{\coqdocnotation{\ensuremath{\rightarrow}}} \coqref{SECF.Equiv.com:133}{\coqdocinductive{com}} \coqexternalref{::type scope:x '->' x}{http://rocq-prover.org/doc/V9.1.1/stdlib//Corelib.Init.Logic}{\coqdocnotation{\ensuremath{\rightarrow}}} \coqref{SECF.Equiv.com:133}{\coqdocinductive{com}}\coqdoceol
\coqdocindent{1.00em}
\ensuremath{|} \coqdef{SECF.Equiv.Himp.CHavoc}{CHavoc}{\coqdocconstructor{CHavoc}} : \coqdocinductive{string} \coqexternalref{::type scope:x '->' x}{http://rocq-prover.org/doc/V9.1.1/stdlib//Corelib.Init.Logic}{\coqdocnotation{\ensuremath{\rightarrow}}} \coqref{SECF.Equiv.com:133}{\coqdocinductive{com}}. \coqdocemptyline
\coqdocnoindent
\coqdockw{Notation} \coqdef{SECF.Equiv.Himp.:com::'havoc' x}{"}{"}'havoc' l" := (\coqref{SECF.Equiv.Himp.CHavoc}{\coqdocconstructor{CHavoc}} \coqdocvar{l})\coqdoceol
\coqdocindent{13.00em}
(\coqdoctac{in} \coqdocvar{custom} \coqdocvar{com} \coqdoctac{at} \coqdockw{level} 60, \coqdocvar{l} \coqdockw{constr} \coqdoctac{at} \coqdockw{level} 0).\coqdoceol
\coqdocnoindent
\coqdockw{Notation} \coqdef{SECF.Equiv.Himp.:com:com scope:'skip'}{"}{"}'skip'"  := \coqref{SECF.Equiv.Himp.CSkip}{\coqdocconstructor{CSkip}}\coqdoceol
\coqdocindent{1.00em}
(\coqdoctac{in} \coqdocvar{custom} \coqdocvar{com} \coqdoctac{at} \coqdockw{level} 0) : \coqdocvar{com\_scope}.\coqdoceol
\coqdocnoindent
\coqdockw{Notation} \coqdef{SECF.Equiv.Himp.:com:com scope:x ':=' x}{"}{"}x := y"  := (\coqref{SECF.Equiv.Himp.CAsgn}{\coqdocconstructor{CAsgn}} \coqdocvar{x} \coqdocvar{y})\coqdoceol
\coqdocindent{1.00em}
(\coqdoctac{in} \coqdocvar{custom} \coqdocvar{com} \coqdoctac{at} \coqdockw{level} 0, \coqdocvar{x} \coqdockw{constr} \coqdoctac{at} \coqdockw{level} 0, \coqdocvar{y} \coqdoctac{at} \coqdockw{level} 85, \coqdockw{no} \coqdockw{associativity},\coqdoceol
\coqdocindent{2.00em}
\coqdocvar{format} "x  :=  y") : \coqdocvar{com\_scope}.\coqdoceol
\coqdocnoindent
\coqdockw{Notation} \coqdef{SECF.Equiv.Himp.:com:com scope:x ';' x}{"}{"}x ; y" := (\coqref{SECF.Equiv.Himp.CSeq}{\coqdocconstructor{CSeq}} \coqdocvar{x} \coqdocvar{y})\coqdoceol
\coqdocindent{1.00em}
(\coqdoctac{in} \coqdocvar{custom} \coqdocvar{com} \coqdoctac{at} \coqdockw{level} 90,\coqdoceol
\coqdocindent{2.00em}
\coqdoctac{right} \coqdockw{associativity},\coqdoceol
\coqdocindent{2.00em}
\coqdocvar{format} "'[v' x ; '/' y ']'") : \coqdocvar{com\_scope}.\coqdoceol
\coqdocnoindent
\coqdockw{Notation} \coqdef{SECF.Equiv.Himp.:com:com scope:'if' x 'then' x 'else' x 'end'}{"}{"}'if' x 'then' y 'else' z 'end'" := (\coqref{SECF.Equiv.Himp.CIf}{\coqdocconstructor{CIf}} \coqdocvar{x} \coqdocvar{y} \coqdocvar{z})\coqdoceol
\coqdocindent{1.00em}
(\coqdoctac{in} \coqdocvar{custom} \coqdocvar{com} \coqdoctac{at} \coqdockw{level} 89, \coqdocvar{x} \coqdoctac{at} \coqdockw{level} 99, \coqdocvar{y} \coqdoctac{at} \coqdockw{level} 99, \coqdocvar{z} \coqdoctac{at} \coqdockw{level} 99,\coqdoceol
\coqdocindent{2.00em}
\coqdocvar{format} "'[v' 'if'  x  'then' '/  ' y '/' 'else' '/  ' z '/' 'end' ']'") : \coqdocvar{com\_scope}.\coqdoceol
\coqdocnoindent
\coqdockw{Notation} \coqdef{SECF.Equiv.Himp.:com:com scope:'while' x 'do' x 'end'}{"}{"}'while' x 'do' y 'end'" := (\coqref{SECF.Equiv.Himp.CWhile}{\coqdocconstructor{CWhile}} \coqdocvar{x} \coqdocvar{y})\coqdoceol
\coqdocindent{1.00em}
(\coqdoctac{in} \coqdocvar{custom} \coqdocvar{com} \coqdoctac{at} \coqdockw{level} 89, \coqdocvar{x} \coqdoctac{at} \coqdockw{level} 99, \coqdocvar{y} \coqdoctac{at} \coqdockw{level} 99,\coqdoceol
\coqdocindent{2.00em}
\coqdocvar{format} "'[v' 'while'  x  'do' '/  ' y '/' 'end' ']'") : \coqdocvar{com\_scope}.\coqdoceol
\coqdocemptyline
\end{coqdoccode}
\paragraph{Exercise: 2 stars, standard (himp\_ceval)}




    Now, we must extend the operational semantics. We have provided
   a template for the \coqref{SECF.Hoare.HoareAssertAssume.ceval}{\coqdocinductive{ceval}} relation below, specifying the big-step
   semantics. What rule(s) must be added to the definition of \coqref{SECF.Hoare.HoareAssertAssume.ceval}{\coqdocinductive{ceval}}
   to formalize the behavior of the \coqdocvar{HAVOC} command? 
\begin{coqdoccode}
\coqdocemptyline
\coqdocnoindent
\coqdockw{Reserved Notation}\coqdoceol
\coqdocindent{4.50em}
"st0 '=[' c ']=>' st1"\coqdoceol
\coqdocindent{4.50em}
(\coqdoctac{at} \coqdockw{level} 40, \coqdocvar{c} \coqdocvar{custom} \coqdocvar{com} \coqdoctac{at} \coqdockw{level} 99,\coqdoceol
\coqdocindent{5.00em}
\coqdocvar{st0} \coqdockw{constr}, \coqdocvar{st1} \coqdockw{constr} \coqdoctac{at} \coqdocvar{next} \coqdockw{level},\coqdoceol
\coqdocindent{5.00em}
\coqdocvar{format} "'[hv' st0  =[ '/  ' '[' c ']' '/' ]=>  st1 ']'").\coqdoceol
\coqdocemptyline
\coqdocnoindent
\coqdockw{Inductive} \coqdef{SECF.Equiv.Himp.ceval}{ceval}{\coqdocinductive{ceval}} : \coqref{SECF.Equiv.Himp.com}{\coqdocinductive{com}} \coqexternalref{::type scope:x '->' x}{http://rocq-prover.org/doc/V9.1.1/stdlib//Corelib.Init.Logic}{\coqdocnotation{\ensuremath{\rightarrow}}} \coqref{SECF.Imp.state}{\coqdocdefinition{state}} \coqexternalref{::type scope:x '->' x}{http://rocq-prover.org/doc/V9.1.1/stdlib//Corelib.Init.Logic}{\coqdocnotation{\ensuremath{\rightarrow}}} \coqref{SECF.Imp.state}{\coqdocdefinition{state}} \coqexternalref{::type scope:x '->' x}{http://rocq-prover.org/doc/V9.1.1/stdlib//Corelib.Init.Logic}{\coqdocnotation{\ensuremath{\rightarrow}}} \coqdockw{Prop} :=\coqdoceol
\coqdocindent{1.00em}
\ensuremath{|} \coqdef{SECF.Equiv.Himp.E Skip}{E\_Skip}{\coqdocconstructor{E\_Skip}} : \coqdockw{\ensuremath{\forall}} \coqdef{SECF.Equiv.st:137}{st}{\coqdocbinder{st}},\coqdoceol
\coqdocindent{3.00em}
\coqref{SECF.Equiv.st:137}{\coqdocvariable{st}} \coqref{SECF.Equiv.Himp.:::x '=[' x ']=>' x}{\coqdocnotation{=[}} \coqref{SECF.Equiv.Himp.:com:com scope:'skip'}{\coqdocnotation{skip}} \coqref{SECF.Equiv.Himp.:::x '=[' x ']=>' x}{\coqdocnotation{]=>}} \coqref{SECF.Equiv.st:137}{\coqdocvariable{st}}\coqdoceol
\coqdocindent{1.00em}
\ensuremath{|} \coqdef{SECF.Equiv.Himp.E Asgn}{E\_Asgn}{\coqdocconstructor{E\_Asgn}}  : \coqdockw{\ensuremath{\forall}} \coqdef{SECF.Equiv.st:138}{st}{\coqdocbinder{st}} \coqdef{SECF.Equiv.a:139}{a}{\coqdocbinder{a}} \coqdef{SECF.Equiv.n:140}{n}{\coqdocbinder{n}} \coqdef{SECF.Equiv.x:141}{x}{\coqdocbinder{x}},\coqdoceol
\coqdocindent{3.00em}
\coqref{SECF.Imp.aeval}{\coqdocdefinition{aeval}} \coqref{SECF.Equiv.st:138}{\coqdocvariable{st}} \coqref{SECF.Equiv.a:139}{\coqdocvariable{a}} \coqexternalref{::type scope:x '=' x}{http://rocq-prover.org/doc/V9.1.1/stdlib//Corelib.Init.Logic}{\coqdocnotation{=}} \coqref{SECF.Equiv.n:140}{\coqdocvariable{n}} \coqexternalref{::type scope:x '->' x}{http://rocq-prover.org/doc/V9.1.1/stdlib//Corelib.Init.Logic}{\coqdocnotation{\ensuremath{\rightarrow}}}\coqdoceol
\coqdocindent{3.00em}
\coqref{SECF.Equiv.st:138}{\coqdocvariable{st}} \coqref{SECF.Equiv.Himp.:::x '=[' x ']=>' x}{\coqdocnotation{=[}} \coqref{SECF.Equiv.x:141}{\coqdocvariable{x}} \coqref{SECF.Equiv.Himp.:com:com scope:x ':=' x}{\coqdocnotation{:=}} \coqref{SECF.Equiv.a:139}{\coqdocvariable{a}} \coqref{SECF.Equiv.Himp.:::x '=[' x ']=>' x}{\coqdocnotation{]=>}} \coqref{SECF.Equiv.Himp.:::x '=[' x ']=>' x}{\coqdocnotation{(}}\coqref{SECF.Equiv.x:141}{\coqdocvariable{x}} \coqref{SECF.Maps.:::x '!->' x ';' x}{\coqdocnotation{!->}} \coqref{SECF.Equiv.n:140}{\coqdocvariable{n}} \coqref{SECF.Maps.:::x '!->' x ';' x}{\coqdocnotation{;}} \coqref{SECF.Equiv.st:138}{\coqdocvariable{st}}\coqref{SECF.Equiv.Himp.:::x '=[' x ']=>' x}{\coqdocnotation{)}}\coqdoceol
\coqdocindent{1.00em}
\ensuremath{|} \coqdef{SECF.Equiv.Himp.E Seq}{E\_Seq}{\coqdocconstructor{E\_Seq}} : \coqdockw{\ensuremath{\forall}} \coqdef{SECF.Equiv.c1:142}{c1}{\coqdocbinder{c1}} \coqdef{SECF.Equiv.c2:143}{c2}{\coqdocbinder{c2}} \coqdef{SECF.Equiv.st:144}{st}{\coqdocbinder{st}} \coqdef{SECF.Equiv.st':145}{st'}{\coqdocbinder{st'}} \coqdef{SECF.Equiv.st'':146}{st'{}'}{\coqdocbinder{st'{}'}},\coqdoceol
\coqdocindent{3.00em}
\coqref{SECF.Equiv.st:144}{\coqdocvariable{st}}  \coqref{SECF.Equiv.Himp.:::x '=[' x ']=>' x}{\coqdocnotation{=[}} \coqref{SECF.Equiv.c1:142}{\coqdocvariable{c1}} \coqref{SECF.Equiv.Himp.:::x '=[' x ']=>' x}{\coqdocnotation{]=>}} \coqref{SECF.Equiv.st':145}{\coqdocvariable{st'}}  \coqexternalref{::type scope:x '->' x}{http://rocq-prover.org/doc/V9.1.1/stdlib//Corelib.Init.Logic}{\coqdocnotation{\ensuremath{\rightarrow}}}\coqdoceol
\coqdocindent{3.00em}
\coqref{SECF.Equiv.st':145}{\coqdocvariable{st'}} \coqref{SECF.Equiv.Himp.:::x '=[' x ']=>' x}{\coqdocnotation{=[}} \coqref{SECF.Equiv.c2:143}{\coqdocvariable{c2}} \coqref{SECF.Equiv.Himp.:::x '=[' x ']=>' x}{\coqdocnotation{]=>}} \coqref{SECF.Equiv.st'':146}{\coqdocvariable{st'{}'}} \coqexternalref{::type scope:x '->' x}{http://rocq-prover.org/doc/V9.1.1/stdlib//Corelib.Init.Logic}{\coqdocnotation{\ensuremath{\rightarrow}}}\coqdoceol
\coqdocindent{3.00em}
\coqref{SECF.Equiv.st:144}{\coqdocvariable{st}}  \coqref{SECF.Equiv.Himp.:::x '=[' x ']=>' x}{\coqdocnotation{=[}} \coqref{SECF.Equiv.c1:142}{\coqdocvariable{c1}} \coqref{SECF.Equiv.Himp.:com:com scope:x ';' x}{\coqdocnotation{;}} \coqref{SECF.Equiv.c2:143}{\coqdocvariable{c2}} \coqref{SECF.Equiv.Himp.:::x '=[' x ']=>' x}{\coqdocnotation{]=>}} \coqref{SECF.Equiv.st'':146}{\coqdocvariable{st'{}'}}\coqdoceol
\coqdocindent{1.00em}
\ensuremath{|} \coqdef{SECF.Equiv.Himp.E IfTrue}{E\_IfTrue}{\coqdocconstructor{E\_IfTrue}} : \coqdockw{\ensuremath{\forall}} \coqdef{SECF.Equiv.st:147}{st}{\coqdocbinder{st}} \coqdef{SECF.Equiv.st':148}{st'}{\coqdocbinder{st'}} \coqdef{SECF.Equiv.b:149}{b}{\coqdocbinder{b}} \coqdef{SECF.Equiv.c1:150}{c1}{\coqdocbinder{c1}} \coqdef{SECF.Equiv.c2:151}{c2}{\coqdocbinder{c2}},\coqdoceol
\coqdocindent{3.00em}
\coqref{SECF.Imp.beval}{\coqdocdefinition{beval}} \coqref{SECF.Equiv.st:147}{\coqdocvariable{st}} \coqref{SECF.Equiv.b:149}{\coqdocvariable{b}} \coqexternalref{::type scope:x '=' x}{http://rocq-prover.org/doc/V9.1.1/stdlib//Corelib.Init.Logic}{\coqdocnotation{=}} \coqref{SECF.Imp.:::'true'}{\coqdocnotation{true}} \coqexternalref{::type scope:x '->' x}{http://rocq-prover.org/doc/V9.1.1/stdlib//Corelib.Init.Logic}{\coqdocnotation{\ensuremath{\rightarrow}}}\coqdoceol
\coqdocindent{3.00em}
\coqref{SECF.Equiv.st:147}{\coqdocvariable{st}} \coqref{SECF.Equiv.Himp.:::x '=[' x ']=>' x}{\coqdocnotation{=[}} \coqref{SECF.Equiv.c1:150}{\coqdocvariable{c1}} \coqref{SECF.Equiv.Himp.:::x '=[' x ']=>' x}{\coqdocnotation{]=>}} \coqref{SECF.Equiv.st':148}{\coqdocvariable{st'}} \coqexternalref{::type scope:x '->' x}{http://rocq-prover.org/doc/V9.1.1/stdlib//Corelib.Init.Logic}{\coqdocnotation{\ensuremath{\rightarrow}}}\coqdoceol
\coqdocindent{3.00em}
\coqref{SECF.Equiv.st:147}{\coqdocvariable{st}} \coqref{SECF.Equiv.Himp.:::x '=[' x ']=>' x}{\coqdocnotation{=[}} \coqref{SECF.Equiv.Himp.:com:com scope:'if' x 'then' x 'else' x 'end'}{\coqdocnotation{if}} \coqref{SECF.Equiv.b:149}{\coqdocvariable{b}} \coqref{SECF.Equiv.Himp.:com:com scope:'if' x 'then' x 'else' x 'end'}{\coqdocnotation{then}} \coqref{SECF.Equiv.c1:150}{\coqdocvariable{c1}} \coqref{SECF.Equiv.Himp.:com:com scope:'if' x 'then' x 'else' x 'end'}{\coqdocnotation{else}} \coqref{SECF.Equiv.c2:151}{\coqdocvariable{c2}} \coqref{SECF.Equiv.Himp.:com:com scope:'if' x 'then' x 'else' x 'end'}{\coqdocnotation{end}} \coqref{SECF.Equiv.Himp.:::x '=[' x ']=>' x}{\coqdocnotation{]=>}} \coqref{SECF.Equiv.st':148}{\coqdocvariable{st'}}\coqdoceol
\coqdocindent{1.00em}
\ensuremath{|} \coqdef{SECF.Equiv.Himp.E IfFalse}{E\_IfFalse}{\coqdocconstructor{E\_IfFalse}} : \coqdockw{\ensuremath{\forall}} \coqdef{SECF.Equiv.st:152}{st}{\coqdocbinder{st}} \coqdef{SECF.Equiv.st':153}{st'}{\coqdocbinder{st'}} \coqdef{SECF.Equiv.b:154}{b}{\coqdocbinder{b}} \coqdef{SECF.Equiv.c1:155}{c1}{\coqdocbinder{c1}} \coqdef{SECF.Equiv.c2:156}{c2}{\coqdocbinder{c2}},\coqdoceol
\coqdocindent{3.00em}
\coqref{SECF.Imp.beval}{\coqdocdefinition{beval}} \coqref{SECF.Equiv.st:152}{\coqdocvariable{st}} \coqref{SECF.Equiv.b:154}{\coqdocvariable{b}} \coqexternalref{::type scope:x '=' x}{http://rocq-prover.org/doc/V9.1.1/stdlib//Corelib.Init.Logic}{\coqdocnotation{=}} \coqref{SECF.Imp.:::'false'}{\coqdocnotation{false}} \coqexternalref{::type scope:x '->' x}{http://rocq-prover.org/doc/V9.1.1/stdlib//Corelib.Init.Logic}{\coqdocnotation{\ensuremath{\rightarrow}}}\coqdoceol
\coqdocindent{3.00em}
\coqref{SECF.Equiv.st:152}{\coqdocvariable{st}} \coqref{SECF.Equiv.Himp.:::x '=[' x ']=>' x}{\coqdocnotation{=[}} \coqref{SECF.Equiv.c2:156}{\coqdocvariable{c2}} \coqref{SECF.Equiv.Himp.:::x '=[' x ']=>' x}{\coqdocnotation{]=>}} \coqref{SECF.Equiv.st':153}{\coqdocvariable{st'}} \coqexternalref{::type scope:x '->' x}{http://rocq-prover.org/doc/V9.1.1/stdlib//Corelib.Init.Logic}{\coqdocnotation{\ensuremath{\rightarrow}}}\coqdoceol
\coqdocindent{3.00em}
\coqref{SECF.Equiv.st:152}{\coqdocvariable{st}} \coqref{SECF.Equiv.Himp.:::x '=[' x ']=>' x}{\coqdocnotation{=[}} \coqref{SECF.Equiv.Himp.:com:com scope:'if' x 'then' x 'else' x 'end'}{\coqdocnotation{if}} \coqref{SECF.Equiv.b:154}{\coqdocvariable{b}} \coqref{SECF.Equiv.Himp.:com:com scope:'if' x 'then' x 'else' x 'end'}{\coqdocnotation{then}} \coqref{SECF.Equiv.c1:155}{\coqdocvariable{c1}} \coqref{SECF.Equiv.Himp.:com:com scope:'if' x 'then' x 'else' x 'end'}{\coqdocnotation{else}} \coqref{SECF.Equiv.c2:156}{\coqdocvariable{c2}} \coqref{SECF.Equiv.Himp.:com:com scope:'if' x 'then' x 'else' x 'end'}{\coqdocnotation{end}} \coqref{SECF.Equiv.Himp.:::x '=[' x ']=>' x}{\coqdocnotation{]=>}} \coqref{SECF.Equiv.st':153}{\coqdocvariable{st'}}\coqdoceol
\coqdocindent{1.00em}
\ensuremath{|} \coqdef{SECF.Equiv.Himp.E WhileFalse}{E\_WhileFalse}{\coqdocconstructor{E\_WhileFalse}} : \coqdockw{\ensuremath{\forall}} \coqdef{SECF.Equiv.b:157}{b}{\coqdocbinder{b}} \coqdef{SECF.Equiv.st:158}{st}{\coqdocbinder{st}} \coqdef{SECF.Equiv.c:159}{c}{\coqdocbinder{c}},\coqdoceol
\coqdocindent{3.00em}
\coqref{SECF.Imp.beval}{\coqdocdefinition{beval}} \coqref{SECF.Equiv.st:158}{\coqdocvariable{st}} \coqref{SECF.Equiv.b:157}{\coqdocvariable{b}} \coqexternalref{::type scope:x '=' x}{http://rocq-prover.org/doc/V9.1.1/stdlib//Corelib.Init.Logic}{\coqdocnotation{=}} \coqref{SECF.Imp.:::'false'}{\coqdocnotation{false}} \coqexternalref{::type scope:x '->' x}{http://rocq-prover.org/doc/V9.1.1/stdlib//Corelib.Init.Logic}{\coqdocnotation{\ensuremath{\rightarrow}}}\coqdoceol
\coqdocindent{3.00em}
\coqref{SECF.Equiv.st:158}{\coqdocvariable{st}} \coqref{SECF.Equiv.Himp.:::x '=[' x ']=>' x}{\coqdocnotation{=[}} \coqref{SECF.Equiv.Himp.:com:com scope:'while' x 'do' x 'end'}{\coqdocnotation{while}} \coqref{SECF.Equiv.b:157}{\coqdocvariable{b}} \coqref{SECF.Equiv.Himp.:com:com scope:'while' x 'do' x 'end'}{\coqdocnotation{do}} \coqref{SECF.Equiv.c:159}{\coqdocvariable{c}} \coqref{SECF.Equiv.Himp.:com:com scope:'while' x 'do' x 'end'}{\coqdocnotation{end}} \coqref{SECF.Equiv.Himp.:::x '=[' x ']=>' x}{\coqdocnotation{]=>}} \coqref{SECF.Equiv.st:158}{\coqdocvariable{st}}\coqdoceol
\coqdocindent{1.00em}
\ensuremath{|} \coqdef{SECF.Equiv.Himp.E WhileTrue}{E\_WhileTrue}{\coqdocconstructor{E\_WhileTrue}} : \coqdockw{\ensuremath{\forall}} \coqdef{SECF.Equiv.st:160}{st}{\coqdocbinder{st}} \coqdef{SECF.Equiv.st':161}{st'}{\coqdocbinder{st'}} \coqdef{SECF.Equiv.st'':162}{st'{}'}{\coqdocbinder{st'{}'}} \coqdef{SECF.Equiv.b:163}{b}{\coqdocbinder{b}} \coqdef{SECF.Equiv.c:164}{c}{\coqdocbinder{c}},\coqdoceol
\coqdocindent{3.00em}
\coqref{SECF.Imp.beval}{\coqdocdefinition{beval}} \coqref{SECF.Equiv.st:160}{\coqdocvariable{st}} \coqref{SECF.Equiv.b:163}{\coqdocvariable{b}} \coqexternalref{::type scope:x '=' x}{http://rocq-prover.org/doc/V9.1.1/stdlib//Corelib.Init.Logic}{\coqdocnotation{=}} \coqref{SECF.Imp.:::'true'}{\coqdocnotation{true}} \coqexternalref{::type scope:x '->' x}{http://rocq-prover.org/doc/V9.1.1/stdlib//Corelib.Init.Logic}{\coqdocnotation{\ensuremath{\rightarrow}}}\coqdoceol
\coqdocindent{3.00em}
\coqref{SECF.Equiv.st:160}{\coqdocvariable{st}}  \coqref{SECF.Equiv.Himp.:::x '=[' x ']=>' x}{\coqdocnotation{=[}} \coqref{SECF.Equiv.c:164}{\coqdocvariable{c}} \coqref{SECF.Equiv.Himp.:::x '=[' x ']=>' x}{\coqdocnotation{]=>}} \coqref{SECF.Equiv.st':161}{\coqdocvariable{st'}} \coqexternalref{::type scope:x '->' x}{http://rocq-prover.org/doc/V9.1.1/stdlib//Corelib.Init.Logic}{\coqdocnotation{\ensuremath{\rightarrow}}}\coqdoceol
\coqdocindent{3.00em}
\coqref{SECF.Equiv.st':161}{\coqdocvariable{st'}} \coqref{SECF.Equiv.Himp.:::x '=[' x ']=>' x}{\coqdocnotation{=[}} \coqref{SECF.Equiv.Himp.:com:com scope:'while' x 'do' x 'end'}{\coqdocnotation{while}} \coqref{SECF.Equiv.b:163}{\coqdocvariable{b}} \coqref{SECF.Equiv.Himp.:com:com scope:'while' x 'do' x 'end'}{\coqdocnotation{do}} \coqref{SECF.Equiv.c:164}{\coqdocvariable{c}} \coqref{SECF.Equiv.Himp.:com:com scope:'while' x 'do' x 'end'}{\coqdocnotation{end}} \coqref{SECF.Equiv.Himp.:::x '=[' x ']=>' x}{\coqdocnotation{]=>}} \coqref{SECF.Equiv.st'':162}{\coqdocvariable{st'{}'}} \coqexternalref{::type scope:x '->' x}{http://rocq-prover.org/doc/V9.1.1/stdlib//Corelib.Init.Logic}{\coqdocnotation{\ensuremath{\rightarrow}}}\coqdoceol
\coqdocindent{3.00em}
\coqref{SECF.Equiv.st:160}{\coqdocvariable{st}}  \coqref{SECF.Equiv.Himp.:::x '=[' x ']=>' x}{\coqdocnotation{=[}} \coqref{SECF.Equiv.Himp.:com:com scope:'while' x 'do' x 'end'}{\coqdocnotation{while}} \coqref{SECF.Equiv.b:163}{\coqdocvariable{b}} \coqref{SECF.Equiv.Himp.:com:com scope:'while' x 'do' x 'end'}{\coqdocnotation{do}} \coqref{SECF.Equiv.c:164}{\coqdocvariable{c}} \coqref{SECF.Equiv.Himp.:com:com scope:'while' x 'do' x 'end'}{\coqdocnotation{end}} \coqref{SECF.Equiv.Himp.:::x '=[' x ']=>' x}{\coqdocnotation{]=>}} \coqref{SECF.Equiv.st'':162}{\coqdocvariable{st'{}'}}\coqdoceol
\coqdocnoindent
\coqdoceol
\coqdocnoindent
\coqdoceol
\coqdocindent{1.00em}
\coqdockw{where} \coqdef{SECF.Equiv.Himp.:::x '=[' x ']=>' x}{"}{"}st =[ c ]=> st'" := (\coqref{SECF.Equiv.ceval:136}{\coqdocinductive{ceval}} \coqdocvar{c} \coqdocvar{st} \coqdocvar{st'}).\coqdoceol
\coqdocemptyline
\end{coqdoccode}
As a sanity check, the following claims should be provable for
    your definition: 
\begin{coqdoccode}
\coqdocemptyline
\coqdocnoindent
\coqdockw{Example} \coqdef{SECF.Equiv.Himp.havoc example1}{havoc\_example1}{\coqdocdefinition{havoc\_example1}} : \coqref{SECF.Imp.empty st}{\coqdocdefinition{empty\_st}} \coqref{SECF.Equiv.Himp.:::x '=[' x ']=>' x}{\coqdocnotation{=[}} \coqref{SECF.Equiv.Himp.:com::'havoc' x}{\coqdocnotation{havoc}} \coqref{SECF.Imp.X}{\coqdocdefinition{X}} \coqref{SECF.Equiv.Himp.:::x '=[' x ']=>' x}{\coqdocnotation{]=>}} \coqref{SECF.Equiv.Himp.:::x '=[' x ']=>' x}{\coqdocnotation{(}}\coqref{SECF.Imp.X}{\coqdocdefinition{X}} \coqref{SECF.Imp.:::x '!->' x}{\coqdocnotation{!->}} 0\coqref{SECF.Equiv.Himp.:::x '=[' x ']=>' x}{\coqdocnotation{)}}.\coqdoceol
\coqdocnoindent
\coqdockw{Proof}.\coqdoceol
 \coqdocvar{Admitted}.\coqdoceol
\coqdocemptyline
\coqdocnoindent
\coqdockw{Example} \coqdef{SECF.Equiv.Himp.havoc example2}{havoc\_example2}{\coqdocdefinition{havoc\_example2}} :\coqdoceol
\coqdocindent{1.00em}
\coqref{SECF.Imp.empty st}{\coqdocdefinition{empty\_st}} \coqref{SECF.Equiv.Himp.:::x '=[' x ']=>' x}{\coqdocnotation{=[}} \coqref{SECF.Equiv.Himp.:com:com scope:'skip'}{\coqdocnotation{skip}}\coqref{SECF.Equiv.Himp.:com:com scope:x ';' x}{\coqdocnotation{;}} \coqref{SECF.Equiv.Himp.:com::'havoc' x}{\coqdocnotation{havoc}} \coqref{SECF.Imp.Z}{\coqdocdefinition{Z}} \coqref{SECF.Equiv.Himp.:::x '=[' x ']=>' x}{\coqdocnotation{]=>}} \coqref{SECF.Equiv.Himp.:::x '=[' x ']=>' x}{\coqdocnotation{(}}\coqref{SECF.Imp.Z}{\coqdocdefinition{Z}} \coqref{SECF.Imp.:::x '!->' x}{\coqdocnotation{!->}} 42\coqref{SECF.Equiv.Himp.:::x '=[' x ']=>' x}{\coqdocnotation{)}}.\coqdoceol
\coqdocnoindent
\coqdockw{Proof}.\coqdoceol
 \coqdocvar{Admitted}.\coqdoceol
\coqdocemptyline
\coqdocnoindent
\coqdockw{Definition} \coqdef{SECF.Equiv.Himp.manual grade for Check rule for HAVOC}{manual\_grade\_for\_Check\_rule\_for\_HAVOC}{\coqdocdefinition{manual\_grade\_for\_Check\_rule\_for\_HAVOC}} : \coqexternalref{option}{http://rocq-prover.org/doc/V9.1.1/stdlib//Corelib.Init.Datatypes}{\coqdocinductive{option}} (\coqexternalref{nat}{http://rocq-prover.org/doc/V9.1.1/stdlib//Corelib.Init.Datatypes}{\coqdocinductive{nat}}\coqexternalref{::type scope:x '*' x}{http://rocq-prover.org/doc/V9.1.1/stdlib//Corelib.Init.Datatypes}{\coqdocnotation{\ensuremath{\times}}}\coqdocinductive{string}) := \coqexternalref{None}{http://rocq-prover.org/doc/V9.1.1/stdlib//Corelib.Init.Datatypes}{\coqdocconstructor{None}}.\coqdoceol
\end{coqdoccode}
\ensuremath{\Box} 

 Finally, we repeat the definition of command equivalence from above: 
\begin{coqdoccode}
\coqdocemptyline
\coqdocnoindent
\coqdockw{Definition} \coqdef{SECF.Equiv.Himp.cequiv}{cequiv}{\coqdocdefinition{cequiv}} (\coqdef{SECF.Equiv.c1:165}{c1}{\coqdocbinder{c1}} \coqdef{SECF.Equiv.c2:166}{c2}{\coqdocbinder{c2}} : \coqref{SECF.Equiv.Himp.com}{\coqdocinductive{com}}) : \coqdockw{Prop} := \coqdockw{\ensuremath{\forall}} \coqdef{SECF.Equiv.st:167}{st}{\coqdocbinder{st}} \coqdef{SECF.Equiv.st':168}{st'}{\coqdocbinder{st'}} : \coqref{SECF.Imp.state}{\coqdocdefinition{state}},\coqdoceol
\coqdocindent{1.00em}
\coqref{SECF.Equiv.st:167}{\coqdocvariable{st}} \coqref{SECF.Equiv.Himp.:::x '=[' x ']=>' x}{\coqdocnotation{=[}} \coqref{SECF.Equiv.c1:165}{\coqdocvariable{c1}} \coqref{SECF.Equiv.Himp.:::x '=[' x ']=>' x}{\coqdocnotation{]=>}} \coqref{SECF.Equiv.st':168}{\coqdocvariable{st'}} \coqexternalref{::type scope:x '<->' x}{http://rocq-prover.org/doc/V9.1.1/stdlib//Corelib.Init.Logic}{\coqdocnotation{\ensuremath{\leftrightarrow}}} \coqref{SECF.Equiv.st:167}{\coqdocvariable{st}} \coqref{SECF.Equiv.Himp.:::x '=[' x ']=>' x}{\coqdocnotation{=[}} \coqref{SECF.Equiv.c2:166}{\coqdocvariable{c2}} \coqref{SECF.Equiv.Himp.:::x '=[' x ']=>' x}{\coqdocnotation{]=>}} \coqref{SECF.Equiv.st':168}{\coqdocvariable{st'}}.\coqdoceol
\coqdocemptyline
\end{coqdoccode}
Let's apply this definition to prove some nondeterministic
    programs equivalent / inequivalent. 

\paragraph{Exercise: 3 stars, standard (havoc\_swap)}




    Are the following two programs equivalent? 
\begin{coqdoccode}
\coqdocemptyline
\coqdocnoindent
\coqdockw{Definition} \coqdef{SECF.Equiv.Himp.pXY}{pXY}{\coqdocdefinition{pXY}} :=\coqdoceol
\coqdocindent{1.00em}
\coqref{SECF.Imp.::com scope:'<x7B' x 'x7D>'}{\coqdocnotation{<\{}} \coqref{SECF.Equiv.Himp.:com::'havoc' x}{\coqdocnotation{havoc}} \coqref{SECF.Imp.X}{\coqdocdefinition{X}} \coqref{SECF.Equiv.Himp.:com:com scope:x ';' x}{\coqdocnotation{;}} \coqref{SECF.Equiv.Himp.:com::'havoc' x}{\coqdocnotation{havoc}} \coqref{SECF.Imp.Y}{\coqdocdefinition{Y}} \coqref{SECF.Imp.::com scope:'<x7B' x 'x7D>'}{\coqdocnotation{\}>}}.\coqdoceol
\coqdocemptyline
\coqdocnoindent
\coqdockw{Definition} \coqdef{SECF.Equiv.Himp.pYX}{pYX}{\coqdocdefinition{pYX}} :=\coqdoceol
\coqdocindent{1.00em}
\coqref{SECF.Imp.::com scope:'<x7B' x 'x7D>'}{\coqdocnotation{<\{}} \coqref{SECF.Equiv.Himp.:com::'havoc' x}{\coqdocnotation{havoc}} \coqref{SECF.Imp.Y}{\coqdocdefinition{Y}}\coqref{SECF.Equiv.Himp.:com:com scope:x ';' x}{\coqdocnotation{;}} \coqref{SECF.Equiv.Himp.:com::'havoc' x}{\coqdocnotation{havoc}} \coqref{SECF.Imp.X}{\coqdocdefinition{X}} \coqref{SECF.Imp.::com scope:'<x7B' x 'x7D>'}{\coqdocnotation{\}>}}.\coqdoceol
\coqdocemptyline
\end{coqdoccode}
If you think they are equivalent, prove it. If you think they are
    not, prove that. 
\begin{coqdoccode}
\coqdocemptyline
\coqdocnoindent
\coqdockw{Theorem} \coqdef{SECF.Equiv.Himp.pXY cequiv pYX}{pXY\_cequiv\_pYX}{\coqdoclemma{pXY\_cequiv\_pYX}} :\coqdoceol
\coqdocindent{1.00em}
\coqref{SECF.Equiv.Himp.cequiv}{\coqdocdefinition{cequiv}} \coqref{SECF.Equiv.Himp.pXY}{\coqdocdefinition{pXY}} \coqref{SECF.Equiv.Himp.pYX}{\coqdocdefinition{pYX}} \coqexternalref{::type scope:x 'x5C/' x}{http://rocq-prover.org/doc/V9.1.1/stdlib//Corelib.Init.Logic}{\coqdocnotation{\ensuremath{\lor}}} \coqexternalref{::type scope:'x7E' x}{http://rocq-prover.org/doc/V9.1.1/stdlib//Corelib.Init.Logic}{\coqdocnotation{\ensuremath{\lnot}}}\coqref{SECF.Equiv.Himp.cequiv}{\coqdocdefinition{cequiv}} \coqref{SECF.Equiv.Himp.pXY}{\coqdocdefinition{pXY}} \coqref{SECF.Equiv.Himp.pYX}{\coqdocdefinition{pYX}}.\coqdoceol
\coqdocnoindent
\coqdockw{Proof}.\coqdoceol
 \coqdocvar{Admitted}.\coqdoceol
\end{coqdoccode}
\ensuremath{\Box} 

\paragraph{Exercise: 4 stars, standard, optional (havoc\_copy)}




    Are the following two programs equivalent? 
\begin{coqdoccode}
\coqdocemptyline
\coqdocnoindent
\coqdockw{Definition} \coqdef{SECF.Equiv.Himp.ptwice}{ptwice}{\coqdocdefinition{ptwice}} :=\coqdoceol
\coqdocindent{1.00em}
\coqref{SECF.Imp.::com scope:'<x7B' x 'x7D>'}{\coqdocnotation{<\{}} \coqref{SECF.Equiv.Himp.:com::'havoc' x}{\coqdocnotation{havoc}} \coqref{SECF.Imp.X}{\coqdocdefinition{X}}\coqref{SECF.Equiv.Himp.:com:com scope:x ';' x}{\coqdocnotation{;}} \coqref{SECF.Equiv.Himp.:com::'havoc' x}{\coqdocnotation{havoc}} \coqref{SECF.Imp.Y}{\coqdocdefinition{Y}} \coqref{SECF.Imp.::com scope:'<x7B' x 'x7D>'}{\coqdocnotation{\}>}}.\coqdoceol
\coqdocemptyline
\coqdocnoindent
\coqdockw{Definition} \coqdef{SECF.Equiv.Himp.pcopy}{pcopy}{\coqdocdefinition{pcopy}} :=\coqdoceol
\coqdocindent{1.00em}
\coqref{SECF.Imp.::com scope:'<x7B' x 'x7D>'}{\coqdocnotation{<\{}} \coqref{SECF.Equiv.Himp.:com::'havoc' x}{\coqdocnotation{havoc}} \coqref{SECF.Imp.X}{\coqdocdefinition{X}}\coqref{SECF.Equiv.Himp.:com:com scope:x ';' x}{\coqdocnotation{;}} \coqref{SECF.Imp.Y}{\coqdocdefinition{Y}} \coqref{SECF.Equiv.Himp.:com:com scope:x ':=' x}{\coqdocnotation{:=}} \coqref{SECF.Imp.X}{\coqdocdefinition{X}} \coqref{SECF.Imp.::com scope:'<x7B' x 'x7D>'}{\coqdocnotation{\}>}}.\coqdoceol
\coqdocemptyline
\end{coqdoccode}
If you think they are equivalent, then prove it. If you think they
    are not, then prove that.  (Hint: You may find the \coqdoctac{assert} tactic
    useful.) 
\begin{coqdoccode}
\coqdocemptyline
\coqdocnoindent
\coqdockw{Theorem} \coqdef{SECF.Equiv.Himp.ptwice cequiv pcopy}{ptwice\_cequiv\_pcopy}{\coqdoclemma{ptwice\_cequiv\_pcopy}} :\coqdoceol
\coqdocindent{1.00em}
\coqref{SECF.Equiv.Himp.cequiv}{\coqdocdefinition{cequiv}} \coqref{SECF.Equiv.Himp.ptwice}{\coqdocdefinition{ptwice}} \coqref{SECF.Equiv.Himp.pcopy}{\coqdocdefinition{pcopy}} \coqexternalref{::type scope:x 'x5C/' x}{http://rocq-prover.org/doc/V9.1.1/stdlib//Corelib.Init.Logic}{\coqdocnotation{\ensuremath{\lor}}} \coqexternalref{::type scope:'x7E' x}{http://rocq-prover.org/doc/V9.1.1/stdlib//Corelib.Init.Logic}{\coqdocnotation{\ensuremath{\lnot}}}\coqref{SECF.Equiv.Himp.cequiv}{\coqdocdefinition{cequiv}} \coqref{SECF.Equiv.Himp.ptwice}{\coqdocdefinition{ptwice}} \coqref{SECF.Equiv.Himp.pcopy}{\coqdocdefinition{pcopy}}.\coqdoceol
\coqdocnoindent
\coqdockw{Proof}.  \coqdocvar{Admitted}.\coqdoceol
\end{coqdoccode}
\ensuremath{\Box} 

 The definition of program equivalence we are using here has some
    subtle consequences on programs that may loop forever.  What
    \coqref{SECF.Equiv.Himp.cequiv}{\coqdocdefinition{cequiv}} says is that the set of possible \textit{terminating} outcomes
    of two equivalent programs is the same. However, in a language
    with nondeterminism, like Himp, some programs always terminate,
    some programs always diverge, and some programs can
    nondeterministically terminate in some runs and diverge in
    others. The final part of the following exercise illustrates this
    phenomenon.


\paragraph{Exercise: 4 stars, advanced (p1\_p2\_term)}




    Consider the following commands: 
\begin{coqdoccode}
\coqdocemptyline
\coqdocnoindent
\coqdockw{Definition} \coqdef{SECF.Equiv.Himp.p1}{p1}{\coqdocdefinition{p1}} : \coqref{SECF.Equiv.Himp.com}{\coqdocinductive{com}} :=\coqdoceol
\coqdocindent{1.00em}
\coqref{SECF.Imp.::com scope:'<x7B' x 'x7D>'}{\coqdocnotation{<\{}} \coqref{SECF.Equiv.Himp.:com:com scope:'while' x 'do' x 'end'}{\coqdocnotation{while}} \coqref{SECF.Imp.:com::'x7E' x}{\coqdocnotation{\ensuremath{\lnot}}} \coqref{SECF.Imp.:com::'(' x ')'}{\coqdocnotation{(}}\coqref{SECF.Imp.X}{\coqdocdefinition{X}} \coqref{SECF.Imp.:com::x '=' x}{\coqdocnotation{=}} 0\coqref{SECF.Imp.:com::'(' x ')'}{\coqdocnotation{)}} \coqref{SECF.Equiv.Himp.:com:com scope:'while' x 'do' x 'end'}{\coqdocnotation{do}}\coqdoceol
\coqdocindent{3.50em}
\coqref{SECF.Equiv.Himp.:com::'havoc' x}{\coqdocnotation{havoc}} \coqref{SECF.Imp.Y}{\coqdocdefinition{Y}}\coqref{SECF.Equiv.Himp.:com:com scope:x ';' x}{\coqdocnotation{;}}\coqdoceol
\coqdocindent{3.50em}
\coqref{SECF.Imp.X}{\coqdocdefinition{X}} \coqref{SECF.Equiv.Himp.:com:com scope:x ':=' x}{\coqdocnotation{:=}} \coqref{SECF.Imp.X}{\coqdocdefinition{X}} \coqref{SECF.Imp.:com::x '+' x}{\coqdocnotation{+}} 1\coqdoceol
\coqdocindent{2.50em}
\coqref{SECF.Equiv.Himp.:com:com scope:'while' x 'do' x 'end'}{\coqdocnotation{end}} \coqref{SECF.Imp.::com scope:'<x7B' x 'x7D>'}{\coqdocnotation{\}>}}.\coqdoceol
\coqdocemptyline
\coqdocnoindent
\coqdockw{Definition} \coqdef{SECF.Equiv.Himp.p2}{p2}{\coqdocdefinition{p2}} : \coqref{SECF.Equiv.Himp.com}{\coqdocinductive{com}} :=\coqdoceol
\coqdocindent{1.00em}
\coqref{SECF.Imp.::com scope:'<x7B' x 'x7D>'}{\coqdocnotation{<\{}} \coqref{SECF.Equiv.Himp.:com:com scope:'while' x 'do' x 'end'}{\coqdocnotation{while}} \coqref{SECF.Imp.:com::'x7E' x}{\coqdocnotation{\ensuremath{\lnot}}} \coqref{SECF.Imp.:com::'(' x ')'}{\coqdocnotation{(}}\coqref{SECF.Imp.X}{\coqdocdefinition{X}} \coqref{SECF.Imp.:com::x '=' x}{\coqdocnotation{=}} 0\coqref{SECF.Imp.:com::'(' x ')'}{\coqdocnotation{)}} \coqref{SECF.Equiv.Himp.:com:com scope:'while' x 'do' x 'end'}{\coqdocnotation{do}}\coqdoceol
\coqdocindent{3.50em}
\coqref{SECF.Equiv.Himp.:com:com scope:'skip'}{\coqdocnotation{skip}}\coqdoceol
\coqdocindent{2.50em}
\coqref{SECF.Equiv.Himp.:com:com scope:'while' x 'do' x 'end'}{\coqdocnotation{end}} \coqref{SECF.Imp.::com scope:'<x7B' x 'x7D>'}{\coqdocnotation{\}>}}.\coqdoceol
\coqdocemptyline
\end{coqdoccode}
Intuitively, \coqref{SECF.Equiv.Himp.p1}{\coqdocdefinition{p1}} and \coqref{SECF.Equiv.Himp.p2}{\coqdocdefinition{p2}} have the same termination behavior:
    either they loop forever, or they terminate in the same state they
    started in.  We can capture the termination behavior of \coqref{SECF.Equiv.Himp.p1}{\coqdocdefinition{p1}} and
    \coqref{SECF.Equiv.Himp.p2}{\coqdocdefinition{p2}} individually with these lemmas: 
\begin{coqdoccode}
\coqdocemptyline
\coqdocnoindent
\coqdockw{Lemma} \coqdef{SECF.Equiv.Himp.p1 may diverge}{p1\_may\_diverge}{\coqdoclemma{p1\_may\_diverge}} : \coqdockw{\ensuremath{\forall}} \coqdef{SECF.Equiv.st:169}{st}{\coqdocbinder{st}} \coqdef{SECF.Equiv.st':170}{st'}{\coqdocbinder{st'}}, \coqref{SECF.Equiv.st:169}{\coqdocvariable{st}} \coqref{SECF.Imp.X}{\coqdocdefinition{X}} \coqexternalref{::type scope:x '<>' x}{http://rocq-prover.org/doc/V9.1.1/stdlib//Corelib.Init.Logic}{\coqdocnotation{\ensuremath{\not=}}} 0 \coqexternalref{::type scope:x '->' x}{http://rocq-prover.org/doc/V9.1.1/stdlib//Corelib.Init.Logic}{\coqdocnotation{\ensuremath{\rightarrow}}}\coqdoceol
\coqdocindent{1.00em}
\coqexternalref{::type scope:'x7E' x}{http://rocq-prover.org/doc/V9.1.1/stdlib//Corelib.Init.Logic}{\coqdocnotation{\ensuremath{\lnot}}} \coqref{SECF.Equiv.st:169}{\coqdocvariable{st}} \coqref{SECF.Equiv.Himp.:::x '=[' x ']=>' x}{\coqdocnotation{=[}} \coqref{SECF.Equiv.Himp.p1}{\coqdocdefinition{p1}} \coqref{SECF.Equiv.Himp.:::x '=[' x ']=>' x}{\coqdocnotation{]=>}} \coqref{SECF.Equiv.st':170}{\coqdocvariable{st'}}.\coqdoceol
\coqdocnoindent
\coqdockw{Proof}.  \coqdocvar{Admitted}.\coqdoceol
\coqdocemptyline
\coqdocnoindent
\coqdockw{Lemma} \coqdef{SECF.Equiv.Himp.p2 may diverge}{p2\_may\_diverge}{\coqdoclemma{p2\_may\_diverge}} : \coqdockw{\ensuremath{\forall}} \coqdef{SECF.Equiv.st:171}{st}{\coqdocbinder{st}} \coqdef{SECF.Equiv.st':172}{st'}{\coqdocbinder{st'}}, \coqref{SECF.Equiv.st:171}{\coqdocvariable{st}} \coqref{SECF.Imp.X}{\coqdocdefinition{X}} \coqexternalref{::type scope:x '<>' x}{http://rocq-prover.org/doc/V9.1.1/stdlib//Corelib.Init.Logic}{\coqdocnotation{\ensuremath{\not=}}} 0 \coqexternalref{::type scope:x '->' x}{http://rocq-prover.org/doc/V9.1.1/stdlib//Corelib.Init.Logic}{\coqdocnotation{\ensuremath{\rightarrow}}}\coqdoceol
\coqdocindent{1.00em}
\coqexternalref{::type scope:'x7E' x}{http://rocq-prover.org/doc/V9.1.1/stdlib//Corelib.Init.Logic}{\coqdocnotation{\ensuremath{\lnot}}} \coqref{SECF.Equiv.st:171}{\coqdocvariable{st}} \coqref{SECF.Equiv.Himp.:::x '=[' x ']=>' x}{\coqdocnotation{=[}} \coqref{SECF.Equiv.Himp.p2}{\coqdocdefinition{p2}} \coqref{SECF.Equiv.Himp.:::x '=[' x ']=>' x}{\coqdocnotation{]=>}} \coqref{SECF.Equiv.st':172}{\coqdocvariable{st'}}.\coqdoceol
\coqdocnoindent
\coqdockw{Proof}.\coqdoceol
 \coqdocvar{Admitted}.\coqdoceol
\end{coqdoccode}
\ensuremath{\Box} 

\paragraph{Exercise: 4 stars, advanced (p1\_p2\_equiv)}




    Use these two lemmas to prove that \coqref{SECF.Equiv.Himp.p1}{\coqdocdefinition{p1}} and \coqref{SECF.Equiv.Himp.p2}{\coqdocdefinition{p2}} are actually
    equivalent. 
\begin{coqdoccode}
\coqdocemptyline
\coqdocnoindent
\coqdockw{Theorem} \coqdef{SECF.Equiv.Himp.p1 p2 equiv}{p1\_p2\_equiv}{\coqdoclemma{p1\_p2\_equiv}} : \coqref{SECF.Equiv.Himp.cequiv}{\coqdocdefinition{cequiv}} \coqref{SECF.Equiv.Himp.p1}{\coqdocdefinition{p1}} \coqref{SECF.Equiv.Himp.p2}{\coqdocdefinition{p2}}.\coqdoceol
\coqdocnoindent
\coqdockw{Proof}.  \coqdocvar{Admitted}.\coqdoceol
\end{coqdoccode}
\ensuremath{\Box} 

\paragraph{Exercise: 4 stars, advanced (p3\_p4\_inequiv)}




    Prove that the following programs are \textit{not} equivalent.  (Hint:
    What should the value of \coqref{SECF.SpecCT.Z}{\coqdocdefinition{Z}} be when \coqref{SECF.Equiv.Himp.p3}{\coqdocdefinition{p3}} terminates?  What about
    \coqref{SECF.Equiv.Himp.p4}{\coqdocdefinition{p4}}?) 
\begin{coqdoccode}
\coqdocemptyline
\coqdocnoindent
\coqdockw{Definition} \coqdef{SECF.Equiv.Himp.p3}{p3}{\coqdocdefinition{p3}} : \coqref{SECF.Equiv.Himp.com}{\coqdocinductive{com}} :=\coqdoceol
\coqdocindent{1.00em}
\coqref{SECF.Imp.::com scope:'<x7B' x 'x7D>'}{\coqdocnotation{<\{}} \coqref{SECF.Imp.Z}{\coqdocdefinition{Z}} \coqref{SECF.Equiv.Himp.:com:com scope:x ':=' x}{\coqdocnotation{:=}} 1\coqref{SECF.Equiv.Himp.:com:com scope:x ';' x}{\coqdocnotation{;}}\coqdoceol
\coqdocindent{2.50em}
\coqref{SECF.Equiv.Himp.:com:com scope:'while' x 'do' x 'end'}{\coqdocnotation{while}} \coqref{SECF.Imp.X}{\coqdocdefinition{X}} \coqref{SECF.Imp.:com::x '<>' x}{\coqdocnotation{\ensuremath{\not=}}} 0 \coqref{SECF.Equiv.Himp.:com:com scope:'while' x 'do' x 'end'}{\coqdocnotation{do}}\coqdoceol
\coqdocindent{3.50em}
\coqref{SECF.Equiv.Himp.:com::'havoc' x}{\coqdocnotation{havoc}} \coqref{SECF.Imp.X}{\coqdocdefinition{X}}\coqref{SECF.Equiv.Himp.:com:com scope:x ';' x}{\coqdocnotation{;}}\coqdoceol
\coqdocindent{3.50em}
\coqref{SECF.Equiv.Himp.:com::'havoc' x}{\coqdocnotation{havoc}} \coqref{SECF.Imp.Z}{\coqdocdefinition{Z}}\coqdoceol
\coqdocindent{2.50em}
\coqref{SECF.Equiv.Himp.:com:com scope:'while' x 'do' x 'end'}{\coqdocnotation{end}} \coqref{SECF.Imp.::com scope:'<x7B' x 'x7D>'}{\coqdocnotation{\}>}}.\coqdoceol
\coqdocemptyline
\coqdocnoindent
\coqdockw{Definition} \coqdef{SECF.Equiv.Himp.p4}{p4}{\coqdocdefinition{p4}} : \coqref{SECF.Equiv.Himp.com}{\coqdocinductive{com}} :=\coqdoceol
\coqdocindent{1.00em}
\coqref{SECF.Imp.::com scope:'<x7B' x 'x7D>'}{\coqdocnotation{<\{}} \coqref{SECF.Imp.X}{\coqdocdefinition{X}} \coqref{SECF.Equiv.Himp.:com:com scope:x ':=' x}{\coqdocnotation{:=}} 0\coqref{SECF.Equiv.Himp.:com:com scope:x ';' x}{\coqdocnotation{;}}\coqdoceol
\coqdocindent{2.50em}
\coqref{SECF.Imp.Z}{\coqdocdefinition{Z}} \coqref{SECF.Equiv.Himp.:com:com scope:x ':=' x}{\coqdocnotation{:=}} 1 \coqref{SECF.Imp.::com scope:'<x7B' x 'x7D>'}{\coqdocnotation{\}>}}.\coqdoceol
\coqdocemptyline
\coqdocnoindent
\coqdockw{Theorem} \coqdef{SECF.Equiv.Himp.p3 p4 inequiv}{p3\_p4\_inequiv}{\coqdoclemma{p3\_p4\_inequiv}} : \coqexternalref{::type scope:'x7E' x}{http://rocq-prover.org/doc/V9.1.1/stdlib//Corelib.Init.Logic}{\coqdocnotation{\ensuremath{\lnot}}} \coqref{SECF.Equiv.Himp.cequiv}{\coqdocdefinition{cequiv}} \coqref{SECF.Equiv.Himp.p3}{\coqdocdefinition{p3}} \coqref{SECF.Equiv.Himp.p4}{\coqdocdefinition{p4}}.\coqdoceol
\coqdocnoindent
\coqdockw{Proof}.  \coqdocvar{Admitted}.\coqdoceol
\end{coqdoccode}
\ensuremath{\Box} 

\paragraph{Exercise: 5 stars, advanced, optional (p5\_p6\_equiv)}




    Prove that the following commands are equivalent.  (Hint: As
    mentioned above, our definition of \coqref{SECF.Equiv.Himp.cequiv}{\coqdocdefinition{cequiv}} for Himp only takes
    into account the sets of possible terminating configurations: two
    programs are equivalent if and only if the set of possible terminating
    states is the same for both programs when given a same starting state
    \coqdocvar{st}.  If \coqref{SECF.Equiv.Himp.p5}{\coqdocdefinition{p5}} terminates, what should the final state be? Conversely,
    is it always possible to make \coqref{SECF.Equiv.Himp.p5}{\coqdocdefinition{p5}} terminate?) 
\begin{coqdoccode}
\coqdocemptyline
\coqdocnoindent
\coqdockw{Definition} \coqdef{SECF.Equiv.Himp.p5}{p5}{\coqdocdefinition{p5}} : \coqref{SECF.Equiv.Himp.com}{\coqdocinductive{com}} :=\coqdoceol
\coqdocindent{1.00em}
\coqref{SECF.Imp.::com scope:'<x7B' x 'x7D>'}{\coqdocnotation{<\{}} \coqref{SECF.Equiv.Himp.:com:com scope:'while' x 'do' x 'end'}{\coqdocnotation{while}} \coqref{SECF.Imp.X}{\coqdocdefinition{X}} \coqref{SECF.Imp.:com::x '<>' x}{\coqdocnotation{\ensuremath{\not=}}} 1 \coqref{SECF.Equiv.Himp.:com:com scope:'while' x 'do' x 'end'}{\coqdocnotation{do}}\coqdoceol
\coqdocindent{3.50em}
\coqref{SECF.Equiv.Himp.:com::'havoc' x}{\coqdocnotation{havoc}} \coqref{SECF.Imp.X}{\coqdocdefinition{X}}\coqdoceol
\coqdocindent{2.50em}
\coqref{SECF.Equiv.Himp.:com:com scope:'while' x 'do' x 'end'}{\coqdocnotation{end}} \coqref{SECF.Imp.::com scope:'<x7B' x 'x7D>'}{\coqdocnotation{\}>}}.\coqdoceol
\coqdocemptyline
\coqdocnoindent
\coqdockw{Definition} \coqdef{SECF.Equiv.Himp.p6}{p6}{\coqdocdefinition{p6}} : \coqref{SECF.Equiv.Himp.com}{\coqdocinductive{com}} :=\coqdoceol
\coqdocindent{1.00em}
\coqref{SECF.Imp.::com scope:'<x7B' x 'x7D>'}{\coqdocnotation{<\{}} \coqref{SECF.Imp.X}{\coqdocdefinition{X}} \coqref{SECF.Equiv.Himp.:com:com scope:x ':=' x}{\coqdocnotation{:=}} 1 \coqref{SECF.Imp.::com scope:'<x7B' x 'x7D>'}{\coqdocnotation{\}>}}.\coqdoceol
\coqdocemptyline
\coqdocnoindent
\coqdockw{Theorem} \coqdef{SECF.Equiv.Himp.p5 p6 equiv}{p5\_p6\_equiv}{\coqdoclemma{p5\_p6\_equiv}} : \coqref{SECF.Equiv.Himp.cequiv}{\coqdocdefinition{cequiv}} \coqref{SECF.Equiv.Himp.p5}{\coqdocdefinition{p5}} \coqref{SECF.Equiv.Himp.p6}{\coqdocdefinition{p6}}.\coqdoceol
\coqdocnoindent
\coqdockw{Proof}.  \coqdocvar{Admitted}.\coqdoceol
\end{coqdoccode}
\ensuremath{\Box} 
\begin{coqdoccode}
\coqdocemptyline
\coqdocnoindent
\coqdockw{End} \coqref{SECF.Equiv.Himp}{\coqdocmodule{Himp}}.\coqdoceol
\coqdocemptyline
\end{coqdoccode}
\section{Additional Exercises}



\paragraph{Exercise: 3 stars, standard, optional (swap\_noninterfering\_assignments)}




    (Hint: You may or may not -- depending how you approach it -- need
    to use \coqdoclemma{functional\_extensionality} explicitly for this one.) 
\begin{coqdoccode}
\coqdocemptyline
\coqdocnoindent
\coqdockw{Theorem} \coqdef{SECF.Equiv.swap noninterfering assignments}{swap\_noninterfering\_assignments}{\coqdoclemma{swap\_noninterfering\_assignments}}: \coqdockw{\ensuremath{\forall}} \coqdef{SECF.Equiv.l1:173}{l1}{\coqdocbinder{l1}} \coqdef{SECF.Equiv.l2:174}{l2}{\coqdocbinder{l2}} \coqdef{SECF.Equiv.a1:175}{a1}{\coqdocbinder{a1}} \coqdef{SECF.Equiv.a2:176}{a2}{\coqdocbinder{a2}},\coqdoceol
\coqdocindent{1.00em}
\coqref{SECF.Equiv.l1:173}{\coqdocvariable{l1}} \coqexternalref{::type scope:x '<>' x}{http://rocq-prover.org/doc/V9.1.1/stdlib//Corelib.Init.Logic}{\coqdocnotation{\ensuremath{\not=}}} \coqref{SECF.Equiv.l2:174}{\coqdocvariable{l2}} \coqexternalref{::type scope:x '->' x}{http://rocq-prover.org/doc/V9.1.1/stdlib//Corelib.Init.Logic}{\coqdocnotation{\ensuremath{\rightarrow}}}\coqdoceol
\coqdocindent{1.00em}
\coqref{SECF.Equiv.var not used in aexp}{\coqdocinductive{var\_not\_used\_in\_aexp}} \coqref{SECF.Equiv.l1:173}{\coqdocvariable{l1}} \coqref{SECF.Equiv.a2:176}{\coqdocvariable{a2}} \coqexternalref{::type scope:x '->' x}{http://rocq-prover.org/doc/V9.1.1/stdlib//Corelib.Init.Logic}{\coqdocnotation{\ensuremath{\rightarrow}}}\coqdoceol
\coqdocindent{1.00em}
\coqref{SECF.Equiv.var not used in aexp}{\coqdocinductive{var\_not\_used\_in\_aexp}} \coqref{SECF.Equiv.l2:174}{\coqdocvariable{l2}} \coqref{SECF.Equiv.a1:175}{\coqdocvariable{a1}} \coqexternalref{::type scope:x '->' x}{http://rocq-prover.org/doc/V9.1.1/stdlib//Corelib.Init.Logic}{\coqdocnotation{\ensuremath{\rightarrow}}}\coqdoceol
\coqdocindent{1.00em}
\coqref{SECF.Equiv.cequiv}{\coqdocdefinition{cequiv}}\coqdoceol
\coqdocindent{2.00em}
\coqref{SECF.Imp.::com scope:'<x7B' x 'x7D>'}{\coqdocnotation{<\{}} \coqref{SECF.Equiv.l1:173}{\coqdocvariable{l1}} \coqref{SECF.Imp.:com:com scope:x ':=' x}{\coqdocnotation{:=}} \coqref{SECF.Equiv.a1:175}{\coqdocvariable{a1}}\coqref{SECF.Imp.:com:com scope:x ';' x}{\coqdocnotation{;}} \coqref{SECF.Equiv.l2:174}{\coqdocvariable{l2}} \coqref{SECF.Imp.:com:com scope:x ':=' x}{\coqdocnotation{:=}} \coqref{SECF.Equiv.a2:176}{\coqdocvariable{a2}} \coqref{SECF.Imp.::com scope:'<x7B' x 'x7D>'}{\coqdocnotation{\}>}}\coqdoceol
\coqdocindent{2.00em}
\coqref{SECF.Imp.::com scope:'<x7B' x 'x7D>'}{\coqdocnotation{<\{}} \coqref{SECF.Equiv.l2:174}{\coqdocvariable{l2}} \coqref{SECF.Imp.:com:com scope:x ':=' x}{\coqdocnotation{:=}} \coqref{SECF.Equiv.a2:176}{\coqdocvariable{a2}}\coqref{SECF.Imp.:com:com scope:x ';' x}{\coqdocnotation{;}} \coqref{SECF.Equiv.l1:173}{\coqdocvariable{l1}} \coqref{SECF.Imp.:com:com scope:x ':=' x}{\coqdocnotation{:=}} \coqref{SECF.Equiv.a1:175}{\coqdocvariable{a1}} \coqref{SECF.Imp.::com scope:'<x7B' x 'x7D>'}{\coqdocnotation{\}>}}.\coqdoceol
\coqdocnoindent
\coqdockw{Proof}.\coqdoceol
 \coqdocvar{Admitted}.\coqdoceol
\end{coqdoccode}
\ensuremath{\Box} 

\paragraph{Exercise: 4 stars, standard, optional (for\_while\_equiv)}




    This exercise extends the optional \coqdocvar{add\_for\_loop} exercise from
    the \coqdoclibrary{Imp} chapter, where you were asked to extend the language
    of commands with C-style \coqdockw{for} loops.  Prove that the command:


      for (c1; b; c2) \{
          c3
      \}


    is equivalent to:


       c1;
       while b do
         c3;
         c2
       end

\begin{coqdoccode}
\coqdocemptyline
\end{coqdoccode}
\paragraph{Exercise: 4 stars, advanced, optional (capprox)}




    In this exercise we define an asymmetric variant of program
    equivalence we call \textit{program approximation}. We say that a
    program \coqdocvar{c1} \textit{approximates} a program \coqdocvar{c2} when, for each of
    the initial states for which \coqdocvar{c1} terminates, \coqdocvar{c2} also terminates
    and produces the same final state. Formally, program approximation
    is defined as follows: 
\begin{coqdoccode}
\coqdocemptyline
\coqdocnoindent
\coqdockw{Definition} \coqdef{SECF.Equiv.capprox}{capprox}{\coqdocdefinition{capprox}} (\coqdef{SECF.Equiv.c1:177}{c1}{\coqdocbinder{c1}} \coqdef{SECF.Equiv.c2:178}{c2}{\coqdocbinder{c2}} : \coqref{SECF.Imp.com}{\coqdocinductive{com}}) : \coqdockw{Prop} := \coqdockw{\ensuremath{\forall}} (\coqdef{SECF.Equiv.st:179}{st}{\coqdocbinder{st}} \coqdef{SECF.Equiv.st':180}{st'}{\coqdocbinder{st'}} : \coqref{SECF.Imp.state}{\coqdocdefinition{state}}),\coqdoceol
\coqdocindent{1.00em}
\coqref{SECF.Equiv.st:179}{\coqdocvariable{st}} \coqref{SECF.Imp.:::x '=[' x ']=>' x}{\coqdocnotation{=[}} \coqref{SECF.Equiv.c1:177}{\coqdocvariable{c1}} \coqref{SECF.Imp.:::x '=[' x ']=>' x}{\coqdocnotation{]=>}} \coqref{SECF.Equiv.st':180}{\coqdocvariable{st'}} \coqexternalref{::type scope:x '->' x}{http://rocq-prover.org/doc/V9.1.1/stdlib//Corelib.Init.Logic}{\coqdocnotation{\ensuremath{\rightarrow}}} \coqref{SECF.Equiv.st:179}{\coqdocvariable{st}} \coqref{SECF.Imp.:::x '=[' x ']=>' x}{\coqdocnotation{=[}} \coqref{SECF.Equiv.c2:178}{\coqdocvariable{c2}} \coqref{SECF.Imp.:::x '=[' x ']=>' x}{\coqdocnotation{]=>}} \coqref{SECF.Equiv.st':180}{\coqdocvariable{st'}}.\coqdoceol
\coqdocemptyline
\end{coqdoccode}
For example, the program


  c1 = while X <> 1 do
         X := X - 1
       end


    approximates \coqdocvar{c2} = \coqref{SECF.SpecCT.X}{\coqdocdefinition{X}} := 1, but \coqdocvar{c2} does not approximate \coqdocvar{c1}
    since \coqdocvar{c1} does not terminate when \coqref{SECF.SpecCT.X}{\coqdocdefinition{X}} = 0 but \coqdocvar{c2} does.  If two
    programs approximate each other in both directions, then they are
    equivalent. 

 Find two programs \coqref{SECF.Equiv.c3}{\coqdocdefinition{c3}} and \coqref{SECF.Equiv.c4}{\coqdocdefinition{c4}} such that neither approximates
    the other. 
\begin{coqdoccode}
\coqdocemptyline
\coqdocnoindent
\coqdockw{Definition} \coqdef{SECF.Equiv.c3}{c3}{\coqdocdefinition{c3}} : \coqref{SECF.Imp.com}{\coqdocinductive{com}}\coqdoceol
\coqdocindent{1.00em}
. \coqdocvar{Admitted}.\coqdoceol
\coqdocnoindent
\coqdockw{Definition} \coqdef{SECF.Equiv.c4}{c4}{\coqdocdefinition{c4}} : \coqref{SECF.Imp.com}{\coqdocinductive{com}}\coqdoceol
\coqdocindent{1.00em}
. \coqdocvar{Admitted}.\coqdoceol
\coqdocemptyline
\coqdocnoindent
\coqdockw{Theorem} \coqdef{SECF.Equiv.c3 c4 different}{c3\_c4\_different}{\coqdoclemma{c3\_c4\_different}} : \coqexternalref{::type scope:'x7E' x}{http://rocq-prover.org/doc/V9.1.1/stdlib//Corelib.Init.Logic}{\coqdocnotation{\ensuremath{\lnot}}} \coqref{SECF.Equiv.capprox}{\coqdocdefinition{capprox}} \coqref{SECF.Equiv.c3}{\coqdocaxiom{c3}} \coqref{SECF.Equiv.c4}{\coqdocaxiom{c4}} \coqexternalref{::type scope:x '/x5C' x}{http://rocq-prover.org/doc/V9.1.1/stdlib//Corelib.Init.Logic}{\coqdocnotation{\ensuremath{\land}}} \coqexternalref{::type scope:'x7E' x}{http://rocq-prover.org/doc/V9.1.1/stdlib//Corelib.Init.Logic}{\coqdocnotation{\ensuremath{\lnot}}} \coqref{SECF.Equiv.capprox}{\coqdocdefinition{capprox}} \coqref{SECF.Equiv.c4}{\coqdocaxiom{c4}} \coqref{SECF.Equiv.c3}{\coqdocaxiom{c3}}.\coqdoceol
\coqdocnoindent
\coqdockw{Proof}.  \coqdocvar{Admitted}.\coqdoceol
\coqdocemptyline
\end{coqdoccode}
Find a program \coqref{SECF.Equiv.cmin}{\coqdocdefinition{cmin}} that approximates every other program. 
\begin{coqdoccode}
\coqdocemptyline
\coqdocnoindent
\coqdockw{Definition} \coqdef{SECF.Equiv.cmin}{cmin}{\coqdocdefinition{cmin}} : \coqref{SECF.Imp.com}{\coqdocinductive{com}}\coqdoceol
\coqdocindent{1.00em}
. \coqdocvar{Admitted}.\coqdoceol
\coqdocemptyline
\coqdocnoindent
\coqdockw{Theorem} \coqdef{SECF.Equiv.cmin minimal}{cmin\_minimal}{\coqdoclemma{cmin\_minimal}} : \coqdockw{\ensuremath{\forall}} \coqdef{SECF.Equiv.c:181}{c}{\coqdocbinder{c}}, \coqref{SECF.Equiv.capprox}{\coqdocdefinition{capprox}} \coqref{SECF.Equiv.cmin}{\coqdocaxiom{cmin}} \coqref{SECF.Equiv.c:181}{\coqdocvariable{c}}.\coqdoceol
\coqdocnoindent
\coqdockw{Proof}.  \coqdocvar{Admitted}.\coqdoceol
\coqdocemptyline
\end{coqdoccode}
Finally, find a non-trivial property which is preserved by
    program approximation (when going from left to right). 
\begin{coqdoccode}
\coqdocemptyline
\coqdocnoindent
\coqdockw{Definition} \coqdef{SECF.Equiv.zprop}{zprop}{\coqdocdefinition{zprop}} (\coqdef{SECF.Equiv.c:182}{c}{\coqdocbinder{c}} : \coqref{SECF.Imp.com}{\coqdocinductive{com}}) : \coqdockw{Prop}\coqdoceol
\coqdocindent{1.00em}
. \coqdocvar{Admitted}.\coqdoceol
\coqdocemptyline
\coqdocnoindent
\coqdockw{Theorem} \coqdef{SECF.Equiv.zprop preserving}{zprop\_preserving}{\coqdoclemma{zprop\_preserving}} : \coqdockw{\ensuremath{\forall}} \coqdef{SECF.Equiv.c:183}{c}{\coqdocbinder{c}} \coqdef{SECF.Equiv.c':184}{c'}{\coqdocbinder{c'}},\coqdoceol
\coqdocindent{1.00em}
\coqref{SECF.Equiv.zprop}{\coqdocaxiom{zprop}} \coqref{SECF.Equiv.c:183}{\coqdocvariable{c}} \coqexternalref{::type scope:x '->' x}{http://rocq-prover.org/doc/V9.1.1/stdlib//Corelib.Init.Logic}{\coqdocnotation{\ensuremath{\rightarrow}}} \coqref{SECF.Equiv.capprox}{\coqdocdefinition{capprox}} \coqref{SECF.Equiv.c:183}{\coqdocvariable{c}} \coqref{SECF.Equiv.c':184}{\coqdocvariable{c'}} \coqexternalref{::type scope:x '->' x}{http://rocq-prover.org/doc/V9.1.1/stdlib//Corelib.Init.Logic}{\coqdocnotation{\ensuremath{\rightarrow}}} \coqref{SECF.Equiv.zprop}{\coqdocaxiom{zprop}} \coqref{SECF.Equiv.c':184}{\coqdocvariable{c'}}.\coqdoceol
\coqdocnoindent
\coqdockw{Proof}.  \coqdocvar{Admitted}.\coqdoceol
\end{coqdoccode}
\ensuremath{\Box} 
\begin{coqdoccode}
\coqdocemptyline
\end{coqdoccode}
\end{document}
